\part{问题定义、历史与研究方法}
\chapter{研究对象、读者与证据方法}
\label{chap:scope-evidence-method}

\section{本书所称“具身智能”}

本书把具身智能定义为：\textbf{通过可观测、可执行的身体与环境形成持续闭环，并在任务、物理和安全约束下感知、估计、决策、行动和更新的智能系统}。这个定义是为比较技术路线而设的操作定义，不替代 ISO 8373 等标准术语\cite{iso8373_2021}。它包含五个必要条件：有明确身体和作用域；有传感与执行接口；动作会改变后续观测；系统针对任务选择行为；失败由物理、任务或安全判据定义。

因此，离线识图模型不是具身系统；把 VLM 接到机械臂上也不自动构成可靠具身智能。反过来，不使用大模型的移动机器人仍可具备闭环自主性。本书关心的不是“是否含 AI”这一标签，而是系统如何在身体和环境约束下形成可复核能力。

\begin{table}[H]
\begingroup
\hbadness=10000
\centering
\caption{研究对象与相邻领域的边界}
\label{tab:scope-boundaries}
\begin{tabularx}{\textwidth}{p{0.16\textwidth}YYY}
\toprule
对象 & 与本书交集 & 本书主讲 & 不在主线的内容 \\
\midrule
机器人学 & 本体、感知、规划、控制 & 学习与基础模型如何进入闭环 & 纯机构设计手册 \\
工业自动化 & 设备、工艺、实时与安全 & 自主策略怎样接入生产系统 & 全套 PLC 编程教程 \\
软件智能体 & 规划、记忆、工具调用 & 动作经身体作用于物理世界 & 纯网页/办公自动化 \\
自动驾驶 & 多传感器、规划、安全 & 可迁移的系统方法 & 道路法规与车辆工程全景 \\
多模态大模型 & 视觉、语言、序列建模 & 与动作、状态和反馈的接口 & 通用聊天模型训练大全 \\
\bottomrule
\end{tabularx}
\endgroup
\end{table}

默认读者具备线性代数、概率和 Python 基础，但未系统学习机器人学。第 4–13 章建立数学、控制和学习基础；第 14–25 章进入多模态、系统架构与 VLA；第 26–30 章讨论数据、评测、部署和安全。产业部分只回答技术交付、供应链和规模化如何受这些约束影响。

\section{本书回答什么，不回答什么}

技术主线回答四类问题：概念和公式如何定义；算法如何训练或求解；代码和系统接口如何实现；证据在哪些任务、硬件和条件下成立。产业主线只在可核事实基础上讨论交付、成本、运维、供应链和政策约束。对于尚无公开证据的商业能力，本书记录“不知道”或“证据不足”，不以融资、演示视频或宣传语补齐结论。

\begingroup\hbadness=10000
本书不是穷尽全部数据库的系统综述，也不进行可合并效应量的元分析（meta-analysis）。PRISMA 2020 强调透明报告检索、筛选和结果\cite{page2021prisma}；本书借鉴其可追溯原则，采用问题导向的教材型技术综述：核心论断优先追溯原始论文、标准、官方代码和可复现实验，并保留检索日期、版本和排除理由。由于机器人任务、硬件和成功判据高度异质，许多结果只能做条件化比较，不能强行合并成统一排名。
\endgroup

\section{证据等级不是一条万能排行榜}

证据强度取决于论断类型。算法机制应看原始论文和源码；接口行为应看对应版本官方文档；标准状态应看标准组织；商业交付应看客户、部署和审计材料。二手综述适合建立广度，不能替代核心模型的原始实验。

\begin{table}[H]
\centering
\caption{论断类型与首选证据}
\label{tab:claim-evidence-map}
\begin{tabularx}{\textwidth}{p{0.20\textwidth}YYp{0.22\textwidth}}
\toprule
论断 & 首选来源 & 必核字段 & 常见误用 \\
\midrule
原理/训练目标 & 原始论文、补充材料 & 假设、公式、数据、消融 & 用新闻稿解释机制 \\
实现/接口 & 官方仓库、固定 commit、文档 & 文件、符号、版本、许可证 & 链接默认分支 \\
性能/SOTA & 原始结果、独立复现 & 任务、硬件、分母、区间 & 跨 benchmark 排名 \\
标准/政策 & 发布机构正式文本 & 编号、版本、状态、范围 & 引用已撤销版本 \\
部署/商业 & 客户验证、运行记录、审计 & 地点、时间、暴露量、介入 & 把试点写成规模交付 \\
图像/代码转载 & 权利页、许可证、作者许可 & 资产、版本、署名、改动 & 把公开可见当可复制 \\
\bottomrule
\end{tabularx}
\end{table}

一个核心论断 $c$ 的最小证据记录写为
\begin{equation}
E(c)=(s,v,l,a,t,r),
\label{eq:claim-evidence-record}
\end{equation}
其中 $s$ 是来源，$v$ 是版本或日期，$l$ 是页码/图号/代码行等定位，$a$ 是适用条件，$t$ 是论断类型，$r$ 是权利和复用状态。只有 URL 而没有定位和版本，难以支撑会变化的实现事实。

\begin{figure}[H]
\centering
\begin{tabular}{c@{\ $\rightarrow$\ }c@{\ $\rightarrow$\ }c@{\ $\rightarrow$\ }c}
\fbox{论断与边界} & \fbox{原始来源/版本} & \fbox{公式、代码或实验} & \fbox{复核与评价}
\end{tabular}
\caption{本书的最小可追溯链。评价只能位于机制和证据之后。}
\label{fig:claim-traceability}
\end{figure}

\section{事实、推断与评价分层}

“论文报告在某任务成功 $k/n$ 次”是事实；“差异可能来自动作频率和复位协议”是推断；“该证据不足以支持通用部署”是评价。三者可以出现在同一节，但不能伪装成同一种确定性。事实使用明确引用；推断列出依据和替代解释；评价说明判断标准。涉及未来、未公开实现或商业数字时，标注时间和不确定性。

\begin{lstlisting}[caption={章节证据包的最低检查；这是写作流程伪代码，不是检索算法}]
for claim in chapter.core_claims:
    require claim.type in {fact, inference, evaluation}
    if claim.type == fact:
        require primary_source(claim) and version_and_locator(claim)
    if claim.concerns_number_or_capability:
        require task_hardware_denominator_and_conditions(claim)
for asset in chapter.figures_tables_code:
    require provenance(asset) and reuse_status(asset)
reject bare_url_as_academic_evidence()
\end{lstlisting}

\section{知识小结与方法判断}

本书的对象是受身体、任务、反馈和安全共同约束的闭环系统。证据按论断匹配来源，并记录版本、定位、适用条件和权利状态。方法判断是：综述的可信度不由参考文献数量单独决定，而由关键论断能否沿“来源—机制/实现—实验—边界”回溯决定；无法回溯的确定语气应被删除或降级。

\chapter{从经典机器人到基础模型：范式与分类坐标}
\label{chap:paradigms-taxonomy}

\section{范式变化是接口和证据变化，不是模型年表}

经典机器人、学习型机器人和基础模型机器人并非前后替代。高精度伺服、状态估计和安全互锁仍是基础模型系统的底座；变化主要发生在任务如何指定、知识从哪里获得、动作如何生成以及泛化如何评测。Brooks 的行为式路线强调智能应通过真实世界的感知—行动耦合逐步构建\cite{brooks1991representation}，它挑战的是必须先建立完整符号世界模型的假设，而不是否认一切内部状态。

基础模型通常从广泛数据训练，再适配多个下游问题\cite{bommasani2021foundation}。进入机器人后，适配对象不再只是文本任务，还包括本体、动作空间、控制频率、相机布局和安全包络。RT-1 以多任务机器人数据训练动作 token 策略\cite{brohan2022rt1}，PaLM-E 把连续传感输入接入语言模型\cite{driess2023palme}，RT-2 探索视觉语言知识到动作的迁移\cite{brohan2023rt2}；这些工作改变了数据和模型接口，却没有取消动力学、控制与真机评测。

\begin{table}[H]
\begingroup
\hbadness=10000
\centering
\caption{五类范式按问题设定比较，而非按年代互相取代}
\label{tab:robot-paradigms}
\begin{tabularx}{\textwidth}{p{0.16\textwidth}YYYY}
\toprule
范式 & 任务/知识来源 & 核心接口 & 主要证据 & 典型边界 \\
\midrule
固定自动化 & 工艺与工程师规则 & I/O、轨迹、互锁 & 节拍、精度、故障率 & 变化需重编程 \\
感知—规划—控制 & 地图、模型、目标 & 状态、规划、控制器 & 完备性、误差、闭环试验 & 建模与开放世界差距 \\
行为式/反应式 & 局部感知与行为组合 & 感知—行动通道 & 实时环境行为 & 长期任务与全局约束 \\
学习型机器人 & 示教、奖励、仿真交互 & 数据集、策略、价值 & 分布内外成功率、消融 & 数据偏移与安全探索 \\
基础模型机器人 & 跨任务/跨模态预训练 & token/latent、动作头、上下文 & 迁移、组合泛化、系统评测 & 身体差异、时延、证据可比性 \\
\bottomrule
\end{tabularx}
\endgroup
\end{table}

\section{用五个坐标描述机器人任务}

“通用机器人”缺少可计算边界。一个评测对象至少应写成
\begin{equation}
\mathcal R=(B,T,C,O,A),
\label{eq:robot-taxonomy-tuple}
\end{equation}
其中 $B$ 是身体与动作能力，$T$ 是任务族，$C$ 是接触和工具关系，$O$ 是环境开放度，$A$ 是自主组织方式。

$B$ 包括移动形式、自由度、负载、传感器和动作接口；$T$ 区分导航、抓取、装配、移动操作和人机协作等任务结构；$C$ 记录自由空间、刚性接触、柔顺接触、双臂/多体和工具使用；$O$ 记录对象、场景、人员和扰动是否在训练覆盖内；$A$ 不宜只给单一“自主等级”，还要说明任务分配、监督、恢复和人工接管分别由谁承担。

\begin{figure}[H]
\centering
\setlength{\tabcolsep}{2pt}
\begin{tabular}{c@{\ $\rightarrow$\ }c@{\ $\rightarrow$\ }c}
\fbox{\begin{minipage}{0.15\textwidth}\centering 身体 $B$\\与任务 $T$\end{minipage}} &
\fbox{\begin{minipage}{0.24\textwidth}\centering 接触 $C$、开放度 $O$\\自主组织 $A$\end{minipage}} &
\fbox{\begin{minipage}{0.18\textwidth}\centering 评测协议\\与证据边界\end{minipage}}
\end{tabular}
\caption{分类坐标的目的不是贴标签，而是生成可比较的评测契约。}
\label{fig:taxonomy-to-evaluation}
\end{figure}

例如“在桌面抓取上成功”至少缺少夹爪、对象分布、是否允许碰撞、相机、动作频率、复位和人工介入。把这些字段补齐后，模型结果才可能进入第 24、28 章的比较框架。

\section{能力是条件化系统量}

模型参数 $\theta$ 不能脱离身体与运行时谈能力。更合适的表示是
\begin{equation}
P(Y=1\mid \theta,B,T,C,O,A,D,V),
\label{eq:conditional-system-capability}
\end{equation}
其中 $D$ 是训练/适配数据，$V$ 是软件、硬件、标定和策略版本。更换机械臂、相机或动作归一化会改变条件分布；允许人工恢复会改变 $Y$ 的含义。因此“模型会装配”“模型具备空间推理”只能作为假设，必须落到可观测成功判据和失败分层。

Open X-Embodiment 通过统一数据组织推动跨本体训练\cite{openx2024}，但统一 schema 不等于统一物理语义：同名末端位姿可能使用不同坐标系，同为夹爪动作可能对应位置、速度或开合事件。基础模型扩大了知识和数据复用范围，也把接口语义、数据治理和部署验证变成主要研究对象。

\section{端到端与模块化是一条连续谱}

端到端通常指损失从观测直接监督动作，模块化指中间状态与组件接口显式。实际系统介于两端：视觉编码器、语言模型和动作头可以联合训练，但外部仍有时间同步、控制器、安全过滤和恢复状态机；传统规划系统也可能使用学习感知或代价函数。比较路线时应问哪些中间量可观测、哪些模块可替换、误差如何定位、闭环如何恢复，而不是只问“是不是端到端”。

\begin{table}[H]
\begingroup
\hbadness=10000
\centering
\caption{路线选择应由失效代价和证据需求驱动}
\label{tab:end-to-end-modular-tradeoff}
\begin{tabularx}{\textwidth}{p{0.22\textwidth}YYY}
\toprule
问题 & 更偏端到端的收益 & 更偏显式接口的收益 & 必做验证 \\
\midrule
感知到动作映射 & 联合优化、少手工标签 & 可诊断状态与约束 & 遮挡、漂移、闭环偏移 \\
跨任务复用 & 共享表示和上下文 & 技能可组合与替换 & 新任务/新本体迁移 \\
实时与安全 & 可压缩单次推理链 & 硬截止期与独立监督 & 最坏时延、回退和故障注入 \\
运维更新 & 统一模型发布 & 局部升级和责任隔离 & schema、兼容与回滚 \\
\bottomrule
\end{tabularx}
\endgroup
\end{table}

\section{知识小结与综述判断}

机器人范式的连续演进体现在知识来源、接口、反馈和评测发生变化，而不是新模型抹去旧控制栈。五维坐标 $(B,T,C,O,A)$ 把抽象能力还原到身体和任务条件。综述判断是：所谓“机器人基础模型”只有在跨任务/跨本体的数据复用、动作接口、闭环泛化和系统约束上给出证据时才是技术类别；仅把通用模型接到机器人演示上，不足以证明范式转变。

\chapter{具身系统的公共架构与核心矛盾}
\label{chap:common-architecture-tensions}

\section{公共栈由数据流和责任边界定义}

不同机器人可以使用不同模型，但都必须把环境变化转成观测、状态、决策和执行，再把结果回流。公共架构不是固定产品框图，而是一组责任：感知提供带时间和不确定度的测量；状态层维护决策可消费的状态；世界表征支持预测和记忆；任务层产生受约束意图；技能层把意图实例化；控制层满足动力学和截止期；安全层拥有过滤、降级和停止权限；数据层记录全过程以供评测和更新。

\begin{figure}[H]
\centering
\begin{tabular}{c@{\ $\rightarrow$\ }c@{\ $\rightarrow$\ }c@{\ $\rightarrow$\ }c}
\fbox{环境/身体} & \fbox{感知与状态} & \fbox{世界模型/任务规划} & \fbox{技能与控制}
\end{tabular}

\vspace{0.7em}
\begin{tabular}{c@{\quad}c@{\quad}c}
\fbox{安全监督：过滤、抢占、停止} & \fbox{运行日志与评测} & \fbox{数据治理与版本更新}
\end{tabular}
\caption{全书公共系统栈。安全和数据不是末端附加模块，而是横跨闭环的旁路。}
\label{fig:common-embodied-stack}
\end{figure}

接口记录可写成
\begin{equation}
m=(y,t_s,t_r,f,\Sigma,q,v),
\label{eq:interface-message-contract}
\end{equation}
其中 $y$ 是值，$t_s/t_r$ 是采样/接收时刻，$f$ 是坐标系或语义 frame，$\Sigma$ 是不确定度，$q$ 是质量/有效标志，$v$ 是 schema 与生产者版本。许多集成故障并非模型不会推理，而是坐标系、单位、时间、新鲜度或版本没有进入接口合同。

\section{多时间尺度系统不能由单一循环解释}

\begin{table}[H]
\begingroup
\hbadness=10000
\centering
\caption{具身系统的多时间尺度及主要失效}
\label{tab:common-timescales}
\begin{tabularx}{\textwidth}{p{0.19\textwidth}p{0.18\textwidth}YYp{0.20\textwidth}}
\toprule
层级 & 典型尺度 & 输入/输出 & 首要约束 & 超时后行为 \\
\midrule
驱动与伺服 & 亚毫秒至毫秒 & 编码器/力—命令 & 抖动、稳定性 & 制动或保持 \\
估计与局部控制 & 毫秒至百毫秒 & 测量—状态/轨迹 & 新鲜度、可观测性 & 降速、拒绝旧状态 \\
技能与动作策略 & 百毫秒至秒 & 目标—动作块 & 闭环纠错、可取消 & 重试或切换技能 \\
任务与人机交互 & 秒至分钟 & 指令—任务图 & 语义、资源、权限 & 澄清、重规划 \\
数据与模型更新 & 小时至月 & 运行记录—版本 & 治理、回归、回滚 & 保持已知良好版本 \\
\bottomrule
\end{tabularx}
\endgroup
\end{table}

\texttt{ros2\_control} 的 read--update--write 把硬件状态、控制器与命令接口分开\cite{ros2control2026}，是毫秒级闭环的一种实现骨架；它不能替代秒级技能取消、分钟级任务恢复或长期数据治理。多速率系统应在接口上声明最大年龄、截止期和取消语义。

\begin{lstlisting}[caption={多速率具身运行时的教学骨架；各回路必须在独立调度与测试中实现}]
while robot_enabled:
    state = realtime_read_and_estimate()
    if state.stale or safety_violation(state): safe_fallback()
    realtime_control_update(state, latest_bounded_skill_command())
    if skill_tick_due():
        skill = update_or_recover_skill(state, task_intent)
    if task_tick_due():
        task_intent = plan_or_request_clarification(state, memory)
    append_only_log(state, task_intent, skill, interventions())
\end{lstlisting}

\section{四组核心矛盾没有一次性答案}

\begin{table}[H]
\begingroup
\hbadness=10000
\centering
\caption{贯穿全书的设计矛盾及可证伪问题}
\label{tab:core-system-tensions}
\begin{tabularx}{\textwidth}{p{0.22\textwidth}YYp{0.24\textwidth}}
\toprule
矛盾 & 一侧收益 & 另一侧收益 & 必须用什么证据裁决 \\
\midrule
开放泛化—安全包络 & 覆盖未知任务 & 行为可预测、风险可控 & OOD、约束接近、接管与事故率 \\
端到端—模块化 & 联合优化、少接口损失 & 可诊断、替换与认证 & 同预算消融、故障定位、恢复 \\
模型规模—实时部署 & 知识与容量 & 时延、功耗与可用性 & 热稳态 p99、任务收益、回退 \\
统一接口—身体差异 & 数据与技能复用 & 保留真实动力学语义 & 跨本体迁移与接口误差 \\
\bottomrule
\end{tabularx}
\endgroup
\end{table}

开放泛化不能只用更多成功视频证明，因为未知行为的尾部风险可能同时上升；端到端与模块化要在相同数据和计算预算下比较；大模型必须报告整链时延和功耗；统一动作空间必须说明如何映射到不同自由度、限位和控制模式。这些矛盾决定后续章节的证据问题，而不是在本章预先选边。

\section{安全旁路与恢复必须拥有真实权限}

\begingroup\hbadness=10000
安全过滤可以在名义动作外增加约束，控制屏障函数是代表性方法\cite{ames2017cbf}；Simplex 思路则在复杂控制器失去可信条件时切换到更强验证的基线控制器\cite{sha2001simplex}。二者共同说明：安全模块必须能拒绝或替换动作，而不只是记录警告。恢复同样不是重新提示模型，而是包含停止、回撤、重新感知、重规划、请求人工和任务终止的状态机。
\endgroup

数据闭环要记录过滤前后动作、触发原因、人工介入和终止语义，否则训练只看到“机器人最终成功”，会把安全监督器和人工恢复的贡献错误归给策略。第 21、26、28、30 章分别展开运行时恢复、数据血缘、评测统计和安全放行。

\section{全书公共符号与接口约定}

\begin{table}[H]
\centering
\caption{后续章节反复使用的最小符号集}
\label{tab:common-symbols}
\begin{tabularx}{\textwidth}{p{0.16\textwidth}Yp{0.18\textwidth}Y}
\toprule
符号 & 含义 & 符号 & 含义 \\
\midrule
$x_t$ & 时刻 $t$ 的真实/估计状态，语境中注明 & $o_t$ & 带时间戳的观测 \\
$a_t,u_t$ & 高层动作/控制输入 & $g$ & 任务目标或条件 \\
$\pi_\theta$ & 参数化策略 & $f_\phi$ & 动力学、世界或预测模型 \\
$\tau$ & 轨迹或 episode & $\mathcal C$ & 约束或安全集合 \\
$D,V$ & 数据与版本集合 & $\mathcal E$ & 评测契约 \\
\bottomrule
\end{tabularx}
\end{table}

状态和观测不得混用；动作与控制输入在动作接口章进一步区分；成功率必须连同评测契约；版本 $V$ 至少覆盖模型、代码、硬件、标定和环境。局部章节若改变符号含义，应显式声明。

\section{知识小结与综述判断}

公共架构由感知、状态、预测、任务、技能、控制、安全和数据闭环的责任与接口组成，并运行在多个时间尺度。四组核心矛盾要求用失效、实时、迁移和恢复证据裁决。综述判断是：架构图的价值不在框的数量，而在每条边是否携带坐标、时间、不确定度、版本和失败语义，以及安全与恢复是否拥有实际抢占权限。


\part{机器人数学与控制基础}
\chapter{数学与计算预备知识}
\label{chap:math-computation}

\section{学习目标、前置知识与问题边界}

本章只讲后续机器人章节会反复调用的计算工具。读者应能在读完后完成四件事：判断一个矩阵问题是否病态；把传感噪声传播为状态不确定性；把估计问题写成概率或最小二乘目标；识别约束优化的可行性、最优性与数值收敛是三件不同的事。前置知识是微积分、向量运算和 Python。$SO(3)$、$SE(3)$ 在第~\ref{chap:kinematics-lie} 章主讲，贝叶斯滤波在第 8 章主讲，本章只建立它们共同使用的符号和数值约定。

机器人程序很少因“不知道矩阵乘法”而失败，更常见的失败是坐标、量纲、随机性与数值条件混在一起。例如，相机外参误差进入抓取位姿后会沿运动链传播；雅可比接近奇异时，毫米级末端修正可能对应极大的关节速度；优化器返回“收敛”也只说明停止条件满足，不保证任务约束真实可行。图~\ref{fig:math-compute-chain} 把本章工具放进一条共同链路。

\begin{figure}[H]
\centering
\small
\fbox{\parbox{0.17\textwidth}{\centering 状态与参数\\$x,\theta$}}
$\longrightarrow$
\fbox{\parbox{0.18\textwidth}{\centering 传感与噪声\\$z=h(x)+v$}}
$\longrightarrow$
\fbox{\parbox{0.18\textwidth}{\centering 线性化与传播\\$J\Sigma J^\top$}}
$\longrightarrow$
\fbox{\parbox{0.18\textwidth}{\centering 目标与约束\\$\min f,\ c=0$}}
$\longrightarrow$
\fbox{\parbox{0.13\textwidth}{\centering 动作\\$u$}}
\caption{机器人计算中的变量、噪声、局部模型与优化接口。作者教学示意；箭头表示计算依赖，不表示固定的软件模块。}
\label{fig:math-compute-chain}
\end{figure}

\section{向量、矩阵与数值敏感性}

设 $x\in\mathbb{R}^{n}$、$A\in\mathbb{R}^{m\times n}$。内积 $x^\top y$ 衡量同一坐标基下的方向一致性；二范数 $\lVert x\rVert_2=\sqrt{x^\top x}$ 衡量欧氏长度；加权范数
\begin{equation}
\lVert x\rVert_W^2=x^\top W x,\qquad W\succ0
\label{eq:weighted-norm}
\end{equation}
把不同方向的单位与可信度写入同一标量。若 $x$ 是位置与角度拼成的误差向量，直接使用单位阵意味着把 $1$ m 与 $1$ rad 当作同等代价；这不是“中性选择”，而是一个隐含权重决定。

线性方程 $Ax=b$ 的数学解写作 $x=A^{-1}b$，实现时通常应调用分解后的线性求解，而不是显式构造 $A^{-1}$。NumPy 的官方线性代数接口也把 `solve`、奇异值分解和条件数作为不同任务提供；病态矩阵即使没有触发奇异异常，结果仍可能因浮点误差而失真\cite{harris2020numpy,numpy_linalg_docs}。二范数条件数为
\begin{equation}
\kappa_2(A)=\lVert A\rVert_2\lVert A^{-1}\rVert_2
=\frac{\sigma_{\max}(A)}{\sigma_{\min}(A)}.
\label{eq:condition-number}
\end{equation}
这里 $\sigma_{\max}$、$\sigma_{\min}$ 是最大和最小奇异值。$\kappa$ 不是误差本身，而是相对输入扰动可能被放大的上界尺度；$\sigma_{\min}\to0$ 时，某个输入方向几乎无法从输出中分辨。

取
\[
A=\begin{bmatrix}1&1\\1&1.0001\end{bmatrix},\quad
b=\begin{bmatrix}2\\2.0001\end{bmatrix}.
\]
精确解为 $x=[1,1]^\top$，但 $A$ 的两列几乎共线，$\kappa_2(A)\approx4.0\times10^4$。若第二个观测仅改变 $10^{-4}$ 的量级，解的分量就可能出现数量级更大的相对变化。后续第 5 章的雅可比奇异、第 8 章的不可观测方向，本质上都会再次出现这种“小奇异值对应弱信息方向”的结构。

\begin{table}[H]
\centering
\caption{常见矩阵问题的诊断量、失败信号与处理方向}
\label{tab:linear-algebra-failures}
\small
\begin{tabularx}{\textwidth}{p{0.19\textwidth}p{0.19\textwidth}p{0.25\textwidth}Y}
\toprule
问题 & 先看什么 & 失败信号 & 处理方向\\
\midrule
$Ax=b$ & 条件数、残差 & 解随微小输入剧烈变化 & 缩放变量、分解求解、正则化\\
最小二乘 & 奇异值谱、秩 & 参数不可辨、协方差巨大 & 删除冗余参数或加入可信先验\\
协方差矩阵 & 对称性、特征值 & 出现负方差或非正定 & 检查线性化、数值对称化与模型\\
梯度/Hessian & 梯度范数、曲率 & 振荡、停在鞍点、步长崩溃 & 线搜索、信赖域、重新参数化\\
\bottomrule
\end{tabularx}
\end{table}

\section{概率、贝叶斯更新与不确定性传播}

机器人状态不是“真值加一个误差条”，而是基于观测和模型形成的条件分布。贝叶斯公式
\begin{equation}
p(x\mid z)=\frac{p(z\mid x)p(x)}{p(z)}
\label{eq:bayes-rule}
\end{equation}
把先验 $p(x)$ 与似然 $p(z\mid x)$ 组合为后验。分母只负责归一化；进行最大后验估计时，可最小化负对数而不显式计算它。若先验与观测噪声均为高斯，负对数后验变成加权平方残差，这就是滤波、束调整和许多轨迹优化目标共享的数学骨架\cite{barfoot2024state}。

设 $x\sim\mathcal{N}(\mu_x,\Sigma_x)$，输出为 $y=f(x)$。在 $\mu_x$ 附近一阶展开
\begin{equation}
f(x)\approx f(\mu_x)+J_f(x-\mu_x),\qquad
J_f=\left.\frac{\partial f}{\partial x}\right|_{\mu_x},
\end{equation}
可得
\begin{equation}
\mu_y\approx f(\mu_x),\qquad
\Sigma_y\approx J_f\Sigma_xJ_f^\top.
\label{eq:first-order-covariance}
\end{equation}
$J_f\in\mathbb{R}^{p\times n}$ 把 $n$ 维输入扰动映射到 $p$ 维输出扰动。式~\eqref{eq:first-order-covariance} 是局部近似：非线性强、分布多峰或角度跨越分支切换时，均值和协方差不足以描述真实输出。

最小算例取 $y=x_1+2x_2$、$\Sigma_x=\operatorname{diag}(0.01,0.04)$。此时 $J_f=[1,2]$，所以
\[
\sigma_y^2=[1,2]\begin{bmatrix}0.01&0\\0&0.04\end{bmatrix}[1,2]^\top=0.17,
\quad \sigma_y\approx0.412.
\]
$x_2$ 的系数为 2，方差贡献因此乘以 $2^2$ 而不是 2。工程上常见的“把标准差直接相加”会在这里给出错误结果。

\section{最大似然、最大后验与非线性最小二乘}

给定观测 $z_i=h_i(x)+v_i$，且 $v_i\sim\mathcal{N}(0,R_i)$，最大似然估计等价于
\begin{equation}
\hat{x}_{\mathrm{ML}}=arg\min_x\frac{1}{2}\sum_i
r_i(x)^\top R_i^{-1}r_i(x),\qquad r_i(x)=z_i-h_i(x).
\label{eq:nonlinear-least-squares}
\end{equation}
$R_i^{-1}$ 是信息矩阵：噪声方差越大，观测权重越小。若增加高斯先验 $x\sim\mathcal{N}(\bar{x},P)$，目标中再加入 $\frac12\lVert x-\bar{x}\rVert_{P^{-1}}^2$，得到最大后验估计。所谓“融合多个传感器”不是把数值平均，而是在共同状态与坐标约定下组合各自的残差模型和不确定性。

非线性最小二乘在当前估计 $x_k$ 处写 $r(x_k+\Delta x)\approx r_k+J_k\Delta x$。Gauss--Newton 步满足
\begin{equation}
(J_k^\top WJ_k)\Delta x=-J_k^\top Wr_k.
\label{eq:gauss-newton}
\end{equation}
前式的输出 $J_k$ 和 $r_k$ 成为后式的线性系统。若 $J_k^\top WJ_k$ 病态，更新会不稳定；Levenberg--Marquardt 通过加入 $\lambda I$ 抑制弱方向，但会牺牲局部收敛速度。优化教材通常把线搜索、信赖域与约束条件分开处理，因为“目标下降”“局部模型可信”和“约束可行”不能由同一个停止标志代替\cite{nocedal2006numerical}。

\section{梯度、约束与拉格朗日乘子}

考虑等式约束问题
\begin{equation}
\min_x f(x)\quad\text{s.t.}\quad c(x)=0.
\label{eq:equality-constrained-problem}
\end{equation}
拉格朗日函数 $\mathcal{L}(x,\lambda)=f(x)+\lambda^\top c(x)$。在满足约束正则性的局部最优点，一阶必要条件为
\begin{equation}
\nabla_x f(x^*)+J_c(x^*)^\top\lambda^*=0,
\qquad c(x^*)=0.
\label{eq:kkt-equality}
\end{equation}
第一式表示目标梯度在可行切空间上不能继续下降；$\lambda$ 可解释为约束收紧一单位时的局部代价敏感度。它不是一个“惩罚系数调参”：乘子由最优性条件共同求解，而惩罚法中的权重由设计者选择。

以 $\min_{x,y}x^2+y^2$ 且 $x+y=1$ 为例，$\mathcal{L}=x^2+y^2+\lambda(x+y-1)$。由 $2x+\lambda=0$、$2y+\lambda=0$ 得 $x=y$，再代入约束得到 $x=y=0.5$、$\lambda=-1$。若约束改为非凸集合，式~\eqref{eq:kkt-equality} 通常只提供局部必要条件；初始化、可行域拓扑和求解器容差仍决定实际结果。

\section{NumPy 可复算例与实现约定}

下面代码把本章三个计算动作放在同一处：不用显式逆求解线性方程；用奇异值计算条件数；用雅可比传播协方差。它是可运行的教学代码，不是 NumPy 源码片段；API 行为依据 NumPy 官方文档，数组计算模型依据 NumPy 论文\cite{harris2020numpy,numpy_linalg_docs}。

\begin{lstlisting}[language=Python,caption={条件数、线性求解与一阶协方差传播的 NumPy 教学代码},label={lst:ch04-numpy}]
import numpy as np

A = np.array([[1.0, 1.0], [1.0, 1.0001]])
b = np.array([2.0, 2.0001])
s = np.linalg.svd(A, compute_uv=False)
kappa = s.max() / s.min()
x = np.linalg.solve(A, b)

Sigma_x = np.diag([0.01, 0.04])
J = np.array([[1.0, 2.0]])
Sigma_y = J @ Sigma_x @ J.T
print(kappa, x, Sigma_y.item())
\end{lstlisting}

实现中还应统一四项约定。第一，向量默认列向量，批量维度写在最前或最后必须全书一致；第二，位置用 m、角度用 rad、时间用 s，日志不得省略单位；第三，协方差必须注明对应误差坐标；第四，任何 `solve`、SVD 或优化调用都记录残差、迭代次数与终止原因。没有这些元数据，数值结果无法用于失败诊断。

\section{工程失效与知识小结}

本章工具的共同限制是局部性。条件数描述局部线性敏感性，不告诉系统如何重新规划；一阶协方差传播不能忠实表示遮挡造成的多峰假设；Gauss--Newton 可能在错误数据关联下稳定收敛到错误答案；KKT 条件也不保证非凸问题的全局最优。工程程序应把“求解器成功”“残差足够小”“约束可行”“状态在业务边界内”分别检查。

\authorjudgement{机器人数学不是为了增加符号密度，而是为了把接口假设显式化。后续章节若出现逆矩阵、协方差、梯度或约束，却没有说明坐标、量纲、条件数与停止条件，就不能把该结果视为可部署证据。}

\chapter{位姿、李群、运动学与雅可比}
\label{chap:kinematics-lie}

\section{学习目标、前置知识与问题设定}

本章回答一个贯穿全书的问题：高层系统给出的“把末端向左移动 3 cm 并旋转 10$^\circ$”如何变成定义明确、可计算、能被关节执行的量。读者应能区分坐标变换与物体运动，能在旋转矩阵、四元数和李代数增量之间选择接口，能从关节链计算末端位姿与雅可比，并能解释逆运动学不收敛究竟来自不可达、奇异、约束冲突还是数值设置。

前置知识是第~\ref{chap:math-computation} 章的矩阵乘法、奇异值、条件数和局部线性化。本章不讲力矩和接触力，它们在第~\ref{chap:dynamics-contact-compliance} 章主讲；也不比较 VLA 的动作 token，它们最终仍须经过本章的位姿或关节接口落到机器人上。

\section{坐标系、齐次变换与变换链}

约定 ${}^{A}p\in\mathbb{R}^3$ 表示点 $p$ 在坐标系 $A$ 中的坐标，${}^{A}T_B$ 表示把 $B$ 系坐标变换到 $A$ 系。齐次变换为
\begin{equation}
{}^{A}T_B=
\begin{bmatrix}
{}^{A}R_B & {}^{A}t_B\\
0_{1\times3} & 1
\end{bmatrix}\in SE(3),\qquad
\begin{bmatrix}{}^{A}p\\1\end{bmatrix}
={}^{A}T_B\begin{bmatrix}{}^{B}p\\1\end{bmatrix}.
\label{eq:homogeneous-transform}
\end{equation}
${}^{A}R_B\in SO(3)$ 无单位，${}^{A}t_B$ 的单位通常为 m。上标与下标不是装饰：中间坐标系必须相消，因而
\begin{equation}
{}^{W}T_E={}^{W}T_B\,{}^{B}T_1\cdots{}^{n}T_E.
\label{eq:transform-chain}
\end{equation}
若把某一项写反，矩阵维度仍可能合法，但几何语义已经错误。Modern Robotics 将刚体运动、运动学、雅可比和控制放在统一 $SE(3)$ 语言下，正是为了让这种接口关系可计算且可检查\cite{lynch2017modern,modern_robotics}。

取二维简化例：$B$ 系相对 $W$ 系旋转 $90^\circ$，并平移 $(1,0)$ m；点在 $B$ 系为 $(1,0)$ m。旋转后该点相对 $B$ 原点位于 $(0,1)$ m，再加平移得 ${}^{W}p=(1,1)$ m。若先平移点再旋转，会得到不同结果；“旋转和平移差不多可以交换”只在特定小量近似中成立。

\begin{figure}[H]
\centering
\small
\fbox{\parbox{0.16\textwidth}{\centering 关节状态\\$q$}}
$\xrightarrow{\text{正运动学}}$
\fbox{\parbox{0.19\textwidth}{\centering 当前位姿\\$T(q)$}}
$\xrightarrow{\log(T^{-1}T_d)}$
\fbox{\parbox{0.17\textwidth}{\centering 位姿误差\\$e\in\mathbb{R}^6$}}
$\xrightarrow{J^\#}$
\fbox{\parbox{0.17\textwidth}{\centering 关节增量\\$\Delta q$}}
$\xrightarrow{\operatorname{integrate}}$
\fbox{\parbox{0.11\textwidth}{\centering 新状态}}
\caption{从关节状态到李群位姿误差、雅可比逆和流形积分的闭环运动学计算链。作者依据 Modern Robotics 与 Pinocchio 官方逆运动学示例改绘\cite{lynch2017modern,pinocchio340ik}。}
\label{fig:kinematic-compute-chain}
\end{figure}

\section{旋转矩阵、欧拉角与四元数}

旋转矩阵集合
\begin{equation}
SO(3)=\{R\in\mathbb{R}^{3\times3}\mid R^\top R=I,\ \det R=1\}
\label{eq:so3-definition}
\end{equation}
有 9 个元素但只有 3 个自由度。正交约束使它可以直接作用于向量并稳定组合，但数值积分后需要重新正交化。欧拉角只用 3 个数，适合人机界面，却依赖旋转顺序并在特定姿态发生参数奇异；同一组“roll、pitch、yaw”若未说明内禀/外禀和轴顺序，并不是完整数据。

单位四元数 $q=[q_w,q_x,q_y,q_z]^\top$ 满足 $\lVert q\rVert=1$。它避免欧拉角的参数奇异，插值与积分也更方便，但 $q$ 与 $-q$ 表示同一旋转，比较或学习时必须处理双覆盖。传感器融合中若把四元数四个分量当作无约束欧氏状态直接相加，既可能破坏单位范数，也会把符号跳变误判为大旋转。Solà 对四元数运动学与误差状态的系统说明表明，常见做法是让名义姿态保留在单位四元数上，把小误差放在三维切空间中更新\cite{sola2017quaternion}。

\begin{table}[H]
\centering
\caption{旋转表示的接口代价与典型失效}
\label{tab:rotation-representations}
\small
\begin{tabularx}{\textwidth}{p{0.16\textwidth}p{0.15\textwidth}p{0.24\textwidth}p{0.23\textwidth}Y}
\toprule
表示 & 存储量 & 主要用途 & 约束/歧义 & 典型失效\\
\midrule
欧拉角 & 3 & 人机界面、有限范围标定 & 旋转顺序必须声明 & 万向节锁、跨界跳变\\
旋转矩阵 & 9 & 坐标变换、组合、几何计算 & $R^\top R=I,\det R=1$ & 数值漂移、冗余优化\\
单位四元数 & 4 & 插值、惯导、姿态状态 & 单位范数、$q\sim -q$ & 未归一化、符号不连续\\
李代数增量 & 3 & 局部误差、优化、滤波更新 & 依赖线性化点和对数分支 & 大角度线性化失真\\
\bottomrule
\end{tabularx}
\end{table}

\section{李群、指数映射与局部误差}

反对称矩阵算子把 $\omega=[\omega_x,\omega_y,\omega_z]^\top$ 写为
\begin{equation}
[\omega]_\times=
\begin{bmatrix}
0&-\omega_z&\omega_y\\
\omega_z&0&-\omega_x\\
-\omega_y&\omega_x&0
\end{bmatrix},\qquad [\omega]_\times v=\omega\times v.
\end{equation}
$\mathfrak{so}(3)$ 是这些反对称矩阵组成的切空间。指数映射
\begin{equation}
R=\operatorname{Exp}(\phi)=\exp([\phi]_\times)
=I+\frac{\sin\theta}{\theta}[\phi]_\times
+\frac{1-\cos\theta}{\theta^2}[\phi]_\times^2,
\quad \theta=\lVert\phi\rVert
\label{eq:rodrigues}
\end{equation}
把轴角向量 $\phi$ 转成有限旋转。$\theta\to0$ 时系数应使用级数或专用 `Exp` 实现，直接除以极小 $\theta$ 会产生数值问题。对数映射执行反向操作，但在旋转角接近 $\pi$ 时存在分支和轴方向敏感性。

$SE(3)$ 的切空间增量写作 $\xi=[\rho^\top,\phi^\top]^\top\in\mathbb{R}^6$，其中 $\rho$ 是与平移相关的局部量，$\phi$ 是旋转向量。位姿误差不宜简单写成两个齐次矩阵逐元素相减；更有几何意义的主体误差为
\begin{equation}
e=\operatorname{Log}(T(q)^{-1}T_d)^\vee\in\mathbb{R}^6.
\label{eq:se3-pose-error}
\end{equation}
这里 $T(q)^{-1}T_d$ 先表达从当前位姿到目标位姿的相对变换，$\operatorname{Log}$ 再把它拉回局部切空间。左误差 $T_dT(q)^{-1}$ 与右误差 $T(q)^{-1}T_d$ 属于不同参考系；代码中的雅可比表达系必须与误差定义一致。Barfoot 对三维状态估计的处理同样强调，流形状态的扰动坐标是优化与不确定性表达的组成部分，而非实现细节\cite{barfoot2024state}。

\section{正运动学、速度运动学与雅可比}

串联机械臂正运动学把关节向量 $q\in\mathbb{R}^n$ 映射为末端位姿 $T(q)\in SE(3)$。在乘积指数形式下，每个关节用螺旋轴 $S_i\in\mathbb{R}^6$ 表示：
\begin{equation}
T(q)=e^{[S_1]q_1}e^{[S_2]q_2}\cdots e^{[S_n]q_n}M,
\label{eq:poe-forward-kinematics}
\end{equation}
$M$ 是零位构型的末端位姿。式~\eqref{eq:poe-forward-kinematics} 输出的位姿进入式~\eqref{eq:se3-pose-error}；对它局部求导得到雅可比，把关节速度映射为末端 twist：
\begin{equation}
V=J(q)\dot q,qquad J\in\mathbb{R}^{6\times n}.
\label{eq:velocity-jacobian}
\end{equation}
雅可比上三行和下三行究竟对应角速度还是线速度，取决于库的约定；空间雅可比与主体雅可比也表达在不同坐标系。只抄一个 $6\times n$ 数组而不记录参考系，会让控制方向在姿态变化后悄然错误。

二连杆平面臂的末端位置为
\begin{equation}
x=l_1\cos q_1+l_2\cos(q_1+q_2),\qquad
y=l_1\sin q_1+l_2\sin(q_1+q_2),
\end{equation}
因此
\begin{equation}
J(q)=
\begin{bmatrix}
-l_1\sin q_1-l_2\sin(q_1+q_2)&-l_2\sin(q_1+q_2)\\
l_1\cos q_1+l_2\cos(q_1+q_2)&l_2\cos(q_1+q_2)
\end{bmatrix}.
\label{eq:two-link-jacobian}
\end{equation}
取 $l_1=l_2=1$ m、$q_1=30^\circ$、$q_2=60^\circ$，得到末端 $(x,y)=(0.866,1.5)$ m，$J=[[-1.5,-1],[0.866,0]]$ m/rad。若期望末端速度 $[0.01,0]^\top$ m/s，解得有限关节速度；当 $q_2\to0$ 时，机械臂伸直，雅可比两列对某些任务方向失去独立性，最小奇异值趋近 0。

\section{逆运动学、阻尼伪逆与约束}

逆运动学不是一般意义上的矩阵求逆，而是寻找满足 $T(q)\approx T_d$ 且符合关节、碰撞和任务约束的 $q$。闭环迭代可写为
\begin{equation}
\Delta q=J^\#e,qquad q_{k+1}=\operatorname{integrate}(q_k,\alpha\Delta q),
\label{eq:ik-update}
\end{equation}
其中 $e$ 来自式~\eqref{eq:se3-pose-error}，$J^\#$ 是广义逆，$\alpha$ 是步长。Moore--Penrose 伪逆在奇异值分解 $J=U\Sigma V^\top$ 下将非零奇异值取倒数；小奇异值会把误差放大为巨大关节速度。阻尼最小二乘改为
\begin{equation}
J^\#_\lambda=J^\top(JJ^\top+\lambda^2I)^{-1},
\label{eq:damped-pseudoinverse}
\end{equation}
在弱方向上限制增益。$\lambda$ 增大使更新更稳但残差更大；它不能让不可达目标变得可达，也不能自动满足碰撞约束。

冗余机械臂还可利用零空间：
\begin{equation}
\dot q=J^\#V+(I-J^\#J)\dot q_0.
\label{eq:nullspace-ik}
\end{equation}
第一项完成主任务，第二项在局部零空间中优化关节居中、远离限位或避障。但数值阻尼使投影不再严格，多个任务的优先级也可能在奇异附近相互污染。需要硬约束时，应把逆运动学写成带边界的二次规划或非线性优化，而不是继续堆叠伪逆。

\section{官方实现导读：Pinocchio 的闭环逆运动学}

Pinocchio 是面向刚体运动学、动力学及其导数的开源库，其论文描述了正/逆动力学和解析导数等统一接口\cite{carpentier2019pinocchio}。官方 v3.4.0 逆运动学示例按“正运动学—李群误差—雅可比—阻尼求解—流形积分”执行，与图~\ref{fig:kinematic-compute-chain} 对应\cite{pinocchio340ik}。下面保留决定算法行为的核心路径，删除模型创建、显示与打印分支。

\begin{lstlisting}[language=Python,caption={Pinocchio v3.4.0 闭环逆运动学核心路径（BSD-2-Clause；据官方示例节选）},label={lst:pinocchio-clik}]
pinocchio.forwardKinematics(model, data, q)
iMd = data.oMi[JOINT_ID].actInv(oMdes)
err = pinocchio.log(iMd).vector
if norm(err) < eps:
    success = True
    break

J = pinocchio.computeJointJacobian(model, data, q, JOINT_ID)
J = -pinocchio.Jlog6(iMd.inverse()) @ J
v = -J.T @ solve(J @ J.T + damp * np.eye(6), err)
q = pinocchio.integrate(model, q, v * DT)
\end{lstlisting}

代码第 2 行把目标表达为当前关节系下的相对位姿，第 3 行取 $SE(3)$ 对数误差；第 8 行的 `Jlog6` 把运动学雅可比转换到误差函数的微分；第 9 行对应式~\eqref{eq:damped-pseudoinverse}，但示例变量 `damp` 直接加到 $JJ^\top$，因此它对应公式中的 $\lambda^2$；第 10 行使用库的流形积分，而不是假设所有构型都能用普通向量相加。

\sourceanchor{仓库 \path{stack-of-tasks/pinocchio}；许可证 BSD-2-Clause；release \texttt{v3.4.0}；官方文档 \path{doc/d-practical-exercises/3-invkine.html} 与 \path{examples/inverse-kinematics.py}；符号 \texttt{forwardKinematics}、\texttt{computeJointJacobian}、\texttt{Jlog6}、\texttt{integrate}。2026-07-14 远端 commit 查询受网络连接失败影响，当前以不可变 release tag 固定；冻结前补记解析后的 commit。}

\section{失效诊断、动作接口与知识小结}

逆运动学失败应按可观测信号分类。误差始终很大且关节触限，首先检查目标是否可达；最小奇异值下降、关节速度上升，说明接近机构奇异；误差方向突然翻转，检查四元数符号与对数分支；位置收敛而姿态不收敛，检查线速度/角速度权重和单位；数值收敛但碰撞，说明优化问题没有包含环境约束。把所有失败都归因于“初值不好”会掩盖接口设计错误。

高层策略输出笛卡尔增量时，至少要声明参考系、平移和旋转单位、动作周期、限幅与控制语义。输出关节目标时，则要声明关节顺序、本体版本和位置/速度/力矩模式。第 22 章会比较不同动作表示，但无论动作来自规则、规划器还是 VLA，最终都必须通过这些接口约定进入第 6 章的动力学与控制闭环。

\authorjudgement{李群与雅可比的价值不在于术语先进，而在于它们把有限位姿、局部误差、速度映射和数值迭代连成同一条可审计链。模型可以学习逆运动学的近似，但只要部署仍需要坐标变换、限位、碰撞和实时诊断，这条几何链就不能被一句“端到端输出动作”省略。}

\chapter{动力学、接触与柔顺控制}
\label{chap:dynamics-contact-compliance}

\section{学习目标、前置知识与问题边界}

读完本章，读者应能从机械臂的能量出发写出关节空间动力学，解释接触互补条件和摩擦锥的物理意义，推导任务空间阻抗控制律，并把一项柔顺插装任务拆成观测、控制、阈值与恢复四个相互独立的设计问题。前置知识是第 5 章中的齐次变换、雅可比矩阵和速度映射。本章不把“接触成功”归因于某一个控制器：孔位估计错误、工件公差、末端标定、传感延迟和恢复策略都可能决定结果。

刚性自由空间轨迹跟踪要求机器人逼近给定位置；接触操作则要求机器人在几何不确定时限制交互力。二者的差别不是“是否加一个力传感器”，而是受控对象从位姿误差扩展为机器人与环境组成的耦合动力系统。Hogan 的阻抗控制把目标写成力与运动之间的动态关系，而不是把位置或力之一孤立地设为唯一输出\cite{hogan1985impedance}；Khatib 的操作空间公式进一步给出了任务坐标中的惯量和力映射\cite{khatib1987operational}。

\section{从能量到关节空间动力学}

设 $q\in\mathbb{R}^{n}$ 为关节位置，$\dot q$、$\ddot q$ 分别为速度和加速度。机械臂的动能和势能写为
\begin{equation}
T(q,\dot q)=\frac{1}{2}\dot q^\top M(q)\dot q,\qquad V(q),
\end{equation}
其中 $M(q)$ 是对称正定惯量矩阵。对拉格朗日量 $L=T-V$ 应用 Euler--Lagrange 方程，可得
\begin{equation}
M(q)\ddot q+C(q,\dot q)\dot q+g(q)+\tau_f
=\tau+J_c(q)^\top\lambda .
\label{eq:joint-dynamics}
\end{equation}
$\tau$ 是执行器力矩，$g(q)=\partial V/\partial q$，$\tau_f$ 汇集关节摩擦等未建模项，$J_c^\top\lambda$ 把接触点的广义力映射到关节空间。$C$ 的矩阵表示不是唯一的，但应满足 $\dot M-2C$ 反对称；该性质使无耗散、无外力系统的能量变化与输入功率一致，是许多稳定性推导的关键\cite{featherstone2008rigid,lynch2017modern}。

式中各项的量纲都必须是力矩：$M\ddot q$ 是产生关节加速度所需的惯性力矩，$C\dot q$ 汇集由速度引起的科氏与离心力矩，$g$ 抵消重力，$\tau_f$ 表示摩擦和传动损失。接触项来自虚功等价：接触点虚位移满足 $\delta x_c=J_c\delta q$，外力做功
\begin{equation}
\delta W=\lambda^\top\delta x_c
=\lambda^\top J_c\delta q
=(J_c^\top\lambda)^\top\delta q,
\end{equation}
所以与接触力等价的关节广义力必为 $J_c^\top\lambda$。这也解释了为什么末端力传感器读数不能直接加到关节力矩上：必须先统一坐标系、符号和力矩参考点，再乘雅可比转置。

反对称性质还给出一个有用的能量检查。令机械能 $E=T+V$，把式\eqref{eq:joint-dynamics}左乘 $\dot q^\top$，利用 $\dot q^\top(\dot M-2C)\dot q=0$，可得
\begin{equation}
\dot E=\dot q^\top\tau+(J_c\dot q)^\top\lambda-\dot q^\top\tau_f.
\label{eq:robot-power-balance}
\end{equation}
右侧依次是执行器输入功率、环境通过接触注入的功率和摩擦耗散。仿真器若在无输入、无接触、无摩擦时让 $E$ 持续增长，通常意味着积分器、$C$ 的实现或坐标变换存在错误。

式\eqref{eq:joint-dynamics}对应三类计算。逆动力学已知 $(q,\dot q,\ddot q)$ 求 $\tau$，用于前馈补偿；正动力学已知 $(q,\dot q,\tau)$ 求 $\ddot q$，用于仿真和预测；参数辨识则由多段运动和力矩数据估计质量、质心、惯量及摩擦参数。三者不能混用：一个能在低速轨迹上拟合力矩的模型，未必能预测碰撞冲量，因为后者还依赖接触几何、恢复系数和采样时间。

\subsection{二连杆算例：耦合项如何进入控制}

对平面二连杆，若忽略摩擦，惯量矩阵可写成
\begin{equation}
M(q)=\begin{bmatrix}
a+2b\cos q_2 & d+b\cos q_2\\
d+b\cos q_2 & d
\end{bmatrix},
\end{equation}
其中 $a,d$ 由连杆质量、长度和转动惯量组成，$b=m_2l_1l_{c2}$。取 $a=3.2$、$b=0.6$、$d=0.5\ \mathrm{kg\,m^2}$，当 $q_2=0$ 时
$M=\left[\begin{smallmatrix}4.4&1.1\\1.1&0.5\end{smallmatrix}\right]$；当 $q_2=\pi/2$ 时变为
$M=\left[\begin{smallmatrix}3.2&0.5\\0.5&0.5\end{smallmatrix}\right]$。同一关节加速度命令在不同姿态需要不同力矩。只使用常数增益而没有模型补偿时，姿态改变会表现为响应速度和超调变化；这不是神经网络特有的“分布偏移”，而是机械系统本身的构型依赖。

例如忽略 $C$ 与 $g$，要求 $\ddot q=[1,0]^\top\ \mathrm{rad/s^2}$，伸直姿态需要惯性力矩 $[4.4,1.1]^\top\ \mathrm{N\,m}$，而 $q_2=\pi/2$ 时只需 $[3.2,0.5]^\top\ \mathrm{N\,m}$。第二关节虽然没有加速度，仍因非对角项收到 $1.1$ 或 $0.5\ \mathrm{N\,m}$；这就是动态耦合。真实逆动力学还必须把 $C\dot q$、$g$ 和期望接触力加入，算例只隔离惯量项。

\section{单边接触、摩擦与混合动力学}

令 $\phi(q)$ 为两物体的法向间隙，$\phi>0$ 表示分离，$\phi=0$ 表示接触。刚性、不可穿透接触的理想化法向条件是
\begin{equation}
0\leq \lambda_n\ \perp\ \phi(q)\geq 0,
\label{eq:complementarity}
\end{equation}
即 $\lambda_n\geq0$、$\phi\geq0$ 且 $\lambda_n\phi=0$。它排除了“尚未接触却产生推力”和“接触时仍相互穿透”两种状态。切向力满足库仑摩擦锥
\begin{equation}
\|\lambda_t\|_2\leq \mu\lambda_n.
\label{eq:friction-cone}
\end{equation}
严格位于锥内可以粘着；达到边界后可能滑动，方向与相对切向速度相反。数值算法常把圆锥多面体化以转成线性互补问题，但边数越少，摩擦方向误差越大。

互补式只允许两种法向模式：分离时 $\phi>0,\lambda_n=0$；承压接触时 $\phi=0,\lambda_n\ge0$。$\phi=0,\lambda_n=0$ 是即将建立或断开的边界状态，而 $\phi>0,\lambda_n>0$ 和 $\phi<0$ 都不合法。摩擦锥还要结合切向速度理解：$\|\lambda_t\|<\mu\lambda_n$ 时，约束求解器可以选择足以阻止滑动的切向力；到锥面时，若仍有相对运动，摩擦方向与切向速度相反。式\eqref{eq:friction-cone}给的是可行集合，不是独立计算 $\lambda_t$ 的控制律。

\begin{figure}[H]
\centering
\includegraphics[width=\textwidth]{v0.3-06-接触互补与摩擦锥.pdf}
\caption{法向互补的合法/非法状态与库仑摩擦锥截面。作者依据刚性接触模型绘制；图为概念解释，不是实验数据\cite{featherstone2008rigid,lynch2017modern}。}
\label{fig:contact-complementarity-friction}
\end{figure}

接触建立或断开会改变约束集合，碰撞还会使速度发生跃迁。因此接触操作是混合系统：连续阶段用式\eqref{eq:joint-dynamics}演化，离散事件更新接触模式或施加冲量。柔顺材料和有限采样控制器常用罚函数 $\lambda_n=k_p[-\phi]_++d_p[-\dot\phi]_+$ 近似接触；它容易实现，却把允许穿透量、刚度和积分步长绑定在一起。减小步长并不自动消除模型偏差，增大刚度也可能造成数值发散。

\begin{table}[H]
\centering
\caption{接触建模选择及其代价}
\label{tab:contact-models}
\begin{tabularx}{\textwidth}{p{0.18\textwidth}YYY}
\toprule
方法 & 主要变量 & 优点 & 典型失效\\
\midrule
刚性互补 & 间隙、接触力/冲量 & 不允许静态穿透，接触逻辑明确 & 求解不光滑，多接触时计算昂贵\\
罚函数 & 穿透量、相对速度 & 易并入常微分方程和自动微分 & 刚度与步长耦合，参数缺乏真实对应\\
可变形接触 & 形变场或局部材料参数 & 能表示软物体和接触斑 & 参数多，识别与实时计算困难\\
数据驱动残差 & 解析模型加学习修正 & 可补偿稳定的系统误差 & 外推不可控，不能替代不可穿透约束\\
\bottomrule
\end{tabularx}
\end{table}

\section{从关节力矩到任务空间阻抗}

末端速度满足 $\dot x=J(q)\dot q$。忽略奇异点并选定任务子空间后，操作空间惯量为
\begin{equation}
\Lambda(q)=\left(JM^{-1}J^\top\right)^{-1}.
\label{eq:operational-inertia}
\end{equation}
$\Lambda$ 说明同一笛卡尔方向在不同姿态呈现不同“等效质量”。若希望位姿误差 $e=x-x_d$ 在外力 $F_{ext}$ 下满足
\begin{equation}
M_d\ddot e+D_d\dot e+K_de=F_{ext},
\label{eq:desired-impedance}
\end{equation}
则最简任务空间弹簧--阻尼力为 $F_c=-K_de-D_d\dot e$，再经 $\tau=J^\top F_c+C\dot q+g$ 映射到关节。完整的操作空间控制还应补偿 $\Lambda$、任务空间科氏/重力项，并处理零空间任务\cite{khatib1987operational}。

式\eqref{eq:operational-inertia}可以从关节加速度关系看出。暂忽略 $\dot J\dot q$、重力和科氏项，若关节力矩由任务力产生，$\tau=J^\top F$，则 $\ddot x=JM^{-1}J^\top F=\Lambda^{-1}F$，因此 $F=\Lambda\ddot x$。当 $J$ 行满秩时该逆存在；接近奇异位形时，某些方向的等效惯量会急剧增大，直接求逆放大噪声。工程实现通常采用阻尼伪逆、限制任务方向或避开奇异区，而不是假定公式在所有构型都数值良好。

阻抗控制给定位移到力的映射；导纳控制则测量外力，再由虚拟动力学积分得到位置或速度命令。高带宽力矩接口常适合外环阻抗，只有位置接口的工业机械臂更常用导纳外环。两者的稳定性都受环境刚度、采样延迟和内环动态影响\cite{deschutter1997force}。“把 $K_d$ 调低就安全”并不成立：过软会造成定位误差和长时间摩擦，过低阻尼会在接触切换时振荡。

\begin{figure}[H]
\centering
\includegraphics[width=\textwidth]{v0.3-06-阻抗与导纳接口.pdf}
\caption{阻抗与导纳的信号方向和执行器接口。作者依据经典力控制分类绘制\cite{hogan1985impedance,deschutter1997force}。}
\label{fig:impedance-admittance-interface}
\end{figure}

图\ref{fig:impedance-admittance-interface}中的差别决定代码落点。阻抗外环从位置/速度误差计算期望交互力或关节力矩，依赖可控力矩接口；导纳外环从测得外力积分出新的位置/速度命令，可包在厂商位置伺服器之外。后者并不天然更稳定：力传感噪声会被积分，位置内环延迟又会改变虚拟质量、阻尼和环境刚度组成的闭环极点。

\subsection{一维参数算例}

设期望接触方向的虚拟质量 $m_d=2\ \mathrm{kg}$、刚度 $k_d=800\ \mathrm{N/m}$。临界阻尼为 $d_d=2\sqrt{m_dk_d}=80\ \mathrm{N\,s/m}$，固有频率 $\omega_n=\sqrt{k_d/m_d}=20\ \mathrm{rad/s}$。若环境位置误差为 $2\ \mathrm{mm}$，静态弹簧力约为 $1.6\ \mathrm{N}$；若标定误差增至 $10\ \mathrm{mm}$，则为 $8\ \mathrm{N}$。控制律没有改变，但几何误差把交互力放大了五倍。这就是为什么接触控制章节必须同时讨论标定、阈值和回退。

\section{官方实现导读：方程如何落到实时回调}

libfranka 0.15.0 的官方示例在 1 个实时回调中读取科氏项、末端雅可比和关节速度，构造弹簧--阻尼任务力矩，并把科氏补偿相加\cite{libfranka015}。以下片段保留其关键计算，省略连接、姿态误差和异常处理：

\begin{lstlisting}[language=C++,caption={libfranka 0.15.0 笛卡尔阻抗示例的关键力矩计算（Apache-2.0）}]
auto coriolis_array = model.coriolis(robot_state);
auto jacobian_array = model.zeroJacobian(
    franka::Frame::kEndEffector, robot_state);
Eigen::Map<const Eigen::Matrix<double, 7, 1>>
    coriolis(coriolis_array.data());
Eigen::Map<const Eigen::Matrix<double, 6, 7>>
    jacobian(jacobian_array.data());
Eigen::Map<const Eigen::Matrix<double, 7, 1>>
    dq(robot_state.dq.data());
tau_task = jacobian.transpose()
         * (-stiffness * error - damping * (jacobian * dq));
tau_d = tau_task + coriolis;
\end{lstlisting}
\sourceanchor{仓库 \texttt{frankarobotics/libfranka}；tag \texttt{0.15.0}；文件 \texttt{examples/cartesian\_impedance\_control.cpp}；许可证 Apache-2.0；符号 \texttt{impedance\_control\_callback}。}

代码与式\eqref{eq:desired-impedance}之间存在三个值得检查的差异。第一，示例明确“不做惯量整形”，所以 $M_d$ 不是任意指定的常数矩阵；实际闭环惯量随构型变化。第二，示例的旋转误差由四元数构造，不能把欧氏角度误差原样复制。第三，官方文档警告该示例把碰撞阈值设得较高，并要求操作者握住急停。示例能解释算法，不等于通过生产安全验证。

\section{完整案例：带搜索与回退的柔顺插装}

考虑圆柱销插入孔。已知末端标称孔位 $\hat x_h$，但横向误差可能达到 $\pm2\ \mathrm{mm}$；六维力/力矩传感器提供 $w=[f_x,f_y,f_z,m_x,m_y,m_z]$。任务不是“沿 $z$ 轴下降”，而是一个带观测判据的状态机：

\begin{figure}[H]
\centering
\includegraphics[width=\textwidth]{v0.3-06-柔顺插装闭环.png}
\caption{柔顺插装的观测、控制与异常恢复链。作者自绘；概念依据为阻抗/力控制文献\cite{hogan1985impedance,deschutter1997force}，不复制论文原图。}
\label{fig:compliant-insertion-loop}
\end{figure}

\begin{enumerate}
  \item \textbf{接近：}在自由空间移动到孔上方，限制速度和动能；检查视觉置信度、工具标定版本和工件是否到位。
  \item \textbf{探触：}沿 $-z$ 低速移动。当 $f_z>f_{contact}$ 持续 $N$ 个采样周期时确认接触，避免单点噪声触发。
  \item \textbf{搜索：}保持法向期望力或低法向刚度，在 $xy$ 平面执行螺旋/栅格微动；横向力和力矩用于判断卡边。
  \item \textbf{插入：}检测到 $z$ 位移增加且横向力下降后，切换到较高横向刚度和受限轴向速度。
  \item \textbf{验证：}以深度、末端残余力和工件存在信号共同判定成功；单一位置阈值可能把工件变形误判为到位。
  \item \textbf{回退：}若 $\|f_{xy}\|>f_{jam}$、$|m_{xy}|>m_{jam}$ 或搜索超时，立即停止下压，沿已知安全方向撤回，重新观测孔位；重复次数超过上限则转人工。
\end{enumerate}

\begin{lstlisting}[language=Python,caption={教学伪代码：插装状态机；非任何厂商源码}]
while not terminal:
    wrench = low_pass(read_wrench())
    if state == APPROACH:
        command_pose(pre_insert)
    elif state == TOUCH and persistent(wrench.fz > f_contact, N):
        state = SEARCH
    elif state == SEARCH:
        command_impedance(spiral_xy(), vz=-v_probe, K=K_search)
        if insertion_signature(depth, wrench): state = INSERT
        if jam(wrench) or timeout(): state = RETRACT
    elif state == RETRACT:
        stop_downward_motion(); retract(); reacquire_hole()
\end{lstlisting}

这段伪代码故意把控制命令与判据分开。$f_{contact}$、$f_{jam}$ 和滤波窗口不能凭经验永久写死：它们需要由传感器噪声、正常接触分布、工件公差和允许风险共同标定。部署日志至少应保存原始/滤波力、位姿、状态转换原因、控制器参数和失败图像，否则失败样本无法回流到几何估计或策略训练。

\section{知识小结与综述判断}

\textbf{知识小结。}关节空间动力学把执行器力矩、惯性、科氏/离心、重力和接触广义力置于同一方程；单边接触由间隙、法向力和互补条件描述，摩擦锥区分粘着与滑动；任务空间惯量说明末端方向上的动态响应随构型变化；阻抗与导纳控制分别规定运动--力关系及力到运动的外环映射。柔顺插装需要状态机、阈值和恢复链，不能只给出一条控制方程。

\authorjudgement{经典阻抗和操作空间理论仍是接触操作的主干，不是等待被 VLA 淘汰的“传统模块”。学习策略最有价值的切入点通常是难以建模的搜索轨迹、接触状态估计和参数调节；安全力限、实时内环、急停与确定性回退仍应由可验证的控制层承担。若一项工作只报告成功率而不报告接触力、超时、卡死和人工恢复，它不足以证明接触操作能力。}

\chapter[运动规划、轨迹优化与最优控制]{运动规划、轨迹优化\\与最优控制}
\label{chap:planning-optimal-control}

\section{学习目标、前置知识与问题边界}

本章把“机器人怎样到达目标”拆成三个不同问题。运动规划寻找不碰撞的几何路径；轨迹生成给路径加上时间、速度和加速度；反馈控制在模型误差与扰动下持续修正执行。读完本章，读者应能写出三者各自的决策变量和约束，推导有限时域 LQR 的反向递推，解释轨迹优化与 MPC 的联系，并为一个受约束任务选择可验证的组合，而不是把 A*、RRT、LQR 与 MPC 当作并列名词。

前置知识是第~\ref{chap:math-computation}章的约束优化、第~\ref{chap:kinematics-lie}章的构型与雅可比，以及第~\ref{chap:dynamics-contact-compliance}章的动力学。高层任务规划和 VLA 只向本章提供目标、代价或参考轨迹；它们不能替代碰撞约束、执行器边界和毫秒级反馈。

\begin{figure}[H]
\centering
\setlength{\tabcolsep}{5pt}
\begin{tabular}{c@{\ $\longrightarrow$\ }c@{\ $\longrightarrow$\ }c}
\fbox{\begin{minipage}{0.25\textwidth}\centering 几何路径\\$q(s)\in\mathcal C_{\mathrm{free}}$\\不含时间\end{minipage}} &
\fbox{\begin{minipage}{0.25\textwidth}\centering 带时轨迹\\$q(t),\dot q(t),\ddot q(t)$\\满足运动边界\end{minipage}} &
\fbox{\begin{minipage}{0.25\textwidth}\centering 闭环控制\\$u_t=\pi(x_t,r_t)$\\抵抗扰动\end{minipage}}
\end{tabular}
\caption{从几何可行到动态可执行的三层接口。作者教学绘制；箭头表示上游结果成为下游参考，不表示一次求解即可完成任务。}
\label{fig:path-trajectory-feedback}
\end{figure}

\section{构型空间、碰撞与采样规划}

令机器人构型为 $q\in\mathcal C$，碰撞构型集合为 $\mathcal C_{\mathrm{obs}}$，则自由空间为 $\mathcal C_{\mathrm{free}}=\mathcal C\setminus\mathcal C_{\mathrm{obs}}$。几何规划寻找连续曲线
\begin{equation}
q:[0,1]\rightarrow\mathcal C_{\mathrm{free}},\qquad q(0)=q_s,\quad q(1)\in\mathcal Q_g.
\label{eq:geometric-path}
\end{equation}
这里的碰撞检查不是图形显示功能，而是可行性判定器。机械臂的一个末端位姿通常对应多个关节构型；只在笛卡尔空间画直线，既可能穿过障碍，也可能越过关节限位或奇异区。

离散栅格或低维图上可用 Dijkstra/A*；高维连续空间常用 PRM、RRT 及其变体。RRT 反复采样 $q_{\mathrm{rand}}$，从树中最近节点向样本延伸并做局部碰撞检查。它在常见条件下具有概率完备性：若存在具有正裕量的可行路径，样本数增加时找到路径的概率趋近于一；但原始 RRT 不保证路径代价趋近最优。RRT* 增加邻域选父与重连，使代价在相应假设下几乎必然趋近最优\cite{karaman2011sampling}。

“渐近最优”不等于有限时间最好。狭窄通道的有效体积很小，均匀采样可能长期碰不到入口；碰撞检测往往占据主要计算量；错误的机器人几何、膨胀半径或场景时间戳会让规划器对错误世界给出数学上正确的路径。因而应记录采样预算、碰撞模型版本、最小净空和终止原因，而不只记录“规划成功”。

\section{轨迹优化：从一条初值到约束局部解}

将轨迹离散为 $Q=(q_0,\ldots,q_N)$，典型目标为
\begin{equation}
\min_Q\;\underbrace{\frac12\sum_{k=0}^{N-1}\|q_{k+1}-q_k\|_W^2}_{\text{平滑/控制代价}}
+\alpha\underbrace{\sum_{k=0}^{N}c_{\mathrm{obs}}(q_k)}_{\text{障碍代价}},
\quad\text{s.t. }q_0=q_s,\;q_N\in\mathcal Q_g,\;g(Q)\le0.
\label{eq:trajectory-optimization}
\end{equation}
$W$ 决定不同关节和时间差分的尺度；$c_{\mathrm{obs}}$ 常由有符号距离场构造；$g$ 收纳关节、速度、加速度和任务约束。CHOMP 使用轨迹空间的平滑度量和障碍代价梯度更新整条轨迹\cite{ratliff2009chomp}；TrajOpt 把非凸问题逐次凸化，并强调连续时间碰撞检查，避免只检查离散路点而在两点之间穿越障碍\cite{schulman2014motion}。

式~\eqref{eq:trajectory-optimization} 是局部优化：初值位于障碍哪一侧，常决定最终同伦类。一个常用工程组合是先用采样规划给出拓扑上可行的粗路径，再由轨迹优化增加净空和平滑性，最后做时间参数化。若直接对一条穿障碍直线做梯度下降，距离场梯度在障碍内部可能不稳定，甚至把轨迹推入错误局部极小值。

路径参数 $s$ 与时间 $t$ 不能混用。给定 $q(s)$，有
\begin{equation}
\dot q=q'(s)\dot s,\qquad
\ddot q=q''(s)\dot s^2+q'(s)\ddot s.
\label{eq:time-parameterization}
\end{equation}
速度、加速度、力矩约束因此转化为对 $\dot s,\ddot s$ 的限制。简单地把相邻路点按固定周期发送给控制器，会在曲率大处产生速度/加速度跳变；几何无碰撞并不意味着动态可执行。

\section{LQR 与 DDP：为什么要反向传播一个二次值函数}

考虑线性离散系统 $x_{k+1}=A_kx_k+B_ku_k$，有限时域二次代价
\begin{equation}
J=x_N^\top Q_Nx_N+\sum_{k=0}^{N-1}
\left(x_k^\top Q_kx_k+u_k^\top R_ku_k\right),
\quad Q_k\succeq0,\;R_k\succ0.
\label{eq:lqr-cost}
\end{equation}
假设后续最优代价为 $V_{k+1}(x)=x^\top P_{k+1}x$，把动力学代入并对 $u_k$ 求极小，得到
\begin{align}
K_k&=(R_k+B_k^\top P_{k+1}B_k)^{-1}B_k^\top P_{k+1}A_k,\\
u_k^*&=-K_kx_k,\\
P_k&=Q_k+A_k^\top P_{k+1}A_k-A_k^\top P_{k+1}B_kK_k,
\quad P_N=Q_N.
\label{eq:riccati-recursion}
\end{align}
这不是凭经验选择反馈增益，而是动态规划下二次值函数的系数递推。$R$ 越大，控制动作越昂贵；$Q$ 越大，状态偏差越昂贵。二者量纲必须经缩放后再比较。

以标量系统 $x_{k+1}=x_k+u_k$、两步时域、$Q_0=Q_1=0$、$Q_2=1$、$R_0=R_1=1$ 为例。由 $P_2=1$ 得 $K_1=1/2,P_1=1/2$；再得 $K_0=(1+1/2)^{-1}(1/2)=1/3$。若 $x_0=3$，则 $u_0=-1,x_1=2$，随后 $u_1=-1,x_2=1$。有限时域和控制代价使系统不必在第一步把误差清零。

非线性系统可围绕名义轨迹线性化并反复执行上述反向/正向过程。iLQR/DDP 的核心是对动力学和代价作局部展开，再用线搜索接受新轨迹；控制受限 DDP 还在反向步显式处理动作边界\cite{tassa2014control}。它们仍是局部方法：接触模式突变、坏初值和不准确二阶近似都可能破坏收敛。

\section{MPC：把有限时域优化变成反馈}

模型预测控制在每个时刻求解
\begin{equation}
\begin{aligned}
\min_{x_{0:N},u_{0:N-1}}\quad&\ell_N(x_N)+\sum_{k=0}^{N-1}\ell(x_k,u_k)\\
\text{s.t.}\quad&x_{k+1}=f(x_k,u_k),\quad x_k\in\mathcal X,\quad u_k\in\mathcal U,\\
&x_0=\hat x_t,
\end{aligned}
\label{eq:mpc}
\end{equation}
只执行最优序列的第一个动作，收到新状态估计后平移时域并重算。与离线轨迹优化相比，MPC 通过重规划形成反馈，并可显式处理状态、动作与终端约束；其稳定性并非来自“每次都优化”这句话，而依赖终端代价/终端集合、递归可行性和模型假设\cite{mayne2000mpc}。

预测时域短会近视，时域长则增加求解时间和模型误差累积。若求解器超过控制周期，系统使用的是过期状态对应的动作；若约束突然不可行，必须有软约束、备份控制器或安全停车，而不能把上一次解无限复用。

\begin{lstlisting}[caption={滚动时域受约束控制的教学伪代码；非厂商或开源项目源码}]
while task_active:
    x_hat, stamp = estimator.latest()
    if now() - stamp > freshness_limit: safe_stop()
    problem.set_initial_state(x_hat)
    problem.set_reference(local_path, goal)
    sol = solver.solve(warm_start=shift(previous_sol))
    if not sol.feasible or sol.solve_time > deadline:
        execute(backup_controller(x_hat)); continue
    execute(rate_limit(sol.u[0]))
    previous_sol = sol
\end{lstlisting}

\section{同一受约束任务的接口比较}

设七轴机械臂要从货架外移动到箱内抓取位，场景有狭窄入口、关节速度/力矩限制和移动中的人。表~\ref{tab:planning-control-comparison} 的方法不是互斥竞品，而是处理不同变量与时间尺度。

\begin{table}[H]
\centering
\caption{同一任务中规划与控制方法的职责、证据与失效}
\label{tab:planning-control-comparison}
\begin{tabularx}{\textwidth}{p{0.15\textwidth}YYYY}
\toprule
方法 & 决策变量 & 主要约束/证据 & 是否反馈 & 典型失效 \\
\midrule
图搜索 & 离散节点路径 & 栅格连通与边代价 & 否 & 维数灾难、离散化漏障碍 \\
RRT/RRT* & 连续构型树 & 局部连接与碰撞检查 & 否 & 狭窄通道慢、有限预算路径差 \\
轨迹优化 & 整段路点/控制 & 距离场、平滑和边界 & 通常否 & 初值敏感、局部极小 \\
LQR/iLQR & 反馈增益/名义轨迹 & 局部动力学与二次代价 & 是 & 线性化失真、动作饱和 \\
MPC & 有限时域状态/动作 & 动力学、状态和动作集合 & 是 & 超时、不可行、模型偏差 \\
\bottomrule
\end{tabularx}
\end{table}

一个可审计组合是：RRT* 找到进入箱体的可行同伦类；轨迹优化提高净空并平滑；时间参数化满足速度/加速度边界；MPC 依据第~\ref{chap:state-estimation-slam}章的最新状态跟踪并避开动态障碍；第~\ref{chap:dynamics-contact-compliance}章的阻抗控制处理抓取前后的接触。若状态估计过期或动态障碍使问题不可行，系统退回保持位或安全撤离。这里每一层都有独立退出条件，因而能定位失败究竟发生在世界模型、几何可行性、动态可行性还是执行反馈。

\section{知识小结与综述判断}

知识小结：路径只描述构型随路径参数的变化，轨迹还包含时间导数，控制则把状态误差闭环到动作。采样规划擅长寻找非凸自由空间中的可行同伦类；轨迹优化擅长从初值改进平滑性和约束；LQR/DDP 用局部值函数得到反馈结构；MPC 以滚动时域把优化转成约束反馈。

综述判断：真实系统不应争论“规划还是学习”，而应明确学习模块输出的是目标、代价、初值、模型还是动作。只要机器人仍受碰撞、执行器和实时性约束，规划与控制的可行性检查就有独立价值；但它们的结论只对输入状态和模型有效，不能掩盖感知延迟、场景错误或接触建模缺失。

\chapter{状态估计、传感器融合与 SLAM}
\label{chap:state-estimation-slam}

\section{学习目标、前置知识与问题设定}

控制器读到的不是“真实状态”，而是由模型、传感器和时间戳共同生成的估计。读完本章，读者应能从贝叶斯递推推导卡尔曼预测与校正，解释 EKF、UKF、粒子滤波和因子图的假设差异，说明 SLAM 中位姿与地图为何耦合，并把时间同步、外参漂移与状态新鲜度写成可监控接口。

离散状态空间模型写为
\begin{equation}
x_k=f(x_{k-1},u_{k-1})+w_{k-1},\qquad
z_k=h(x_k)+v_k,
\label{eq:nonlinear-state-observation}
\end{equation}
$x_k$ 可含位姿、速度、IMU 偏置和地图路标，$z_k$ 是相机、雷达、编码器或力觉观测；$w,v$ 分别是过程和观测噪声。式中下标也是时间语义：若 $z_k$ 实际对应旧时刻却按当前状态融合，正确的矩阵运算仍会得到错误估计。

\begin{figure}[H]
\centering
\setlength{\tabcolsep}{3pt}
\begin{tabular}{c@{\ $\longrightarrow$\ }c@{\ $\longrightarrow$\ }c@{\ $\longrightarrow$\ }c}
\fbox{\begin{minipage}{0.17\textwidth}\centering 上一后验\\$\hat x^+,P^+$\end{minipage}} &
\fbox{\begin{minipage}{0.15\textwidth}\centering 模型预测\\$f,Q$\end{minipage}} &
\fbox{\begin{minipage}{0.17\textwidth}\centering 创新检验\\$r_k,S_k$\end{minipage}} &
\fbox{\begin{minipage}{0.17\textwidth}\centering 当前后验\\$\hat x^+,P^+$\\附时间戳\end{minipage}}
\end{tabular}
\caption{递归状态估计的数据流。作者教学绘制；创新门控与时间戳检查是估计器接口的一部分，不是日志附属项。}
\label{fig:state-estimation-loop}
\end{figure}

\section{线性高斯模型与卡尔曼滤波}

对线性模型
\begin{equation}
x_k=F_{k-1}x_{k-1}+B_{k-1}u_{k-1}+w_{k-1},\quad
z_k=H_kx_k+v_k,
\end{equation}
令 $w\sim\mathcal N(0,Q)$、$v\sim\mathcal N(0,R)$ 且相互独立。高斯分布经线性变换仍为高斯，因此预测为
\begin{align}
\hat x_{k|k-1}&=F_{k-1}\hat x_{k-1|k-1}+B_{k-1}u_{k-1},\\
P_{k|k-1}&=F_{k-1}P_{k-1|k-1}F_{k-1}^\top+Q_{k-1}.
\label{eq:kf-predict}
\end{align}
观测创新及其协方差为 $r_k=z_k-H_k\hat x_{k|k-1}$、$S_k=H_kP_{k|k-1}H_k^\top+R_k$。最小化更新后误差协方差可得\cite{kalman1960new}
\begin{align}
K_k&=P_{k|k-1}H_k^\top S_k^{-1},\\
\hat x_{k|k}&=\hat x_{k|k-1}+K_kr_k,\\
P_{k|k}&=(I-K_kH_k)P_{k|k-1}.
\label{eq:kf-update}
\end{align}
$K_k$ 不是人工信任系数，而是先验不确定性与观测不确定性的矩阵权衡。数值实现宜使用线性求解代替显式 $S^{-1}$，并可用 Joseph 形式
\begin{equation}
P_{k|k}=(I-KH)P^-(I-KH)^\top+KRK^\top
\end{equation}
减轻有限精度造成的非对称或非正定。

一维算例最能说明增益含义。先验位置为 $8\ \mathrm m$、方差 $P^-=4\ \mathrm{m^2}$；测量为 $9\ \mathrm m$、方差 $R=1\ \mathrm{m^2}$，且 $H=1$。则 $K=4/(4+1)=0.8$，后验 $\hat x^+=8+0.8(9-8)=8.8\ \mathrm m$，方差 $P^+=(1-0.8)4=0.8\ \mathrm{m^2}$。若把测量方差错写成 $0.01$，滤波器会几乎完全追随传感器；这不是“融合更灵敏”，而是噪声模型失配。

可观测性决定观测序列能否区分内部状态。线性时不变系统的可观测矩阵
\begin{equation}
\mathcal O=\begin{bmatrix}H^\top&(HF)^\top&\cdots&(HF^{n-1})^\top\end{bmatrix}^\top
\end{equation}
须满列秩才能从有限序列重构全部状态。滤波器能输出一个数，不代表该状态可观；不可观方向的协方差若被错误压小，会形成危险的“自信漂移”。

\section{非线性滤波：线性化、sigma 点与样本}

EKF 在当前估计附近以雅可比 $F_k=\partial f/\partial x$、$H_k=\partial h/\partial x$ 线性化，再使用式~\eqref{eq:kf-predict}--\eqref{eq:kf-update}。它计算紧凑，但线性化点错误、姿态参数化不连续或强非线性会造成不一致。误差状态滤波把姿态等流形状态保留在名义值上，只在切空间估计小误差；第~\ref{chap:kinematics-lie}章的 $SO(3)/SE(3)$ 正是为此提供接口\cite{sola2017quaternion,barfoot2024state}。

UKF 选取一组确定性 sigma 点，使其样本均值与协方差匹配当前高斯分布；每个点通过非线性 $f,h$ 后重新加权均值与协方差，从而避免显式雅可比\cite{julier2004ukf}。它并没有消除高斯假设，也不能自动处理多峰后验。粒子滤波则用带权样本近似后验，可表示多峰和强非线性，但维数升高时需要大量粒子，重采样还会造成样本贫化。

\begin{table}[H]
\centering
\caption{递归估计与批量平滑方法的假设和失效}
\label{tab:estimators-comparison}
\begin{tabularx}{\textwidth}{p{0.14\textwidth}YYYY}
\toprule
方法 & 后验表示 & 非线性处理 & 主要优势 & 典型失效 \\
\midrule
KF & 单高斯 & 线性模型 & 闭式、快速、可解释 & 模型或噪声非线性/非高斯 \\
EKF & 局部高斯 & 一阶雅可比 & 适合高频惯导融合 & 线性化不一致、雅可比错误 \\
UKF & 单高斯 & sigma 点传播 & 无需显式雅可比 & 高维点数增加、仍是单峰 \\
粒子滤波 & 带权样本 & 直接采样传播 & 可表示多峰 & 维数灾难、样本贫化 \\
因子图平滑 & 轨迹联合分布 & 非线性最小二乘 & 可重线性化、用回环修正历史 & 内存/求解开销、错误数据关联 \\
\bottomrule
\end{tabularx}
\end{table}

创新 $r_k$ 是最直接的在线诊断量。若模型正确，归一化创新平方 $r_k^\top S_k^{-1}r_k$ 应与观测维数对应的 $\chi^2$ 分布大致一致。长期偏大表示 $Q/R$ 太小、模型遗漏或观测异常；长期偏小可能表示协方差过于保守。门控只能拒绝离群值，不能修复持续偏置和错误外参。

\section{从递归滤波到因子图与 SLAM}

SLAM 同时估计机器人轨迹 $X=(x_0,\ldots,x_T)$ 与地图 $M$。里程计约束把相邻位姿连接，路标观测把位姿与地图点连接，回环约束把远隔时刻重新连接。定位依赖地图，地图又由带不确定性的位姿观测生成，因此二者不能作为独立黑盒串联\cite{durrantwhyte2006slam}。

在独立高斯测量下，最大后验估计可写成非线性最小二乘
\begin{equation}
X^*=\arg\min_X\sum_i\|r_i(X_i)\|_{\Sigma_i^{-1}}^2,
\label{eq:factor-graph-objective}
\end{equation}
每个因子只依赖少量变量，雅可比和正规方程因此稀疏。批量平滑会重新优化整个窗口；固定滞后平滑只保留最近状态；iSAM2 用 Bayes tree 对受新因子影响的部分增量重线性化和重求解\cite{kaess2012isam2}。相较只保留当前后验的滤波，平滑能让回环观测修正历史轨迹，但计算和内存代价更高。

视觉 SLAM 的典型管线包括特征/直接跟踪、局部地图、关键帧优化、回环检测与全局图优化。ORB-SLAM2 用统一系统覆盖单目、双目和 RGB-D 相机，并公开了关键帧式实现\cite{murartal2017orbslam2}。这类系统的公开精度不能直接迁移到另一相机、镜头、曝光和运动模式：低纹理、重复纹理、运动模糊、动态物体和时间不同步都会改变数据关联质量。

地图也不是单一对象。占据栅格回答哪里可通行，稠密几何回答表面形状，语义地图记录对象与区域，拓扑图记录地点连通。面向移动操作时，控制器通常需要带协方差和时间戳的局部位姿/速度，规划器需要碰撞几何，高层任务需要对象语义；强迫一个地图同时承担全部职责，会把更新频率、坐标系和不确定性混在一起。

\section{时间同步、外参与状态新鲜度}

设机器人以速度 $v$ 运动，传感器时间偏差为 $\Delta t$，则仅时间错位就产生约 $\|v\|\Delta t$ 的位置误差；旋转角速度 $\omega$ 产生约 $\|\omega\|\Delta t$ 的姿态误差。以 $v=1.5\ \mathrm{m/s}$、$\Delta t=80\ \mathrm{ms}$ 为例，位置错位为 $0.12\ \mathrm m$。若规划器的安全净空只有 $0.10\ \mathrm m$，再准确的几何地图也无法补偿这段时间误差。

外参把传感器坐标系变到机体或世界系。相机支架松动造成的 $1^\circ$ 姿态漂移，在 $5\ \mathrm m$ 距离上对应约 $8.7\ \mathrm{cm}$ 的横向误差。若只调大观测噪声，滤波器可能表现得更平滑，却不会消除系统偏置。工程上应把时间偏差、外参和 IMU 偏置视为可标定/可估计状态，并监控它们的可观测条件。

\begin{lstlisting}[caption={带时间戳、创新门控和新鲜度检查的融合教学伪代码}]
for measurement in time_ordered_queue:
    predict_to(measurement.stamp, imu_buffer)
    residual, S = innovation(measurement)
    if residual.T @ solve(S, residual) > chi2_gate:
        reject(measurement); continue
    update(measurement)
    enforce_symmetric_positive_covariance()

state, stamp = estimator.latest()
if now() - stamp > control_freshness_limit:
    controller.enter_safe_hold()
\end{lstlisting}

该伪代码刻意把“融合”和“向控制器发布”分开：估计器即使仍在运行，也可能因队列堵塞发布过期状态。除均值和协方差外，状态包至少应携带参考坐标系、测量时间、发布时间、外参版本、有效性标志和最近一次拒绝/重定位原因。

\section{移动操作案例：误差如何沿系统栈传播}

考虑移动底盘把机械臂送到货架前抓取。底盘定位使用轮速、IMU 和视觉，机械臂末端又依赖底盘到世界的变换。若视觉延迟导致底盘横向位姿偏差 $12\ \mathrm{cm}$，局部规划器可能仍在旧坐标系中认为通道有净空；机械臂到位后，末端相对货架的同量级偏差会触发第~\ref{chap:dynamics-contact-compliance}章的异常接触。此时提高阻抗控制器柔顺性只能降低冲击，不能恢复错误几何。

可验证的回退链为：首先用创新统计和传感器心跳区分离群观测与数据中断；其次检查状态时间戳和外参版本；若协方差或新鲜度越界，停止进入狭窄区并退到已知安全位；只有在局部重定位或重新标定后，才让第~\ref{chap:planning-optimal-control}章的规划器基于新地图重算。日志需保留原始时间戳、接受/拒绝原因、状态协方差、规划版本和最终动作，否则无法判断失败源于估计、规划还是控制。

\section{知识小结与综述判断}

知识小结：状态估计是带模型和不确定性的时间递推。KF 对线性高斯模型给出闭式预测/校正；EKF 用局部雅可比，UKF 用 sigma 点，粒子滤波用样本；因子图把整段轨迹写成稀疏最小二乘，使回环能修正历史。SLAM 的地图、位姿、数据关联和坐标系是一体化估计问题。

综述判断：具身系统最常见的估计风险不是“算法不够新”，而是时间、坐标系、偏置和协方差口径没有成为接口契约。后续的世界模型和长期记忆可以增加语义，但只要动作要落到真实身体，控制器仍需要一个带时间戳、参考系和可信度的物理状态；没有这些元数据，模型输出再丰富也无法形成可审计闭环。


\part{学习型机器人基础}
\chapter{模仿学习、DAgger 与偏好学习}
\label{chap:imitation-preference}

\section{问题设定：学习动作，还是学习目标}

示教数据通常写为 $\mathcal D_E=\{(o_t,a_t^E)\}$。行为克隆直接学习 $\pi_\theta(a\mid o)$；逆强化学习和偏好学习则先从轨迹比较中恢复奖励或效用，再由规划或强化学习求策略。二者回答不同问题：前者问“专家在这个观测下做了什么”，后者问“哪些结果更符合人的意图”。本章只讨论示教与偏好数据如何形成监督信号；MDP 与策略优化留给第~\ref{chap:reinforcement-learning}章，生成式动作分布留给第~\ref{chap:generative-action-policy}章。

\begin{figure}[H]
\centering
\begin{tabular}{c@{\ $\rightarrow$\ }c@{\ $\rightarrow$\ }c@{\ $\rightarrow$\ }c}
\fbox{专家/接管} & \fbox{轨迹与偏好} & \fbox{策略或奖励模型} & \fbox{影子评估与回流}
\end{tabular}
\caption{示教学习的闭环。作者教学绘制；关键不是一次性收集“干净数据”，而是明确哪些状态由当前策略访问、哪些失败需要重新标注。}
\label{fig:imitation-data-loop}
\end{figure}

\section{行为克隆与协变量偏移}

离散动作的行为克隆最小化负对数似然
\begin{equation}
\mathcal L_{\mathrm{BC}}(\theta)
=-\frac1N\sum_{i=1}^{N}\log\pi_\theta(a_i^E\mid o_i),
\label{eq:bc-nll}
\end{equation}
连续确定性动作常用 $\|a_i^E-\mu_\theta(o_i)\|_2^2$。这相当于假设给定观测只有一个条件均值；当“从杯子左侧绕过”和“从右侧绕过”都正确时，均方误差可能预测两条路线的平均值，恰好撞向杯子。第~\ref{chap:generative-action-policy}章将用动作块与生成模型处理这种多峰性。

更根本的问题是训练状态来自专家分布 $d_{\pi_E}$，部署状态来自学习策略自己的分布 $d_{\pi_\theta}$。一次小误差会把机器人带到示教中未出现的状态，后续预测因此更差。若把每一步独立失误率粗略写成 $\epsilon$，长度为 $T$ 的执行至少一次失误概率为 $1-(1-\epsilon)^T$。取 $\epsilon=0.02,T=50$，该概率约为 $63.6\%$；这不是严格的 DAgger 理论界，但直观说明逐帧准确率不能代替长时域成功率。

行为克隆仍是强基线：数据充足、闭环稳定、任务短且训练与部署分布接近时，它避免在线试错并易于调试。合格实验应同时报告轨迹级成功、恢复成功、接管次数、失败状态覆盖和动作延迟，而不只报告验证集 MSE。

\section{DAgger：在策略访问的状态上重新提问}

DAgger 迭代执行当前策略，收集其访问状态，并请求专家对这些状态给出动作，然后聚合到训练集\cite{ross2011dagger}。第 $i$ 轮可写为
\begin{equation}
\pi_i=\operatorname{Train}(\mathcal D_i),\qquad
\mathcal D_{i+1}=\mathcal D_i\cup
\{(o,a^E(o)):o\sim d_{\tilde\pi_i}\},
\end{equation}
其中 $\tilde\pi_i$ 可按概率混合专家与学习策略。它修正的不是模型容量，而是监督数据与部署状态分布的错位。

\begin{lstlisting}[caption={带安全接管和聚合预算的 DAgger 教学伪代码}]
dataset = expert_demonstrations
for round in range(max_rounds):
    policy = train_bc(dataset)
    for episode in shadow_or_guarded_rollouts(policy):
        for observation in episode:
            if supervisor.unsafe(observation):
                expert_action = supervisor.take_over(observation)
                dataset.add(observation, expert_action,
                            reason="safety_intervention")
    stop_if(validation_success_stable and intervention_rate_low)
\end{lstlisting}

真机 DAgger 的代价常被低估。专家必须在偏离状态下给出有意义动作；人机切换有延迟；“专家动作”可能依赖操作者看得到、策略看不到的信息；主动探索还可能损坏设备。因此工程实现常采用影子模式、仿真扰动、人工回放标注或只在安全域内聚合，而不是无条件让未成熟策略自由运行。

\section{逆强化学习与偏好学习}

逆强化学习假设专家行为近似优化某个未知奖励 $r_\phi(s,a)$。经典学徒学习通过匹配专家与策略的特征期望寻找奖励和策略\cite{abbeel2004apprenticeship}。奖励并不唯一：对所有轨迹加常数，甚至作势函数变换，都可能保持最优策略不变。因而恢复到的是与行为等价的一类目标，而不是专家脑中的“真实动机”。

偏好学习让标注者比较两段轨迹 $\tau^A,\tau^B$。若段得分为 $R_\phi(\tau)=\sum_t r_\phi(s_t,a_t)$，Bradley--Terry 模型给出
\begin{equation}
P_\phi(\tau^A\succ\tau^B)
=\frac{\exp R_\phi(\tau^A)}{\exp R_\phi(\tau^A)+\exp R_\phi(\tau^B)}.
\label{eq:preference-probability}
\end{equation}
对人类标签做交叉熵即可训练奖励模型，再用它优化策略。Christiano 等展示了从少量轨迹片段比较学习复杂目标的路径\cite{christiano2017preferences}。但奖励模型会被策略主动寻找漏洞；标注者间分歧、比较顺序、视频视角和未显示的安全信息都会进入标签噪声。

\begin{table}[H]
\centering
\caption{示教与偏好路线的监督对象和主要风险}
\label{tab:imitation-preference-comparison}
\begin{tabularx}{\textwidth}{p{0.15\textwidth}YYYY}
\toprule
方法 & 监督信号 & 是否需在线访问 & 优点 & 典型失效 \\
\midrule
行为克隆 & 状态--动作对 & 否 & 简单、稳定、可离线训练 & 协变量偏移、多峰平均 \\
DAgger & 当前策略状态上的专家动作 & 是/受控 & 直接覆盖恢复状态 & 标注昂贵、接管延迟 \\
逆强化学习 & 专家轨迹 & 常需反复求策略 & 可迁移目标表述 & 奖励不可辨识、计算昂贵 \\
偏好奖励 & 轨迹片段成对排序 & 可离线或主动查询 & 不要求逐帧动作 & 标注分歧、奖励黑客 \\
\bottomrule
\end{tabularx}
\end{table}

\section{示教策略上线：从数据到失败回流}

一个可审计流程包含五个门。第一，采集时保存原始观测、动作、操作者、设备标定、时间戳和任务结果；第二，按场景和操作者划分训练/验证集，避免相邻帧泄漏；第三，影子评估只记录策略建议，不驱动执行器；第四，在低速、限空间和人工接管下灰度；第五，把失败按感知、动作、时延、接触和任务定义分类后回流。只把成功轨迹追加到数据集，会系统性删除恢复行为。

上线指标应区分离线拟合与闭环结果。动作误差可发现单位或坐标问题，轨迹成功率衡量任务完成，接管率反映安全负担，失败恢复率衡量是否能离开分布外状态。任何指标都要按任务难度、物体和环境切片；总体平均数可能掩盖某个高风险子域完全失败。

\section{知识小结与综述判断}

知识小结：行为克隆在专家分布上做监督学习；DAgger 把监督查询移到当前策略访问的状态；逆强化学习和偏好学习把人类行为转成奖励或效用。四者的关键差异是标注对象、数据分布和是否允许在线访问，而不是网络骨干。

综述判断：机器人示教系统的上限常由数据协议而非损失函数决定。若观测不含专家决策所需信息、失败状态从不标注、接管延迟不记录，再复杂的策略也只能拟合一个不可部署的数据集。评价应把知识讲解与上线证据分开：先说明算法在什么假设下成立，再报告真实闭环是否满足这些假设。

\chapter{强化学习：从 MDP 到离线与安全 RL}
\label{chap:reinforcement-learning}

\section{MDP 是交互假设，不是算法名称}

马尔可夫决策过程写为 $\mathcal M=(\mathcal S,\mathcal A,P,r,\gamma)$。策略 $\pi(a\mid s)$ 产生动作，环境按 $P(s'\mid s,a)$ 转移并给奖励 $r$。目标是最大化
\begin{equation}
J(\pi)=\mathbb E_{\tau\sim\pi}\left[\sum_{t=0}^{T}\gamma^t r(s_t,a_t)\right].
\label{eq:rl-objective}
\end{equation}
“状态”必须足以预测未来；若相机图像缺少速度、接触历史或目标身份，实际问题是部分可观测的。第~\ref{chap:state-estimation-slam}章提供物理状态估计，第~\ref{chap:representation-transfer}章讨论学习表示，本章只研究给定状态/历史后怎样优化长期回报。

\begin{figure}[H]
\centering
\begin{tabular}{c@{\ $\rightarrow$\ }c@{\ $\rightarrow$\ }c@{\ $\rightarrow$\ }c}
\fbox{策略 $\pi$} & \fbox{动作 $a_t$} & \fbox{环境 $P$} & \fbox{状态与奖励 $s_{t+1},r_t$}
\end{tabular}
\caption{在线强化学习的数据闭环。作者教学绘制；离线 RL 切断了策略到环境的新交互箭头，只能使用固定日志。}
\label{fig:rl-interaction-loop}
\end{figure}

\section{价值函数、Bellman 方程与策略梯度}

策略价值和动作价值为
\begin{align}
V^\pi(s)&=\mathbb E_\pi\left[\sum_{k=0}^{\infty}\gamma^k r_{t+k}\mid s_t=s\right],\\
Q^\pi(s,a)&=\mathbb E_\pi\left[\sum_{k=0}^{\infty}\gamma^k r_{t+k}\mid s_t=s,a_t=a\right].
\end{align}
按第一步拆分得到 Bellman 期望方程
\begin{equation}
Q^\pi(s,a)=r(s,a)+\gamma\,
\mathbb E_{s'\sim P,a'\sim\pi}[Q^\pi(s',a')].
\label{eq:bellman-expectation}
\end{equation}
它把长时域目标转成一步奖励加后继价值，但也把模型、奖励和价值估计误差递归传播\cite{sutton2018reinforcement}。

以确定性两状态任务为例：在 $s_0$ 执行动作后即时奖励 1 并到 $s_1$，$s_1$ 的后续价值为 5，$\gamma=0.9$，则 $Q(s_0,a)=1+0.9\times5=5.5$，不是即时奖励 1。若价值 5 来自分布外动作的过估计，错误也被折扣后写回 $s_0$。

对可微随机策略，策略梯度定理给出
\begin{equation}
\nabla_\theta J(\pi_\theta)
=\mathbb E_{s,a\sim\pi_\theta}
\left[\nabla_\theta\log\pi_\theta(a\mid s)A^\pi(s,a)\right],
\label{eq:policy-gradient}
\end{equation}
$A=Q-V$ 衡量动作相对基线的增益。actor--critic 用 actor 表示策略、critic 估计 $V$ 或 $Q$。基线降低方差但不消除偏差；critic 错误会直接误导 actor。

最大熵方法在奖励之外鼓励策略熵。Soft Actor-Critic 使用离策略数据和熵正则，在连续控制中兼顾探索与样本复用\cite{haarnoja2018sac}。熵系数不是“随机程度旋钮”这么简单：动作维度、缩放和执行器饱和都会改变实际探索能量。

\section{在线、离线与示教增强：数据假设的差别}

在线 RL 可以用当前策略收集新数据并纠正价值估计，但真机样本昂贵且探索危险。离策略在线方法可复用旧数据，但仍假设能继续交互。离线 RL 完全使用固定数据集 $\mathcal D=\{(s,a,r,s')\}$；其核心风险是策略选择数据中罕见的动作，而 critic 在这些动作上没有可靠监督。

CQL 在 Bellman 误差之外加入保守项，使数据外动作的价值倾向被压低\cite{kumar2020cql}：
\begin{equation}
\mathcal L_{\mathrm{CQL}}(Q)=\mathcal L_{\mathrm{Bellman}}(Q)
+\alpha\left(\mathbb E_{s\sim\mathcal D,a\sim\mu}[Q(s,a)]
-\mathbb E_{(s,a)\sim\mathcal D}[Q(s,a)]\right).
\label{eq:cql-objective}
\end{equation}
$\mu$ 覆盖候选动作。$\alpha$ 太小不能抑制过估计，太大则把潜在优于数据行为的动作一并压低。离线 RL 不是“把 SAC 的 replay buffer 冻结”；必须显式处理数据支持和 OOD 动作。

示教增强可用专家数据初始化 replay buffer、给策略加行为克隆正则，或把示教与在线经验分层采样。它降低早期探索成本，却可能把专家偏差固化。比较方案时必须固定环境交互预算、示教数量和安全终止规则，否则“在线更强”可能只是用了更多真机尝试。

\begin{table}[H]
\centering
\caption{同一机器人任务中的训练数据假设}
\label{tab:rl-data-regimes}
\begin{tabularx}{\textwidth}{p{0.15\textwidth}YYYY}
\toprule
路线 & 新交互 & 主要优势 & 核心风险 & 合理证据 \\
\midrule
在线 RL & 持续 & 可主动纠正策略分布 & 样本成本与危险探索 & 真机步数、损坏/急停 \\
离策略在线 & 持续但可复用旧数据 & 样本效率较高 & replay 分布与当前策略错位 & 有效样本年龄、覆盖率 \\
离线 RL & 无 & 可用历史日志训练 & OOD 动作价值过估计 & 数据支持、保守度、离线选择 \\
示教增强 & 少量/分阶段 & 快速进入有效行为区 & 专家偏差、比较预算不公平 & 示教量、接管和新增交互 \\
\bottomrule
\end{tabularx}
\end{table}

\section{安全 RL：期望回报之外的约束}

安全问题可写成约束 MDP：最大化 $J_R(\pi)$，同时满足一个或多个代价预算
\begin{equation}
J_{C_j}(\pi)=\mathbb E_\pi\left[\sum_t\gamma^t c_j(s_t,a_t)\right]\le d_j.
\label{eq:cmdp}
\end{equation}
拉格朗日法优化 $J_R-\sum_j\lambda_j(J_{C_j}-d_j)$，但期望约束不能保证单次执行不越界；罕见碰撞可被大量安全回合平均掉。CPO 等方法在局部策略更新中近似控制约束变化\cite{achiam2017cpo}，仍依赖代价建模、信赖域近似和采样覆盖。

真实系统通常还需要独立安全层：动作限幅、碰撞监测、控制屏障函数、急停和人工接管。安全层会改变策略实际看到的转移；若日志只记录策略动作而不记录被过滤后的执行动作，训练数据将违反 MDP 语义。

\begin{lstlisting}[caption={带执行动作回写和安全终止的在线训练教学伪代码}]
proposed = policy.sample(state)
executed, intervention = safety_filter(proposed, state)
next_state, reward, cost = robot.step(executed)
replay.add(state, proposed, executed, reward, cost,
           next_state, intervention)
if cost > hard_limit or state_estimate_stale(next_state):
    robot.safe_stop(); mark_terminal(reason="safety")
update_actor_critic(replay, constraint_budget)
\end{lstlisting}

\section{机器人 RL 的实验与失败诊断}

奖励设计必须与传感可测量量一致。只奖励末端靠近目标可能诱导高速撞击；只惩罚能耗可能让机器人不动作；稀疏成功奖励则需要探索或示教。应把任务成功、约束代价、能耗、动作平滑、接管和硬件磨损分别报告，不用一个加权回报掩盖安全退化。

同一任务的最低比较应包含：行为克隆基线、固定数据离线 RL、等交互预算在线/示教增强方案；对每种方法报告数据构成、随机种子、策略选择依据和真机复验。不能用测试环境回报反复挑 checkpoint，再宣称零样本泛化。

\section{知识小结与综述判断}

知识小结：Bellman 方程递归连接即时奖励和未来价值；策略梯度通过优势函数更新动作概率；actor--critic 把策略与价值估计耦合。在线、离策略在线和离线 RL 的根本差别是能否从当前策略获得新数据；安全 RL 又在期望回报之外加入代价约束。

综述判断：机器人 RL 的首要问题不是选择算法缩写，而是声明交互权限、数据支持和安全边界。离线数据不足时，价值函数的精确小数没有物理意义；安全代价不可测时，约束优化也只是形式。可信结论必须把算法性能与真机交互成本、介入机制和失败后果放在同一张证据表中。

\chapter{表示学习、多任务学习与技能迁移}
\label{chap:representation-transfer}

\section{表示的职责与本章边界}

表示 $z=f_\theta(o)$ 的价值不在于“维度更小”，而在于保留下游决策需要的信息，同时抑制光照、背景等无关变化。多任务学习进一步让任务条件 $c$ 与共享表示共同决定动作或价值。本文后续第 14--17 章将讲多模态编码器和基础模型接口；本章只研究表示怎样通过自监督、任务共享和迁移被训练，以及何时共享反而有害。

\begin{figure}[H]
\centering
\begin{tabular}{c@{\ $\rightarrow$\ }c@{\ $\rightarrow$\ }c}
\fbox{多模态观测 $o$} & \fbox{共享表示 $z=f_\theta(o)$} &
\begin{tabular}{c}\fbox{任务头 $g_1(z,c_1)$}\\[4pt]\fbox{任务头 $g_2(z,c_2)$}\\[4pt]\fbox{技能/策略头}\end{tabular}
\end{tabular}
\caption{共享表示与任务专用接口。作者教学绘制；共享编码器减少重复学习，但不保证各任务梯度相容。}
\label{fig:shared-representation-heads}
\end{figure}

\section{自监督表示：什么被视为相同}

对比学习用正样本对 $(o,o^+)$ 和负样本集合学习相似性。InfoNCE 的一类写法为
\begin{equation}
\mathcal L_{\mathrm{NCE}}
=-\log\frac{\exp(\operatorname{sim}(z,z^+)/\tau)}
{\exp(\operatorname{sim}(z,z^+)/\tau)+\sum_j\exp(\operatorname{sim}(z,z_j^-)/\tau)}.
\label{eq:infonce}
\end{equation}
温度 $\tau$ 控制分布尖锐度。关键选择不是公式，而是正样本语义：同一图像的颜色增强鼓励光照不变性，相邻视频帧鼓励短时一致性，同一任务不同视角鼓励视角不变性。若把接触前后两帧误当正样本，表示可能主动丢掉任务关键的接触变化。

重建式目标要求从 $z$ 恢复输入或遮蔽部分，易保留细节，但也可能把容量花在背景像素。时序预测要求 $z_t$ 预测未来 $z_{t+k}$，可学习动力相关信息，却可能忽略对短期预测无帮助但对任务重要的语义。不存在脱离下游任务的“通用最好表示”。

R3M 使用大规模人类第一视角视频，并结合时间对比、视频--语言对齐和稀疏正则学习机器人操作视觉表示；论文在多项模拟和真实操作任务上评估冻结特征的数据效率\cite{nair2022r3m}。这说明人类视频可提供可迁移视觉先验，但不证明其表示天然包含机器人本体状态、力觉或精确动作坐标。

\section{多任务目标、任务条件与负迁移}

设任务 $i$ 的损失为 $\mathcal L_i$，共享参数目标通常写为
\begin{equation}
\mathcal L(\theta,\phi_{1:m})=\sum_{i=1}^{m}w_i
\mathbb E_{(o,a)\sim\mathcal D_i}
[\ell_i(g_{\phi_i}(f_\theta(o),c_i),a)].
\label{eq:multitask-loss}
\end{equation}
$w_i$ 同时受到任务重要性、数据量、损失量纲和采样频率影响。若简单按样本均匀采样，大数据任务会支配共享编码器；若按任务均匀采样，小任务又可能被过采样。

共享是否冲突可先看梯度余弦
\begin{equation}
\rho_{ij}=\frac{g_i^\top g_j}{\|g_i\|_2\|g_j\|_2},
\qquad g_i=\nabla_\theta\mathcal L_i.
\label{eq:gradient-cosine}
\end{equation}
$\rho_{ij}<0$ 表示沿任务 $i$ 的下降方向可能增大任务 $j$ 的一阶损失。以 $g_1=(1,0)$、$g_2=(-0.8,0.2)$ 为例，点积为 $-0.8$；直接求和得到 $(0.2,0.2)$，显著削弱任务 1 的更新。这只是局部诊断，不代表两个任务永远不相容，但比笼统声称“共享促进泛化”更可检验。

MT-Opt 在大规模真实机器人多任务数据上使用任务条件化和数据共享，展示了相关任务数据可帮助长尾任务\cite{kalashnikov2021mtopt}。其结论依赖任务族、数据规模和统一机器人接口；把不同本体、不同动作单位的数据直接混合，并不会自动得到跨本体技能。

\section{技能表示与跨本体迁移}

技能可表示为离散 ID、语言条件、连续潜变量 $z_s$ 或一段闭环策略。连续技能潜变量常通过编码器从轨迹推断，再由低层策略执行 $\pi(a\mid s,z_s)$。若潜变量只优化重建动作，它可能编码操作者身份或速度，而不是可组合的任务语义；需要任务成功、可预测性或互信息约束来校准含义。

跨本体迁移至少要对齐三类接口：观测语义、动作语义和可达能力。把末端增量作为中间接口可弱化关节数差异，但不同机器人工作空间、负载、夹爪和控制频率仍不同。共享策略应显式接收本体描述或使用本体专用适配器，并在第~\ref{chap:kinematics-lie}章定义的任务空间与关节空间之间做受约束映射。

\begin{table}[H]
\centering
\caption{表示和迁移时常见参数复用方式}
\label{tab:representation-transfer-strategies}
\begin{tabularx}{\textwidth}{>{\raggedright\arraybackslash}p{0.17\textwidth}YYYY}
\toprule
方式 & 更新范围 & 数据需求 & 优点 & 主要风险 \\
\midrule
冻结编码器 & 只训任务头 & 少 & 稳定、便于比较 & 表示缺少动作/本体信息 \\
全量微调 & 全部参数 & 多 & 适应能力强 & 灾难遗忘、过拟合 \\
适配器/低秩更新 & 小量任务参数 & 中 & 保留共享骨干 & 容量不足或任务隔离 \\
共享骨干多头 & 编码器与各头 & 多任务同步 & 数据复用 & 梯度冲突、任务不平衡 \\
本体条件化 & 共享策略加本体描述 & 跨本体配对/覆盖 & 接口统一 & 未见本体外推不可靠 \\
\bottomrule
\end{tabularx}
\end{table}

\section{表示坍缩与可审计诊断}

表示坍缩指不同输入映射到近似常数。即使没有完全坍缩，也可能出现任务条件被忽略、少数维度支配或背景捷径。诊断应包含特征方差/秩、任务线性探针、跨场景最近邻、遮挡/背景干预和下游闭环性能。只展示 t-SNE 图不能证明可控性或因果相关性。

\begin{lstlisting}[caption={多任务共享训练中的梯度冲突监控教学伪代码}]
task_batches = sampler.balanced_batches()
task_grads = []
for task, batch in task_batches:
    loss = task_loss(task, batch)
    task_grads.append(grad(loss, shared_encoder))
conflict = pairwise_cosine(task_grads)
log(conflict, feature_rank(shared_encoder), per_task_metrics)
if persistent_negative_transfer(conflict, per_task_metrics):
    split_adapter_or_reweight_tasks()
apply_joint_update(task_grads)
\end{lstlisting}

冻结、微调和多任务训练必须用相同数据预算、增强和任务头容量比较。否则骨干差异会与训练资源混在一起。表示的最终证据不是预训练损失，而是在未见场景和受控干预下，任务成功是否更稳、需要的标注是否更少。

\section{知识小结与综述判断}

知识小结：对比、重建和时序预测分别规定了不同的不变性；多任务目标通过共享参数复用数据，也可能因梯度冲突和任务不平衡产生负迁移；技能迁移还要求观测、动作和本体能力对齐。

综述判断：所谓“通用表示”必须附带任务集合和接口条件。一个在图像语义上强的编码器未必适合毫米级接触控制；一个在同构机械臂间共享的策略也未必能外推到新夹爪。合理做法是把共享范围当作实验变量，逐层证明哪些知识可复用、哪些必须由任务或本体适配器承担。

\chapter{生成式动作策略：动作块、扩散与流匹配}
\label{chap:generative-action-policy}

\section{为什么动作条件分布需要生成模型}

给定观测历史 $o_{t-h:t}$ 和任务条件 $c$，策略学习的是未来动作分布 $p(a_{t:t+H-1}\mid o,c)$。单步均方回归把多峰分布压成条件均值；自回归离散 token、潜变量动作块、扩散和流匹配则用不同方式保留多峰性。本章主讲动作头和滚动执行，不重复第~\ref{chap:imitation-preference}章的数据采集，也不展开第 23 章的 VLA 骨干比较。

\begin{figure}[H]
\centering
\setlength{\tabcolsep}{3pt}
\begin{tabular}{cccc}
\fbox{\begin{minipage}{0.19\textwidth}\centering 单步连续回归\\$a_t=\mu(o_t)$\end{minipage}} &
\fbox{\begin{minipage}{0.19\textwidth}\centering 动作块/潜变量\\$a_{t:t+H-1}$\end{minipage}} &
\fbox{\begin{minipage}{0.19\textwidth}\centering 扩散去噪\\$x_K\rightarrow\cdots\rightarrow x_0$\end{minipage}} &
\fbox{\begin{minipage}{0.19\textwidth}\centering 流 ODE\\$\dot x=v_\theta(x,t,o)$\end{minipage}}
\end{tabular}
\caption{四类动作头的计算对象。作者教学绘制；VLA 离散 token 路线在第 23 章详细比较。}
\label{fig:generative-action-heads}
\end{figure}

\section{动作块：预测窗口与执行窗口不是一回事}

动作块一次预测 $H$ 个动作。其监督目标可写成
\begin{equation}
\mathcal L_{\mathrm{chunk}}=-\mathbb E_{(o,a_{0:H-1})\sim\mathcal D}
\log p_\theta(a_{0:H-1}\mid o,c).
\end{equation}
块内相关性让模型学习“接近--对准--插入”这样的短时结构，并减少逐步推理开销。ACT 以 Transformer 和条件变分模型学习双臂动作块，用时间集成缓解相邻块预测不一致\cite{zhao2023act}。

设预测块长 $H=16$，系统只执行前 $E=8$ 步后重新观测。$H$ 决定模型看多远，$E$ 决定多久闭环一次。$E=H$ 时延低但开环长；$E=1$ 最闭环但计算贵，且相邻预测可能抖动。时间集成可对覆盖同一时刻的多个历史预测加权平均，但在接触模式突然变化时，旧预测会拖慢响应。

潜变量模型引入 $z$ 表示不同动作模式，证据下界为
\begin{equation}
\mathcal L_{\mathrm{CVAE}}
=\mathbb E_{q_\phi(z\mid a,o)}[-\log p_\theta(a\mid z,o)]
+\beta D_{\mathrm{KL}}(q_\phi(z\mid a,o)\|p(z)).
\end{equation}
$\beta$ 太大可能使后验忽略动作而坍缩，太小则训练后验与推理先验不匹配。

\section{扩散策略：从动作噪声反复去噪}

令干净动作块为 $x_0$。前向过程逐步加噪，并可直接采样
\begin{equation}
x_k=\sqrt{\bar\alpha_k}x_0+\sqrt{1-\bar\alpha_k}\,\epsilon,
\qquad \epsilon\sim\mathcal N(0,I).
\label{eq:diffusion-forward}
\end{equation}
网络以观测条件 $c_o$、噪声层 $k$ 和 $x_k$ 预测噪声：
\begin{equation}
\mathcal L_{\mathrm{diff}}=
\mathbb E_{x_0,k,\epsilon}
\|\epsilon-\epsilon_\theta(x_k,k,c_o)\|_2^2.
\label{eq:diffusion-noise-loss}
\end{equation}
推理从高斯噪声开始，按调度器迭代得到动作块。Diffusion Policy 将视觉条件、时序扩散网络与滚动时域执行结合，用动作序列分布表达多峰操作行为\cite{chi2023diffusionpolicy}。

扩散模型的优势不是“动作更平滑”这一口号，而是能在连续高维块上表示复杂分布；代价是多次网络调用、调度器敏感性和闭环延迟。动作归一化尤为关键：平移米、旋转弧度和夹爪开度若共用尺度，噪声日程会对各维度产生完全不同的物理扰动。

\begin{lstlisting}[caption={滚动执行的条件动作扩散教学伪代码}]
while task_active:
    condition = encode(observation_history, task)
    actions = normal_noise(shape=(horizon, action_dim))
    for k in reversed(noise_schedule):
        eps = noise_net(actions, k, condition)
        actions = scheduler.step(eps, k, actions)
    actions = denormalize_and_project_limits(actions)
    execute(actions[:execution_horizon])
    if contact_mode_changed() or state_stale(): replan_immediately()
\end{lstlisting}

投影动作限位只保证单步边界，不能保证整段轨迹无碰撞或动态可行；必要时仍要调用第~\ref{chap:planning-optimal-control}章的约束层。若去噪耗时 120 ms 而控制周期 20 ms，生成结果到达时状态已经变化，离线成功率不能解释在线失败。

\section{Flow Matching：直接学习概率路径的速度场}

Flow Matching 选择从简单分布 $p_0$ 到数据分布 $p_1$ 的条件概率路径 $x_t$，训练向量场匹配目标速度 $u_t$：
\begin{equation}
\mathcal L_{\mathrm{FM}}(\theta)
=\mathbb E_{t,x_t}\|v_\theta(x_t,t,c_o)-u_t(x_t\mid x_1)\|_2^2.
\label{eq:flow-matching}
\end{equation}
推理求解常微分方程 $\mathrm d x/\mathrm dt=v_\theta(x,t,c_o)$ 从噪声到动作。Flow Matching 给出了无需模拟随机扩散路径的向量场回归框架\cite{lipman2023flowmatching}；$\pi_0$ 等机器人策略把流式动作头用于连续动作生成\cite{black2024pi0}。

“ODE 所以一步即可”并不成立。数值积分仍可能需要多个函数评估；直线路径、网络误差和求解器容差决定步数与质量。相比扩散，流模型有机会减少采样步数，但训练目标、路径选择和蒸馏策略不同，不能只比较名义迭代次数。

\section{动作表示与执行接口比较}

\begin{table}[H]
\centering
\caption{动作头的分布能力、计算与闭环权衡}
\label{tab:generative-action-comparison}
\begin{tabularx}{\textwidth}{>{\raggedright\arraybackslash}p{0.18\textwidth}YYYY}
\toprule
路线 & 输出对象 & 多峰表达 & 推理特点 & 主要失效 \\
\midrule
连续单步回归 & 一个连续动作 & 弱 & 一次前向、低延迟 & 均值动作、短视 \\
离散自回归 token & token 序列 & 中/强 & 逐 token 解码 & 量化误差、时延随长度增 \\
潜变量动作块 & 连续动作块 & 中 & 一次或少量采样 & 块间不一致、潜变量坍缩 \\
扩散动作块 & 连续动作块 & 强 & 多步去噪 & 延迟、调度与尺度敏感 \\
流式动作块 & 连续动作块 & 强 & ODE/少步积分潜力 & 路径与积分误差 \\
\bottomrule
\end{tabularx}
\end{table}

比较必须固定观测、骨干、数据、动作频率和执行窗口。只让扩散模型使用动作块而让回归基线单步输出，会把时序结构收益错误归因于生成范式。至少应做单步/动作块消融、推理函数评估次数、端到端延迟、任务成功、恢复能力和动作边界违反率。

\section{官方实现导读与复现锚点}

Diffusion Policy 官方仓库把数据集归一化、workspace 训练入口、policy 条件去噪和环境 runner 分开；正文初稿以其 RSS 论文和官方仓库作为实现入口，冻结前需固定 immutable commit、许可证、策略类和推理函数行段。ACT 的公开实现同样把示教数据加载、CVAE/Transformer 策略和 rollout 分开。代码阅读应沿“数据字段--归一化--损失--采样--执行窗口”追踪，而不是只截取网络定义。

对任何生成动作头，复现记录至少包含：动作坐标系与单位、归一化统计量来源、预测/执行窗口、噪声或流时间离散、推理调用次数、控制周期、动作后处理和随机种子。缺少这些字段，模型权重无法解释为可执行策略。

\section{知识小结与综述判断}

知识小结：动作块用联合序列减少逐步短视；扩散通过前向加噪和反向去噪学习连续多峰动作；Flow Matching 回归概率路径速度场并通过 ODE 生成动作。预测窗口、执行窗口和控制周期共同决定闭环性质。

综述判断：生成模型解决的是条件动作分布和时序一致性，不自动解决物理可行性。论文中的离线动作拟合、仿真 rollout 与真机闭环是三种证据；只有同时报告推理延迟、动作尺度、约束违反和失败恢复，才能判断扩散或流式动作头是否真的优于简单动作块。

\chapter{sim2real、域适应与训练闭环}
\label{chap:sim2real-loop}

\section{现实差距不是一个标量}

sim2real 的目标不是让仿真“看起来像现实”，而是让在仿真分布上训练的系统在真实任务分布上保持可接受风险。差距至少包括观测渲染、传感噪声、动力学参数、接触、执行器、时延、任务对象和人类行为。不同差距需要不同工具：系统辨识缩小参数偏差，域随机化扩大训练覆盖，域适应对齐表示，残差学习补偿稳定模型误差，在线适应根据历史推断当前环境。

本章只连接第~\ref{chap:imitation-preference}--\ref{chap:generative-action-policy}章形成训练与上线闭环，不重新解释其中算法。

\begin{figure}[H]
\centering
\setlength{\tabcolsep}{3pt}
\begin{tabular}{c@{\ $\rightarrow$\ }c@{\ $\rightarrow$\ }c@{\ $\rightarrow$\ }c}
\fbox{\begin{minipage}{0.16\textwidth}\centering 仿真/历史\\数据\end{minipage}} &
\fbox{\begin{minipage}{0.16\textwidth}\centering 训练与\\压力测试\end{minipage}} &
\fbox{\begin{minipage}{0.16\textwidth}\centering 影子/小流量\\真机\end{minipage}} &
\fbox{\begin{minipage}{0.16\textwidth}\centering 失败分类\\与回采\end{minipage}}
\end{tabular}
\caption{从仿真训练到真机回采的版本闭环。作者教学绘制；每个箭头都必须保存数据、配置、模型和退出原因。}
\label{fig:sim2real-release-loop}
\end{figure}

\section{系统辨识、域随机化与鲁棒目标}

令仿真域参数为 $\xi$，包括质量、摩擦、时延、纹理等。域随机化训练
\begin{equation}
\min_\theta\;\mathbb E_{\xi\sim p_{\mathrm{train}}(\xi),\tau\sim\pi_\theta,P_\xi}
[\mathcal L(\tau)].
\label{eq:domain-randomization-risk}
\end{equation}
Tobin 等通过随机纹理、光照和相机参数研究从合成视觉到真实检测/抓取的迁移\cite{tobin2017domainrandomization}；动力学随机化则在质量、阻尼、摩擦等参数上训练控制策略\cite{peng2018dynamicsrandomization}。

$p_{\mathrm{train}}$ 太窄，真实参数落在支持集外；太宽，策略可能为极端域牺牲正常性能，甚至找不到稳定行为。随机化范围应来自标定数据、制造公差和故障假设，并进行逐因素消融，而不是凭经验把所有参数乘上 $[0.5,1.5]$。

系统辨识用真实输入输出估计 $\hat\xi$，再围绕它随机化。残差动力学写成
\begin{equation}
x_{t+1}=f_{\mathrm{sim}}(x_t,u_t;\hat\xi)+r_\psi(x_t,u_t),
\label{eq:residual-dynamics}
\end{equation}
或在控制器输出上加受限残差。残差模型适合补偿可重复偏差，但数据外外推可能破坏稳定性；应限制幅值、频带和适用域，并保留解析基线。

鲁棒训练还可关注最坏域或尾部风险，例如
\begin{equation}
\min_\theta\max_{q\in\mathcal Q}
\mathbb E_{\xi\sim q}[\mathcal L(\pi_\theta;\xi)].
\end{equation}
它提高对指定不确定集的保护，不代表对未建模故障普遍安全。

\section{域适应与在线适应}

视觉域适应试图让仿真与真实观测的特征分布接近，可用风格变换、对抗对齐、自训练或少量真实标签微调。边缘分布对齐不保证任务条件分布对齐：若仿真中红色总对应目标、现实中颜色无关，特征看似对齐仍会保留错误捷径。

在线适应从最近观测/动作历史推断环境潜变量 $z_t$，策略执行 $a_t=\pi(o_t,z_t)$。RMA 把基础策略与适应模块分开，在模拟随机域训练后用历史信号快速推断真实环境变化\cite{kumar2021rma}。这类方法要求适应信号在任务运动中可观；静止时无法辨识地面摩擦，短历史也可能混淆负载与执行器故障。

\begin{table}[H]
\centering
\caption{现实差距、可观测证据与主要补救手段}
\label{tab:sim2real-gap-matrix}
\begin{tabularx}{\textwidth}{p{0.16\textwidth}YYYp{0.23\textwidth}}
\toprule
差距 & 可观测证据 & 主要手段 & 验证切片 & 常见误用 \\
\midrule
视觉/传感 & 像素、噪声谱、缺帧 & 渲染随机化、适应 & 光照、相机、遮挡 & 只追求逼真截图 \\
刚体动力学 & 轨迹、力矩、加速度 & 辨识、参数随机化 & 负载、速度、姿态 & 参数范围无数据依据 \\
接触/摩擦 & 力、滑移、插入结果 & 接触辨识、残差、真机数据 & 材料、磨损、公差 & 用视觉随机化替代接触建模 \\
时延/控制栈 & 时间戳、队列、执行回执 & 延迟注入、预测、限速 & CPU 负载、网络抖动 & 只随机固定延迟 \\
任务/人员 & 对象与操作序列 & 场景生成、真实回采 & 长尾物体、异常流程 & 仿真任务过于整洁 \\
\bottomrule
\end{tabularx}
\end{table}

\section{时延与接触：最容易被低估的差距}

若机器人以 $1\ \mathrm{m/s}$ 运动，总感知--推理--通信时延为 $80\ \mathrm{ms}$，状态错位已达 $8\ \mathrm{cm}$。第~\ref{chap:state-estimation-slam}章说明时间戳错误如何污染估计；这里的训练要求是在仿真中注入与真实分布一致的抖动、丢包和计算超时，而不是只加一个固定延迟常数。

接触差距更难由参数随机化覆盖。摩擦、柔顺、间隙和碰撞求解器的离散步长耦合。插装策略在仿真中成功，可能只是利用了允许穿透或过度光滑的接触模型。应把接触峰值、滑移、插入力、恢复次数和表面损伤纳入真机门，而不只看最终是否插入。

\section{影子验证、灰度放行与失败回流}

训练完成不等于可部署。建议的放行序列为：离线重放检查单位和接口；仿真压力测试覆盖参数尾部；hardware-in-the-loop 验证控制周期与协议；影子模式比较建议动作但不执行；受限工作空间低速小流量；逐步扩大任务分布。每一级必须定义进入条件、退出条件和自动回退版本。

\begin{lstlisting}[caption={模型版本放行与失败回采的教学伪代码}]
candidate = registry.load(version)
assert offline_replay(candidate).passes(interface_and_regression_gates)
assert sim_stress(candidate).passes(tail_risk_gates)
shadow = run_shadow(candidate, production_inputs)
if shadow.disagreement_or_latency > limits: reject(candidate)
for scope in [low_speed_cell, guarded_tasks, limited_fleet]:
    result = deploy(candidate, scope, rollback=stable_version)
    if result.safety_event or result.metric_regression:
        rollback(); collect_failure_bundle(result); break
promote_only_if(all_gates_passed_and_reviewed)
\end{lstlisting}

失败包应包含原始传感、状态估计、模型输入输出、实际执行动作、安全过滤、控制回执、仿真参数、软件/模型版本和操作者说明。回采后先分类根因，再决定补真实示教、扩大随机化、更新辨识或修改安全层。把所有失败都追加为“困难样本”可能让模型学习互相矛盾的标签。

\section{贯通第 9--12 章的闭环案例}

以箱体抓取为例：第~\ref{chap:imitation-preference}章的示教建立基本接近与抓取数据；第~\ref{chap:reinforcement-learning}章可在仿真中优化长时域成功和能耗；第~\ref{chap:representation-transfer}章共享不同物体的视觉表示；第~\ref{chap:generative-action-policy}章生成多峰动作块。sim2real 阶段不再重新发明这些算法，而是逐项核验观测分布、动作单位、动力学、接触、时延与安全接口。

若真机在透明物体上失败，先检查感知与数据覆盖；若到达但夹取滑落，检查接触与夹爪模型；若离线动作正确但在线抖动，检查推理/控制时延和执行窗口；若特定负载失败，检查参数支持与适应可观测性。不同根因对应不同回流，避免用无方向的“继续训练”掩盖工程问题。

\section{知识小结与综述判断}

知识小结：系统辨识缩小模型偏差，域随机化扩大参数覆盖，域适应对齐观测表示，残差与在线适应补偿剩余变化。现实差距必须按观测、动力学、接触、时延和任务分别建模，并通过影子、灰度和回滚形成版本闭环。

综述判断：sim2real 不是某个训练技巧，而是一套证据链。随机化只有在范围有测量依据、真实域位于支持集且策略在尾部仍可控时才成立；“零样本上真机”也不意味着零真实信息，因为标定、任务设计和参数范围本身都来自现实。可信报告应公开这些现实信息怎样进入训练与放行。


\part{多模态与基础模型基础}
\chapter{Transformer 与动作序列建模}
\label{chap:transformer-action-sequence}

\section{注意力计算：内容寻址而非时间模型本身}

给定序列表示 $X\in\mathbb R^{L\times d}$，线性投影得到 $Q=XW_Q,K=XW_K,V=XW_V$，scaled dot-product attention 为
\begin{equation}
\operatorname{Attn}(Q,K,V)=
\operatorname{softmax}\left(\frac{QK^\top}{\sqrt{d_k}}+M\right)V.
\label{eq:scaled-attention}
\end{equation}
$M$ 可屏蔽未来 token 或无效 padding。除以 $\sqrt{d_k}$ 是为了避免维度增大时点积方差过大，使 softmax 过早饱和。多头注意力让不同子空间独立计算，再拼接投影。Transformer 用自注意力与前馈网络取代循环结构，但顺序必须由位置编码、因果掩码或时间戳显式提供\cite{vaswani2017attention}。

注意力权重不是因果解释。一个 token 权重高只说明当前计算中的相似性和归一化结果；层间残差、MLP 和后续头还会改变贡献。机器人报告不应把热力图直接写成“模型关注了正确物体”的行为证据。

\section{机器人序列如何组织}

机器人输入可能包括语言 $l$、图像 patch $v_t$、本体状态 $p_t$、触觉 $h_t$ 和历史动作 $a_{t-1}$。一种交错序列为
\begin{equation}
[l_{1:m},v_0,p_0,h_0,a_0,v_1,p_1,h_1,a_1,\ldots].
\end{equation}
另一种做法先在每个时刻融合多模态，再把时刻 latent 组成时间序列。前者保留细粒度交互但 token 多，后者计算便宜但连接器成为信息瓶颈。

\begin{figure}[H]
\centering
\begin{tabular}{c@{\quad}c@{\quad}c@{\quad}c@{\quad}c}
\fbox{语言前缀} & \fbox{$v_{t-1},p_{t-1}$} & \fbox{$a_{t-1}$} & \fbox{$v_t,p_t$} & \fbox{预测 $a_t$}
\end{tabular}
\caption{机器人因果序列的最小时间轴。作者教学绘制；每个观测块还必须携带采样时间，排列顺序不能修复时间戳错误。}
\label{fig:robot-causal-token-timeline}
\end{figure}

位置编码只区分槽位，不等于物理时间。相机 30 Hz、本体 200 Hz、触觉 1 kHz 时，按到达顺序拼接会把通信抖动误写成动作因果。应先按测量时间对齐或保留连续时间嵌入，并标注缺帧、插值和状态新鲜度。

\section{自回归目标与动作序列的特殊性}

因果模型最小化
\begin{equation}
\mathcal L_{\mathrm{AR}}=-\sum_{i=1}^{L}
\log p_\theta(y_i\mid y_{<i},c).
\label{eq:autoregressive-loss}
\end{equation}
文本 token 错误通常仍可在后续文本中修正；动作 token 会改变世界，后续条件分布随之变化。动作还具有连续单位、速度/加速度约束和不可逆风险。因而相同的交叉熵并不意味着相同的闭环代价。

连续动作可直接回归，也可量化成离散 token。RT-1 把机器人动作离散化并与图像/语言条件一起建模\cite{brohan2022rt1}。离散化便于复用语言模型头，但每维 bin 宽决定物理分辨率；独立量化各维又忽略关节/末端耦合。第 12 章讨论扩散和流式连续动作，第 23 章比较代表性 VLA 的动作接口。

teacher forcing 训练时使用真实前缀，部署时使用模型先前输出，形成 exposure bias。动作块减少自回归步数，却延长开环执行窗口。评估必须报告 rollout 成功和恢复，而非只给 token 准确率。

\section{复杂度、KV cache 与实时性}

标准全注意力计算和注意力矩阵内存约为 $O(L^2d)$ 与 $O(L^2)$。自回归推理若缓存各层历史 key/value，新 token 不必重算全部前缀，但缓存内存约随 $O(N_{\mathrm{layer}}Ld)$ 增长。多相机 patch、长历史和大模型层数会快速耗尽显存。

设控制频率 50 Hz，即周期 20 ms；端到端推理 100 ms 意味着动作使用至少 5 个周期前的状态。即使平均延迟 18 ms，99.9 分位 80 ms 也可能触发接触失败。应分别测量编码、连接器、解码、后处理和通信时间，并报告尾延迟。

\begin{lstlisting}[caption={带时间戳与 KV cache 失效条件的因果推理教学伪代码}]
tokens, stamps = encode_new_observations(sensor_packet)
if stale(stamps) or calibration_version_changed():
    kv_cache.clear(); controller.safe_hold()
tokens = append_time_embeddings(tokens, stamps)
logits, kv_cache = model.decode(tokens, kv_cache=kv_cache)
action = project_units_limits(decode_action(logits))
if inference_deadline_missed(): use_backup_controller()
else: execute(action)
\end{lstlisting}

缓存不是永远有效的。机器人重定位、任务切换、外参更新或安全回滚后，旧 KV 中的状态语义可能与当前世界冲突。需要缓存版本、任务 ID 和物理状态重置条件。

\section{文本序列与动作序列的证据差异}

\begin{table}[H]
\centering
\caption{文本与机器人动作序列的关键差异}
\label{tab:text-action-sequence-differences}
\begin{tabularx}{\textwidth}{p{0.18\textwidth}YY}
\toprule
维度 & 文本生成 & 动作生成 \\
\midrule
token 语义 & 离散符号，错误多可继续阅读 & 连续物理量或量化量，改变环境状态 \\
时间 & 序列位置常足够 & 测量时间、控制周期和延迟直接影响可执行性 \\
约束 & 语法/事实约束为主 & 关节、碰撞、动力学、接触与安全约束 \\
反馈 & 输出后通常无物理反馈 & 每步观测闭环，状态分布由动作改变 \\
评测 & 困惑度、准确率、偏好 & 任务成功、约束违反、延迟、恢复与损坏 \\
\bottomrule
\end{tabularx}
\end{table}

\section{知识小结与综述判断}

知识小结：注意力做内容寻址，位置/时间编码与掩码赋予顺序；自回归目标预测下一个 token；KV cache 用内存换取增量推理。机器人序列还必须组织多率传感、动作单位和状态时间戳。

综述判断：把动作写成 token 只统一了软件接口，没有消除物理约束。Transformer 的主要价值是灵活融合长序列条件；其主要工程代价是 token 预算、尾延迟和缓存一致性。任何“实时”结论都应给出完整控制周期和尾分位，而不是只报模型前向平均值。

\chapter[多模态编码器、连接器与接口]{多模态编码器、连接器\\与接口}
\label{chap:multimodal-connectors}

\section{不同模态解决不同可观测性问题}

RGB 提供纹理与语义，深度/点云提供几何，触觉提供局部接触形变与滑移，本体状态提供关节和执行器信息，音频可提示碰撞或电机异常。多模态并非“信号越多越好”：每种传感器都有坐标系、采样率、延迟和失效模式。编码器的任务是把原始信号变成下游可用特征，连接器则把异构特征映射到共享模型接受的 token 或 latent 接口。

\section{视觉、三维、触觉与本体编码}

ViT 把图像划分为 $P\times P$ patch。对 $H\times W$ 图像，token 数为
\begin{equation}
N_v=\frac{H}{P}\frac{W}{P}.
\end{equation}
若 $H=W=224,P=16$，单帧有 196 个 patch token；两相机、8 帧即 3136 个视觉 token，尚未计入语言和动作。ViT 说明了 patch 序列进入 Transformer 的基本接口\cite{dosovitskiy2021vit}；CLIP 用图文对比预训练获得开放词汇语义对齐\cite{radford2021clip}，但其全局语义不保证毫米级几何和接触状态。

点云编码需处理排列不变性、密度和坐标系；体素/稀疏卷积适合规则空间索引，点式网络保留不规则采样。触觉图像看似 RGB，物理含义却来自弹性体形变、照明和接触标定。GelSight 类传感器从高分辨率表面形变估计几何和力\cite{yuan2017gelsight}；直接套用自然图像增强可能删除微小滑移纹理。本体状态通常经 MLP/傅里叶特征编码，但单位、关节顺序和归一化版本必须写入接口。

\begin{table}[H]
\centering
\caption{模态、编码目标与不可互换的失效}
\label{tab:modality-encoder-failures}
\begin{tabularx}{\textwidth}{p{0.15\textwidth}YYYY}
\toprule
模态 & 主要信息 & 常见编码 & 时间/坐标要求 & 典型失效 \\
\midrule
RGB & 纹理、语义 & CNN/ViT & 相机内外参、曝光时间 & 光照、遮挡、域偏移 \\
深度/点云 & 几何、距离 & 点网络、体素/稀疏卷积 & 深度标定、世界/机体系 & 空洞、反光、稀疏 \\
触觉 & 接触形变、滑移 & CNN/时序编码器 & 传感器局部系、接触时刻 & 漂移、磨损、饱和 \\
本体状态 & 关节、速度、力矩 & MLP/时序编码 & 关节顺序、单位、控制时间 & 编码器回零、延迟 \\
音频 & 冲击、摩擦、电机声 & 频谱/音频 Transformer & 麦克风同步、窗口 & 环境噪声、混响 \\
\bottomrule
\end{tabularx}
\end{table}

\section{连接器内部到底发生什么}

最简单的线性/MLP adapter 对每个视觉 token 投影到语言模型维度，保留 token 数。Cross-attention 连接器用少量可学习 query $Q_l$ 读取大量模态特征 $K_m,V_m$：
\begin{equation}
Z=\operatorname{softmax}\left(\frac{Q_lK_m^\top}{\sqrt d}\right)V_m.
\label{eq:latent-cross-attention}
\end{equation}
若 query 数为 $N_l\ll N_m$，它把 $N_m$ 个输入压成 $N_l$ 个 latent。Perceiver 用 latent bottleneck 迭代读取高维输入\cite{jaegle2021perceiver}；BLIP-2 的 Q-Former 用可学习 query 在冻结图像编码器与冻结语言模型之间提取视觉信息\cite{li2023blip2}。

\begin{figure}[H]
\centering
\setlength{\tabcolsep}{3pt}
\begin{tabular}{c@{\ $\rightarrow$\ }c@{\ $\rightarrow$\ }c@{\ $\rightarrow$\ }c}
\fbox{\begin{minipage}{0.17\textwidth}\centering 模态特征\\$N_m\times d_m$\end{minipage}} &
\fbox{\begin{minipage}{0.14\textwidth}\centering K/V\\投影\end{minipage}} &
\fbox{\begin{minipage}{0.17\textwidth}\centering $N_l$ 个\\latent query\end{minipage}} &
\fbox{\begin{minipage}{0.17\textwidth}\centering 共享 token\\$N_l\times d$\end{minipage}}
\end{tabular}
\caption{latent-query 连接器的内部计算。作者依据 Perceiver/Q-Former 的共同 cross-attention 结构改绘\cite{jaegle2021perceiver,li2023blip2}。}
\label{fig:latent-query-connector}
\end{figure}

压缩不是免费。$N_l$ 太少会丢失小物体或接触瞬间，太多则占据上下文和计算。query 训练目标若只受语言描述监督，可能保留“杯子”语义而丢掉杯沿姿态。连接器是否有效必须按下游定位、动作和接触任务评估。

\section{prefix、交错与 latent slot 接口}

视觉前缀把所有视觉 token 放在语言/动作之前，便于因果解码但弱化帧间对应；interleaving 按时间交错观测和动作，保留轨迹结构但序列更长；latent slot 固定每帧或每模态的 token 数，计算稳定但信息瓶颈更强；action token 把动作离散或嵌入同一词表/隐藏空间。

上下文预算为 $L_{\max}$ 时，若语言占 $L_l$、每帧视觉占 $L_v$、本体/触觉占 $L_s$，可保留帧数近似为
\begin{equation}
T_{\max}=\left\lfloor\frac{L_{\max}-L_l-L_{\mathrm{out}}}{L_v+L_s}\right\rfloor.
\end{equation}
例如 $L_{\max}=4096,L_l=256,L_{\mathrm{out}}=256,L_v=196,L_s=8$，只能容纳 17 帧。提高相机分辨率会直接缩短时间历史。

\begin{lstlisting}[caption={最小 latent-query 连接器教学伪代码}]
features, stamps = modality_encoder(raw_sensor)
assert calibration_and_time_valid(stamps)
keys = key_projection(features)
values = value_projection(features)
queries = learned_queries.expand(batch_size, num_latents, dim)
latents = softmax(queries @ keys.T / sqrt(dim)) @ values
latents = output_projection(latents)
return latents, stamps, modality_valid_mask
\end{lstlisting}

真实实现还需多头、残差、归一化和掩码；伪代码强调连接器输出必须保留时间戳和有效性，而不能只返回浮点张量。

\section{信息损失、模态鸿沟与时间对齐}

模态对齐常用对比或配对损失，但“语义靠近”不等于几何同构。文本“轻轻放下”没有唯一力轨迹，触觉接触也没有唯一语言描述。共享 latent 需要允许一对多关系，并由动作/任务损失补充约束。

传感到达时间不同会形成虚假跨模态关联。相机看到未接触，触觉已报告碰撞时，简单拼接让模型学到矛盾状态。应按测量时间对齐、保留延迟分布，并对缺失模态训练显式 mask。模态 dropout 只能提高已知缺失的鲁棒性，不能修复错误标定或未检测的陈旧数据。

\section{知识小结与综述判断}

知识小结：编码器从各模态提取任务相关特征；adapter 改变维度，Q-Former 与 Perceiver 类连接器用 latent query 压缩大量特征；prefix、交错和 slot 决定 token 与时间组织。token 预算把空间分辨率、相机数与历史长度直接耦合。

综述判断：结构图中的“connector”不是无损电缆，而是决定信息丢失和训练梯度的瓶颈。论文比较应报告输入 token、输出 latent、监督目标和下游任务，而不能只说“使用轻量适配器”。对机器人而言，时间戳、坐标系与失效 mask 与 hidden size 同等重要。

\chapter{从 VLM 到 VLA：语义落地与动作接口}
\label{chap:vlm-to-vla}

\section{视觉语言预训练迁移了什么}

图文预训练可迁移物体类别、属性、关系、开放词汇检索和语言组合先验。CLIP 的对比目标把配对图文拉近，使视觉特征能被自然语言查询\cite{radford2021clip}；PaLM-E 把连续传感表示注入语言模型，探索具身多模态推理\cite{driess2023palme}。这些能力可帮助理解“把红杯放到左侧托盘”，却不自动给出红杯六自由度姿态、机器人可达构型、抓取力和接触恢复。

VLA 必须把语义条件落到动作分布：
\begin{equation}
p(a_{t:t+H}\mid I_{\le t},l,p_{\le t},h_{\le t},e),
\label{eq:vla-conditional-action}
\end{equation}
其中 $I$ 是视觉、$l$ 是语言、$p$ 是本体状态、$h$ 是历史/触觉、$e$ 是本体与控制接口。遗漏 $e$ 会把“同一动作”误认为能跨机器人直接执行。

\begin{figure}[H]
\centering
\setlength{\tabcolsep}{2pt}
\begin{tabular}{c@{\ $\rightarrow$\ }c@{\ $\rightarrow$\ }c@{\ $\rightarrow$\ }c@{\ $\rightarrow$\ }c}
\fbox{\begin{minipage}{0.12\textwidth}\centering 看懂\\目标\end{minipage}} &
\fbox{\begin{minipage}{0.13\textwidth}\centering 定位/\\grounding\end{minipage}} &
\fbox{\begin{minipage}{0.13\textwidth}\centering 判断\\可供性\end{minipage}} &
\fbox{\begin{minipage}{0.14\textwidth}\centering 生成可执行\\动作\end{minipage}} &
\fbox{\begin{minipage}{0.13\textwidth}\centering 观测并\\恢复\end{minipage}}
\end{tabular}
\caption{从语义到闭环动作的五个证据门。作者教学绘制；任一门失败都不能由“语义理解正确”替代。}
\label{fig:semantic-to-action-gates}
\end{figure}

\section{grounding、可供性与本体状态}

grounding 要把词语对应到图像区域、三维对象或持续跟踪 ID。检测框不足以抓取：还需姿态、遮挡、几何和不确定性。对象在图像中可见也不等于机器人可达；可供性取决于本体、夹爪、障碍和当前状态。

SayCan 将语言模型给出的技能相关性与机器人价值/可供性结合，使语言上合理但机器人做不到的技能被降权\cite{ahn2022saycan}。抽象地可写为
\begin{equation}
s(k\mid l,s)=\log p_{\mathrm{LM}}(k\mid l)
+\lambda\log A(k\mid s,e),
\label{eq:language-affordance-score}
\end{equation}
$k$ 为技能，$A$ 是在当前状态和本体上成功的可供性。若语言概率为 0.9、可供性为 0，组合仍必须禁止执行；给 $A$ 加数值下限时尤其要避免把不可行动作重新赋予非零机会。

本体状态提供关节、末端、夹爪和控制模式。只看图像时，相同场景可能对应完全不同动作：机械臂正在向前运动还是静止、夹爪是否已闭合、上一动作是否被安全层截断，都决定下一步。状态还必须带单位、坐标系和时间戳。

\section{动作头、低层控制与闭环缺口}

VLA 可输出离散动作 token、连续末端增量、关节目标、动作块、技能调用或轨迹参数。接口越低层，训练数据越依赖具体本体和控制频率；接口越高层，越依赖可靠的技能库和规划器。没有“天然最统一”的动作空间。

RT-1 展示了语言/视觉条件下的大规模真实机器人动作 token 学习\cite{brohan2022rt1}。但输出动作仍要经过单位还原、限幅、安全过滤和低层控制。接触任务还需力/触觉反馈和恢复状态机；模型生成“向下 2 cm”并不能证明下方无障碍或插孔已对齐。

\begin{table}[H]
\centering
\caption{VLM 能力向 VLA 迁移时的证据缺口}
\label{tab:vlm-vla-transfer-gap}
\begin{tabularx}{\textwidth}{>{\raggedright\arraybackslash}p{0.19\textwidth}YY>{\raggedright\arraybackslash}p{0.28\textwidth}}
\toprule
已有能力 & 可迁移用途 & 仍缺信息 & 必要验证 \\
\midrule
开放词汇识别 & 找到语言目标候选 & 三维姿态、实例持续性 & 遮挡/同类实例 grounding \\
关系与常识 & 任务分解、候选技能 & 本体可达、动态约束 & 当前状态的可供性 \\
语言生成 & 解释或规划草案 & 动作单位与控制周期 & 编译到技能/动作接口 \\
视觉问答 & 描述场景与异常 & 接触力、隐蔽几何 & 触觉/状态闭环 \\
长序列建模 & 使用历史条件 & 状态新鲜度、错误记忆 & 延迟、恢复和回滚一致性 \\
\bottomrule
\end{tabularx}
\end{table}

\section{安全门与失败恢复}

VLA 的动作提议应与执行分离。执行前检查状态新鲜度、坐标/单位、关节/速度边界、碰撞和任务权限；执行中监控接触、超时和偏差；执行后验证预期状态是否发生。验证失败要触发重观测、撤回或请求人工，而不是无条件进入下一个语言步骤。

\begin{lstlisting}[caption={语义动作到可执行命令的安全门教学伪代码}]
proposal = vla.predict(observation, instruction, embodiment)
command = action_adapter.decode_units_frames(proposal)
if stale(observation) or not affordance_model.feasible(command):
    return reobserve_or_replan()
if not collision_and_limit_check(command):
    return safe_reject(reason="physical_constraint")
result = low_level_controller.execute(command, monitors=contact_limits)
if not verify_expected_state(result): recover_or_request_help()
\end{lstlisting}

安全过滤后的实际动作必须回写历史，否则模型以为自己的提议已执行，下一步条件错误。过滤频率高也说明策略与可行域不匹配，不能把安全层拦截当作部署成功。

\section{与 VLA 系统综述的边界}

本章只建立必要接口：语义先验、grounding、本体状态、动作表示和闭环验证。第 22 章建立统一分类，第 23 章比较代表性模型，第 24 章讨论数据/训练/泛化，第 25 章核验开源复现。这里不按模型逐一写架构和排行榜，避免同一论文在基础章与综述章重复介绍。

\section{知识小结与综述判断}

知识小结：VLM 可迁移开放词汇语义、关系和语言组合先验；VLA 还必须加入 grounding、本体状态、动作接口、物理约束和执行后验证。语言相关性与机器人可供性是不同信号。

综述判断：从“看懂”到“做对”的缺口不是再加一个动作头即可填平，而是整条闭环接口。模型论文若只报告语义任务或离线动作准确率，不能证明具身能力；真机证据应说明实际执行动作、低层控制、安全拦截和失败恢复，否则无法判断能力来自模型还是外部系统。

\chapter{上下文、外部记忆与长时程任务}
\label{chap:context-memory}

\section{三种时间尺度需要三种状态}

长时程任务不能把所有历史帧无限塞进上下文。高频物理状态服务控制，保留毫秒到秒；技能级事件记录“抓取失败、已重定位”，保留分钟；任务级状态记录目标、已完成步骤、资源与权限，可跨会话持久化。三者更新频率、可信度和写入权限不同。

\begin{figure}[H]
\centering
\begin{tabular}{ccc}
\fbox{\begin{minipage}{0.27\textwidth}\centering 高频状态\\位姿/速度/接触\\短窗、不可陈旧\end{minipage}} &
\fbox{\begin{minipage}{0.27\textwidth}\centering 事件记忆\\技能结果/异常\\可检索、带证据\end{minipage}} &
\fbox{\begin{minipage}{0.27\textwidth}\centering 任务状态\\目标/步骤/权限\\结构化、可持久化\end{minipage}}
\end{tabular}
\caption{长时程机器人记忆的三层结构。作者教学绘制；事件摘要不能覆盖实时状态，任务状态也不能由自由文本独占。}
\label{fig:three-layer-robot-memory}
\end{figure}

纯长上下文把历史 token 留在模型窗口，保真但成本随长度增加；外部检索只取相关片段，扩展容量但可能漏召回或污染；结构化任务状态用 schema 保存关键变量，可验证但表达灵活性较低。工程系统通常组合三者。

\section{历史压缩、事件摘要与检索}

摘要把长轨迹压成短文本或 latent。压缩必须保留任务所需事实、来源和不确定性。只写“抓取失败”不够，应记录对象 ID、姿态、夹爪命令、接触/视觉证据、失败分类和恢复动作。摘要模型可能把“不确定”改写成确定陈述，因此应链接原始事件包。

检索分数可组合语义、时间和任务过滤：
\begin{equation}
S(m_i,q,t)=\alpha\,\operatorname{sim}(e(m_i),e(q))
-\beta(t-t_i)+\gamma\,\mathbf 1[\operatorname{task}(m_i)=\operatorname{task}(q)].
\label{eq:memory-retrieval-score}
\end{equation}
时间衰减不应统一：安全禁令可能长期有效，瞬时位姿几百毫秒即过期。RAG 展示了参数模型结合外部检索的基本框架\cite{lewis2020rag}；机器人还需物理时间、坐标系和设备版本过滤。

ReAct 把推理轨迹与外部动作/观察交错\cite{yao2023react}，为工具调用和事件写入提供了组织方式。但自由文本思考不应作为唯一控制状态；模型可改写或遗漏先前结论，且文本没有强类型约束。

\section{结构化任务状态与恢复}

任务状态可写成有限状态机、行为树 blackboard 或事件溯源记录。状态转移应由可验证事件触发：
\begin{equation}
z_{k+1}=\delta(z_k,e_k),\qquad
e_k=(\text{type},\text{time},\text{evidence},\text{version}).
\end{equation}
若状态为 INSERTING，只有深度、力和视觉共同满足成功条件才转到 VERIFIED；单一语言模型判断不应覆盖硬传感条件。

恢复必须区分“恢复软件上下文”和“恢复物理状态”。进程重启后数据库可能记得物体已抓取，但机器人实际因断电松爪。系统应重新观测关键不变量，再决定从哪个技能继续。任务状态需包含模型/地图/标定版本，回滚模型时同步检查记忆兼容性。

\begin{lstlisting}[caption={带版本、来源与过期检查的记忆读写教学伪代码}]
event = observe_skill_result(raw_sensors, controller_log)
memory.append(event, source_links=raw_bundle,
              map_version=map.id, calibration=calibration.id)
task_state = transition(task_state, verified(event))

candidates = memory.retrieve(query, task_id=task_state.task_id)
usable = [m for m in candidates
          if not expired(m) and compatible_versions(m, current_system)]
context = compose(recent_physical_state(), usable, task_state)
if critical_invariant_unknown(context): reobserve_before_acting()
\end{lstlisting}

\section{上下文预算与写入治理}

上下文窗口应给当前观测、任务指令、检索历史和动作输出分别预算。若检索片段挤占当前图像 token，模型可能“记得过去却看不见现在”。可先保留不可压缩的当前状态和安全规则，再按得分选择事件，最后留出输出空间。

外部记忆还需要写入权限。传感器事实、控制执行结果、操作者备注和模型推断应使用不同字段；推断不能覆盖事实。错误信息若被模型生成、写入、再次检索，会形成自我强化污染。关键事件应需要规则或人工确认。

\begin{table}[H]
\centering
\caption{长上下文、检索记忆与结构化状态的职责比较}
\label{tab:memory-strategy-comparison}
\begin{tabularx}{\textwidth}{p{0.17\textwidth}YYYY}
\toprule
方式 & 优点 & 代价 & 典型风险 & 适用信息 \\
\midrule
纯长上下文 & 保留原始细节、实现直接 & 计算/内存随长度增长 & 位置稀释、窗口截断 & 最近对话、短轨迹 \\
摘要上下文 & token 少、可读 & 压缩不可逆 & 幻觉、遗漏不确定性 & 技能结果概述 \\
外部向量检索 & 容量大、按需读取 & 索引和检索链复杂 & 漏召回、污染、陈旧 & 历史案例、文档 \\
结构化任务状态 & 强类型、可验证、易恢复 & schema 设计成本 & schema 不全、迁移失败 & 步骤、权限、不变量 \\
事件溯源 & 可审计、可重放 & 日志量与重建成本 & 事件顺序/版本错误 & 控制与任务关键事件 \\
\bottomrule
\end{tabularx}
\end{table}

\section{错误记忆与长时程评测}

评测应主动注入：过期位姿、错误对象 ID、矛盾摘要、缺失恢复事件、地图版本变化和检索无关片段。指标包括关键事实召回、错误记忆采纳率、恢复后状态一致性、上下文延迟和任务成功。只在干净历史上跑长任务，无法证明系统能抵抗记忆污染。

语言模型可帮助解释和检索，但安全关键不变量应由结构化状态和传感证据维护。SayCan 等系统把高层语言规划与可执行技能结合\cite{ahn2022saycan}；长时程扩展仍需明确技能结果怎样写回、失败怎样影响后续候选。

\section{知识小结与综述判断}

知识小结：高频状态、事件记忆和任务状态对应不同时间尺度；长上下文保留细节，检索扩展容量，结构化状态提供可验证持久化。摘要、检索和状态转移都必须携带来源、时间与版本。

综述判断：机器人记忆不是“更长 prompt”，而是状态治理。真正危险的不是忘记一段描述，而是把过期或推断性信息当成当前物理事实。长时程系统的证据应包括断点恢复、版本回滚和错误记忆注入测试，而不只展示一次成功的长链推理。


\part{具身系统架构}
\chapter{感知、融合与状态接口}
\label{chap:perception-state-interface}

\section{本章任务：把感知结果变成可消费状态}

第 15 章讨论原始模态怎样编码为特征，本章只主讲这些特征怎样成为世界模型、规划器和控制器能够安全消费的状态。检测框、分割掩码、点云簇或触觉图像都是\emph{测量解释}，尚不是系统状态：它们通常缺少稳定身份、统一坐标系、测量时刻、估计不确定性和有效期。

一个最小状态包可写为
\begin{equation}
\mathcal{S}_t=\{\hat{\bm{x}}_t,\bm{P}_t,\mathcal{O}_t,\mathcal{C}_t,
\mathcal{F}_t,\mathcal{T}_t,v_{\mathrm{calib}},\bm{m}_t\},
\label{eq:state-packet}
\end{equation}
其中 $\hat{\bm{x}}_t$ 与 $\bm{P}_t$ 是机器人状态及协方差，$\mathcal{O}_t$ 是带身份的对象集合，$\mathcal{C}_t$ 是接触集合，$\mathcal{F}_t$ 是坐标变换，$\mathcal{T}_t$ 保存测量与发布时间，$v_{\mathrm{calib}}$ 是标定版本，$\bm{m}_t$ 是逐字段有效性掩码。式~\eqref{eq:state-packet} 的要点不是字段越多越好，而是每个消费者都能判断“值是什么、在何时何地成立、还能否使用”。

\section{从像素、点和接触到对象状态}

二维检测给出类别与框，实例分割进一步给出像素级对象区域。Mask R-CNN 的并行分类、框回归与掩码分支是典型实例分割结构\cite{he2017maskrcnn}，但类别概率不能直接回答对象的六自由度位姿、遮挡程度或可抓取区域。三维点集编码如 PointNet 通过对称聚合保持输入排列不变\cite{qi2017pointnet}；它仍不自动解决传感器外参、稀疏视角和跨帧身份。

对象状态应至少包含
\begin{equation}
o_i=(\mathrm{id}_i,c_i,{}^W\bm{T}_{o_i},\bm{v}_i,\Sigma_i,
\mathcal{A}_i,t_i,q_i),
\end{equation}
其中 $\mathcal{A}_i$ 表示可供性候选，例如“可抓取区域”或“可倾倒方向”，$q_i$ 表示质量与来源标志。可供性不是类别的固定属性：同一杯子在杯柄被遮挡、内部装有热液体或位于狭窄货架时，可执行动作不同。

触觉又提供另一类局部状态。GelSight 类传感器从弹性体表面形变估计接触几何与力\cite{yuan2017gelsight}；它适合判断接触、滑移和局部形状，却不能替代全局视觉定位。把触觉图像当普通 RGB 图像并做强裁剪，可能恰好删掉决定滑移的微小纹理。

\begin{table}[H]
\centering
\caption{感知输出进入状态包前必须补齐的语义}
\label{tab:perception-state-fields}
\begin{tabularx}{\textwidth}{p{0.13\textwidth}YYYY}
\toprule
来源 & 原始输出 & 必补字段 & 主要消费者 & 不可静默处理的失效 \\
\midrule
二维视觉 & 框、掩码、关键点 & 身份、深度、坐标、遮挡、时间 & 任务规划、抓取 & 类别正确但位姿错误 \\
三维感知 & 点、体素、表面 & 外参、法向、协方差、可见域 & 碰撞检查、定位 & 稀疏区被误当自由空间 \\
触觉/力觉 & 形变、力矩、滑移分数 & 接触点、工具坐标、带宽、饱和 & 柔顺控制、技能终止 & 饱和值被当真实力 \\
本体状态 & 关节、速度、电流 & 单位、顺序、零位、驱动器状态 & 估计器、控制器 & 版本错配仍数值合法 \\
\bottomrule
\end{tabularx}
\end{table}

\section{时间对齐、融合与统计门控}

多模态融合首先是时序问题。若相机帧的测量时刻为 $t_m$、状态包发布时间为 $t_p$，系统必须按运动模型把测量传播到共同参考时刻，或明确保留其原始时刻；直接拼接“最新值”会制造不存在的同步快照。第 8 章的滤波与平滑给出了估计基础，本章强调发布接口。

对测量残差 $\bm{r}=\bm{z}-h(\hat{\bm{x}})$ 和创新协方差 $\bm{S}$，常用 Mahalanobis 门控为
\begin{equation}
d^2=\bm{r}^{\mathsf T}\bm{S}^{-1}\bm{r}\leq \chi^2_{k,1-\alpha}.
\label{eq:mahalanobis-gate}
\end{equation}
式中 $k$ 是残差维数，$\alpha$ 是拒绝概率。该门只能说明测量与当前统计模型是否一致，不能证明对象身份或坐标语义正确；系统性外参偏差甚至可能长期“稳定地错误”。Kalman 滤波的协方差更新\cite{kalman1960new}和机器人状态估计中的流形、时间与噪声建模\cite{barfoot2024state}都要求这些假设显式化。

\begin{figure}[H]
\centering
\setlength{\tabcolsep}{3pt}
\begin{tabular}{c@{\ $\rightarrow$\ }c@{\ $\rightarrow$\ }c@{\ $\rightarrow$\ }c}
\fbox{\begin{minipage}{0.17\textwidth}\centering 带时间戳的\\模态观测\end{minipage}} &
\fbox{\begin{minipage}{0.16\textwidth}\centering 坐标与\\标定校验\end{minipage}} &
\fbox{\begin{minipage}{0.16\textwidth}\centering 关联、门控\\与融合\end{minipage}} &
\fbox{\begin{minipage}{0.18\textwidth}\centering 状态包\\+ 质量元数据\end{minipage}}
\end{tabular}
\caption{从测量到可消费状态的生产链。任何校验失败都应产生降级或拒绝事件，而非填入默认零值。}
\label{fig:state-production-pipeline}
\end{figure}

\section{缺失模态与降级发布}

缺失模态至少分为四类：传感器未发布、消息迟到、数值越界、语义质量不足。四者不能统一写成零。以触觉为例，零接触力与“触觉离线”含义相反；以前一帧填充又会把已结束接触延长。状态包应携带 `valid`、`age`、`source` 和 `reason`，消费者声明自己允许的最大陈旧时间。

\begin{lstlisting}[caption={状态包发布器的教学伪代码：显式处理迟到、标定与降级}]
def build_state_packet(target_time, measurements, calibration):
    packet = StatePacket(target_time=target_time)
    for name, msg in measurements.items():
        if msg is None:
            packet.invalidate(name, reason="missing")
            continue
        if target_time - msg.stamp > MAX_AGE[name]:
            packet.invalidate(name, reason="stale", age=target_time-msg.stamp)
            continue
        if msg.calibration_version != calibration.version:
            packet.invalidate(name, reason="calibration_mismatch")
            continue
        aligned = transform_and_propagate(msg, calibration, target_time)
        if not innovation_gate(aligned):
            packet.invalidate(name, reason="statistical_outlier")
            continue
        packet.update(name, aligned.value, aligned.covariance,
                      source=msg.sensor_id, stamp=msg.stamp)
    packet.mode = choose_degraded_mode(packet.validity_mask)
    return packet
\end{lstlisting}

\section{工程失效：数值合法不等于语义正确}

最危险的接口故障往往不会产生 NaN：毫米与米混用、相机坐标与世界坐标混用、旧外参配新数据、对象 ID 在遮挡后交换，都能产生范围合理的数值。因而测试应同时覆盖单位、坐标变换闭环、时间单调性、标定版本和对象身份连续性。状态包还应记录最近一次重定位、测量拒绝与传感器降级原因，使第 21 章的运行时监督器能够追溯。

\section{知识小结与综述判断}

本章把感知系统的输出从“模型预测张量”提升为带契约的状态产品。核心链条是：模态解释产生候选测量，时间与坐标校验建立可比性，关联与融合给出估计和不确定性，状态包再把有效性与来源交给下游。综述判断是：衡量感知模块不能只看离线 mAP 或分割精度，还要看状态在闭环中的陈旧度、校准性、身份稳定性和降级可用性。


\chapter{世界模型、环境表征与可预测状态}
\label{chap:world-models}

\section{先回答：模型预测什么，供谁决策}

“世界模型”常被泛化为任何能生成未来画面的模型。本章采用更严格的系统定义：若模型学习环境在动作作用下的状态转移，并且预测结果实际进入动作选择、风险评估或恢复，它才是决策链中的世界模型。高保真仿真器、控制器内部动力学模型、视频生成器和语义表征模型都可能提供资产，但用途不同。

\begin{table}[H]
\centering
\caption{四类易被混称为世界模型的系统}
\label{tab:world-model-types}
\begin{tabularx}{\textwidth}{p{0.17\textwidth}YYYY}
\toprule
类型 & 主要输入/输出 & 优化目标 & 决策接入方式 & 典型盲区 \\
\midrule
物理仿真器 & 状态、动作 $\to$ 数值轨迹 & 物理一致与数值稳定 & 离线训练、测试、MPC & 资产与现实偏差 \\
控制模型 & 局部状态、控制 $\to$ 短期状态 & 局部预测误差 & 高频优化与稳定控制 & 长期语义变化 \\
视频生成模型 & 图像、条件 $\to$ 未来帧 & 感知似然或生成质量 & 候选可视化、表征学习 & 动作可控性与物理可行性 \\
决策世界模型 & 历史、动作 $\to$ 潜状态/奖励/终止 & 决策充分的预测 & 想象 rollout、价值或风险 & 模型偏差被策略利用 \\
\bottomrule
\end{tabularx}
\end{table}

\section{潜状态空间：从观测重建到动作条件转移}

对观测 $\bm{o}_t$、动作 $\bm{a}_t$ 和潜状态 $\bm{z}_t$，典型结构分解为
\begin{align}
\bm{z}_t &\sim q_\phi(\bm{z}_t\mid \bm{z}_{t-1},\bm{a}_{t-1},\bm{o}_t),\\
\tilde{\bm{z}}_{t+1} &\sim p_\theta(\bm{z}_{t+1}\mid \bm{z}_t,\bm{a}_t),\\
\hat{\bm{o}}_t &\sim p_\theta(\bm{o}_t\mid\bm{z}_t),\qquad
(\hat r_t,\hat c_t)=g_\theta(\bm{z}_t).
\label{eq:latent-world-model}
\end{align}
$q_\phi$ 把新观测吸收到后验状态，$p_\theta$ 在没有未来观测时执行想象转移，$g_\theta$ 预测奖励和继续概率。PlaNet 用潜动力学从像素执行规划\cite{hafner2019planet}；DreamerV3 则在学习的世界模型中想象未来并训练 actor--critic\cite{hafner2023dreamerv3}。

潜状态的目标不是压缩率最大，而是保留决策所需变量。重建损失偏爱占像素面积大的背景，接触建立、摩擦变化或细小障碍却可能只占少数像素。可增加奖励、终止、接触和约束违反等辅助头，但每个头都会把训练数据中的标注偏差带入状态。

\begin{figure}[H]
\centering
\setlength{\tabcolsep}{3pt}
\begin{tabular}{c@{\ $\rightarrow$\ }c@{\ $\rightarrow$\ }c@{\ $\rightarrow$\ }c}
\fbox{\begin{minipage}{0.16\textwidth}\centering 历史观测\\与状态包\end{minipage}} &
\fbox{\begin{minipage}{0.15\textwidth}\centering 后验潜状态\\$\bm z_t$\end{minipage}} &
\fbox{\begin{minipage}{0.17\textwidth}\centering 候选动作下\\想象 rollout\end{minipage}} &
\fbox{\begin{minipage}{0.18\textwidth}\centering 回报、风险\\与终止预测\end{minipage}}
\end{tabular}
\caption{世界模型进入闭环的最小链路。若最后的预测不改变候选动作排序，模型只是辅助表征而非在线决策组件。}
\label{fig:world-model-decision-loop}
\end{figure}

\section{多步 rollout：模型误差如何累积}

候选动作序列 $\bm{a}_{t:t+H-1}$ 的风险调整目标可写为
\begin{equation}
J=\mathbb{E}\!\left[\sum_{k=0}^{H-1}\gamma^k \hat r_{t+k}\right]
-\lambda\sum_{k=1}^{H}u(\tilde{\bm z}_{t+k})
-\rho\Pr(\text{constraint violation}),
\label{eq:risk-aware-imagination}
\end{equation}
其中 $u(\cdot)$ 是模型集成分歧、预测方差或分布外分数。$H$ 越长并不必然越好：一步小偏差会通过策略选择进入未见区域，随后模型可能给出高回报但不真实的“幻想”。因此工程系统常采用短视窗重规划、真实观测校正和不确定性阈值，而不是一次生成完整长轨迹。

\begin{lstlisting}[caption={带不确定性和硬约束的想象规划教学伪代码}]
best = None
for action_sequence in propose_candidates(state):
    latent = world_model.encode(state)
    score, rejected = 0.0, False
    for action in action_sequence:
        latent_ensemble = world_model.predict_ensemble(latent, action)
        reward = mean_reward(latent_ensemble)
        uncertainty = ensemble_disagreement(latent_ensemble)
        if uncertainty > U_MAX or predicts_hard_violation(latent_ensemble):
            rejected = True
            break
        score += reward - RISK_WEIGHT * uncertainty
        latent = mean_latent(latent_ensemble)
    if not rejected and (best is None or score > best.score):
        best = Candidate(action_sequence, score)
return best.first_action if best else request_safe_recovery()
\end{lstlisting}

\section{像素生成、JEPA 与抽象预测的目标差异}

像素生成要求解释纹理、光照和背景，优点是可视化直观，代价是大量容量用于与动作无关的细节。JEPA 类方法在表征空间预测被遮挡区域；I-JEPA 明确采用非生成式目标，由上下文块预测目标块的表征\cite{assran2023ijepa}。这能学习语义表征，但原始 I-JEPA 没有动作条件、奖励或终止变量，不能因“会预测表征”就直接称为机器人动力学世界模型。

抽象预测还可以面向对象关系、接触模式或任务谓词。例如“杯子在托盘上”比每个像素更适合长时程规划；但抽象器若漏掉杯沿裂纹，规划器看不到安全风险。可靠系统往往同时维护多尺度状态：控制层使用连续局部模型，任务层使用对象/谓词状态，视觉模型提供更新证据。

\section{怎样证明世界模型真的有用}

有效证据至少包含三层：一是单步与多步预测在任务相关变量上的校准误差；二是接入规划后相对无模型、真实模型或消融基线的成功率与样本效率；三是分布外、接触突变和传感缺失时的风险与回退。只展示逼真的未来视频不能证明控制收益，只展示平均回报也不能说明失败是否来自模型、规划器或执行器。

\section{知识小结与综述判断}

世界模型把“当前状态”扩展为“候选动作下的可预测状态”。其技术主线是后验编码、动作条件转移、任务相关预测和想象决策。综述判断是：真正的分水岭不在模型是否生成视频，而在预测是否具有动作条件、是否传播不确定性、是否进入闭环，并能在模型不可信时触发真实观测校正或安全恢复。


\chapter[任务规划、结构化中间表示与具身推理]{任务规划、结构化中间表示\\与具身推理}
\label{chap:task-planning-ir}

\section{规划的产物不应只是自然语言步骤}

自然语言目标“把冰箱里的饮料拿给客厅的人”省略了对象身份、门状态、抓取约束和失败处理。任务规划的职责，是把意图编译成可由技能库执行、可在状态包上检查、可在失败后恢复的中间表示。第 16 章主讲语义怎样落地到动作接口；本章主讲落地之后的结构、约束和执行反馈，不讨论低层控制细节。

把技能 $\sigma$ 写成契约
\begin{equation}
\sigma=(\mathrm{Pre},\mathrm{Param},\mathrm{Run},\mathrm{Post},
\mathrm{Fail},\mathrm{Compensate}),
\label{eq:skill-contract}
\end{equation}
其中前置条件决定能否启动，后置条件决定是否成功，失败集合区分可重试、需重规划和必须停机的情形，补偿动作描述已经改变环境后如何回退。没有这些字段，所谓“任务分解”只是一串无法核验的文本。

\section{四类中间表示及其执行语义}

经典自动规划以状态、动作前置条件和效果搜索计划\cite{ghallab2004automatedplanning}。HTN 通过方法把复合任务递归分解，适合表达领域程序；行为树通过周期 tick 和 Success/Failure/Running 返回值组织反应式执行\cite{colledanchise2018behaviortrees}；程序化表示能表达循环、几何计算和 API 组合，但必须面对类型、权限和终止性。

\begin{table}[H]
\centering
\caption{任务中间表示的能力与代价}
\label{tab:task-ir-comparison}
\begin{tabularx}{\textwidth}{p{0.15\textwidth}YYYY}
\toprule
表示 & 强项 & 执行反馈 & 易验证内容 & 主要风险 \\
\midrule
文本步骤 & 生成灵活、便于人读 & 通常非结构化 & 很少 & 歧义、漏条件、无法确定调用 \\
HTN/符号 & 层级分解、前后条件清晰 & 状态谓词更新 & 可达性、方法约束 & 符号状态与现实脱节 \\
行为树 & 反应式 tick、局部回退 & 三值或扩展状态 & 节点契约、局部安全 & 大树维护与共享状态副作用 \\
受限程序 & 循环、计算、工具组合 & 异常、返回值、日志 & 类型、权限、静态规则 & 任意代码、超时、API 误用 \\
\bottomrule
\end{tabularx}
\end{table}

行为树“可反应”不是因为图画成树，而是控制节点反复 tick 子节点。例如 Sequence 遇到 Failure 立即失败，遇到 Running 保留进度；Fallback 遇到 Success 即返回，失败才尝试下一分支。若叶节点不检查新状态，树仍会机械执行旧计划。

\section{语言相关性必须乘上可执行性}

SayCan 的关键思想是让语言模型对技能与指令的相关性，受到机器人在当前状态下执行该技能的价值/可供性约束\cite{ahn2022saycan}。对候选技能 $\sigma_i$ 可抽象为
\begin{equation}
i^*=\arg\max_i\left[
\log p_{\mathrm{LM}}(\sigma_i\mid g,h)+
\beta\log p_{\mathrm{exec}}(\mathrm{success}\mid \mathcal S_t,\sigma_i)
\right],
\label{eq:ch20-language-affordance-score}
\end{equation}
其中 $g$ 是目标，$h$ 是已执行历史。两项不能互相替代：语言高分不代表机器人够得到，价值高分也不代表动作与任务相关。若安全约束为硬条件，应在排序前直接过滤，而不是仅以小惩罚加入分数。

Code as Policies 进一步让语言模型生成调用感知与控制 API 的程序，可表达反馈循环和几何计算\cite{liang2023codepolicies}。这种路线的关键不是“模型会写 Python”，而是暴露给模型的 API 是否受限、参数是否类型化、执行是否在沙箱内、每次调用是否由状态与安全监督器复核。

\begin{figure}[H]
\centering
\setlength{\tabcolsep}{3pt}
\begin{tabular}{c@{\ $\rightarrow$\ }c@{\ $\rightarrow$\ }c@{\ $\rightarrow$\ }c}
\fbox{\begin{minipage}{0.16\textwidth}\centering 自然语言目标\\与当前状态\end{minipage}} &
\fbox{\begin{minipage}{0.16\textwidth}\centering 候选技能\\及契约\end{minipage}} &
\fbox{\begin{minipage}{0.17\textwidth}\centering 约束检查后的\\BT/HTN/程序\end{minipage}} &
\fbox{\begin{minipage}{0.18\textwidth}\centering 执行事件\\写回与重规划\end{minipage}}
\end{tabular}
\caption{从语言意图到受约束技能调用的编译链。结构化中间表示位于生成模型与运行时之间。}
\label{fig:task-compilation-chain}
\end{figure}

\section{完整例子：取饮料任务的状态与失败分支}

初始状态为“机器人在厨房外、冰箱门未知、目标饮料身份不确定”。规划器不能直接输出 `grasp(drink)`，而应先定位冰箱、确认门状态、打开门、重新感知、绑定对象身份、验证抓取条件，再执行取放。下面的伪代码刻意把感知写回和补偿分支放在主路径内。

\begin{lstlisting}[caption={受限任务程序：显式前置条件、写回与失败分支}]
def fetch_drink(goal, state):
    require(goal.destination == "living_room")
    navigate("fridge", timeout=60)
    state = observe_state(required=["fridge_door", "drink_candidates"])

    if state.fridge_door == "closed":
        result = run_skill("open_fridge", force_limit=SAFE_FORCE)
        if result.code == "handle_not_found":
            return replan_with_active_perception(result.evidence)
        if not result.success:
            return escalate("door_open_failed", result.log)

    state = observe_state(required=["drink_candidates"])
    target = bind_object(goal.description, state.drink_candidates)
    if target.confidence < ID_THRESHOLD:
        return ask_human_to_disambiguate(target.alternatives)

    grasp = plan_grasp(target, constraints=["upright", "collision_free"])
    result = run_skill("pick", object_id=target.id, grasp=grasp)
    if result.code == "slip":
        return run_skill("place_safe_and_regrasp", object_id=target.id)
    require(observe_state().held_object_id == target.id)
    return navigate_and_handover(goal.destination, target.id)
\end{lstlisting}

该程序不是为了展示语法，而是展示五个可审计点：目标被类型化；对象 ID 在感知后绑定；技能调用带约束；失败码决定不同恢复；后置条件由新观测确认。若生成模型只负责候选结构，静态检查器和运行时仍可拒绝越权 API、无界循环或缺失失败分支。

\section{多机器人协调不是多份独立计划}

多个机器人共享对象、通道和工具时，局部可行不等于联合可行。中间表示需增加资源锁、时序约束、通信超时和责任转移。例如一个机器人交付工件后，另一个机器人只有在视觉确认工件进入交接区并取得所有权后才能抓取。语言模型可以生成协作候选，但冲突检测、时间窗和安全互锁应由结构化调度器执行。

\section{知识小结与综述判断}

任务规划的核心产物是可检查的执行结构，而非听起来合理的解释。技能契约定义可执行单元，HTN/行为树/受限程序组织分解与反馈，语言相关性和当前可执行性共同筛选候选，运行事件再写回状态。综述判断是：评价“具身推理”应优先检查计划能否绑定真实对象、通过前后条件、覆盖失败分支并被运行时拒绝，而不是仅评价自然语言推理链是否流畅。

\chapter[技能、控制、安全监督与运行时恢复]{技能、控制、安全监督\\与运行时恢复}
\label{chap:runtime-recovery}

\section{运行时不是一条从模型直达电机的箭头}

第 20 章产生受约束技能调用，本章讨论调用怎样成为电机命令并在异常时停止。完整运行时至少包含：任务执行器、技能服务器、轨迹/策略、低层控制器、安全过滤器、独立监控器、事件日志和恢复管理器。VLA 或任务规划器只生产候选动作/技能，不拥有绕过速度、力、空间和权限约束的权力。

技能接口应声明输入状态版本、参数单位、允许工作区、最大时长、终止判据、抢占语义和失败码。低层控制器则在更高频率上跟踪目标或调节交互。阻抗控制通过期望的质量--阻尼--刚度关系管理位姿误差与外力\cite{hogan1985impedance}，操作空间控制在任务坐标中组织动力学\cite{khatib1987operational}；二者都不能判断“是否抓错人手中的杯子”。

\section{不同时间尺度必须有明确所有权}

\begin{table}[H]
\centering
\caption{具身运行时的时间尺度、权限与失败响应}
\label{tab:runtime-timescales}
\begin{tabularx}{\textwidth}{p{0.17\textwidth}YYYY}
\toprule
层级 & 典型周期 & 输出 & 允许的抢占者 & 典型失败响应 \\
\midrule
任务规划 & 秒至分钟 & 技能图/程序 & 人、策略监督器 & 重规划、请求协助 \\
技能/策略 & 20--500 ms & 动作块、目标位姿 & 安全监督器、超时器 & 重试、切换技能、撤退 \\
轨迹与控制 & 0.5--10 ms & 力矩、速度或位置命令 & 硬实时安全层 & 限幅、保持、受控停机 \\
驱动与硬件保护 & 微秒至毫秒 & 电流、制动 & 硬件联锁 & 断能、制动、故障锁存 \\
\bottomrule
\end{tabularx}
\end{table}

高层“停止”必须被翻译为低层可实现的停机轨迹，而不是突然令速度为零并假设机械系统瞬时静止。反之，驱动器故障也必须向上传播为结构化事件，否则任务执行器可能不断重发同一技能。

\section{安全过滤：把建议动作投影回可行集}

令安全集合为 $\mathcal C=\{\bm x:h(\bm x)\geq0\}$。控制屏障函数方法要求存在 $\alpha$ 使
\begin{equation}
\dot h(\bm x,\bm u)+\alpha(h(\bm x))\geq0.
\end{equation}
对仿射系统 $\dot{\bm x}=f(\bm x)+g(\bm x)\bm u$，可把学习策略建议 $\bm u_{\mathrm{nom}}$ 投影为
\begin{align}
\bm u^*=\arg\min_{\bm u}\;&\frac12\|\bm u-\bm u_{\mathrm{nom}}\|_{\bm W}^2\\
\text{s.t. }\;&\nabla h^{\mathsf T}(f+g\bm u)+\alpha(h)\geq0,\quad
\bm u\in\mathcal U.
\label{eq:cbf-filter}
\end{align}
该二次规划形式把“尽量接近策略”与“保持安全集合”分开\cite{ames2017cbf}。但它依赖状态估计、模型和约束可行性：障碍距离过期、制动模型错误或多个约束冲突时，优化器可能无解。无解分支必须定义受控停机或硬件保护，不能回退到未经检查的原动作。

\begin{figure}[H]
\centering
\setlength{\tabcolsep}{3pt}
\begin{tabular}{c@{\ $\rightarrow$\ }c@{\ $\rightarrow$\ }c@{\ $\rightarrow$\ }c}
\fbox{\begin{minipage}{0.16\textwidth}\centering 规划器/VLA\\候选命令\end{minipage}} &
\fbox{\begin{minipage}{0.16\textwidth}\centering 技能契约\\与状态校验\end{minipage}} &
\fbox{\begin{minipage}{0.16\textwidth}\centering 安全投影\\与控制器\end{minipage}} &
\fbox{\begin{minipage}{0.19\textwidth}\centering 执行反馈、事件\\与恢复状态机\end{minipage}}
\end{tabular}
\caption{运行时闭环。安全监督和恢复从状态与执行反馈独立取证，而不是只相信模型自报成功。}
\label{fig:runtime-supervision-loop}
\end{figure}

\section{Simplex 思路：复杂性能控制与简单安全控制分离}

Simplex 架构用高性能但难以完全验证的控制器处理正常任务，同时保留可验证的基线控制器；决策模块在安全边界前切换\cite{sha2001simplex}。在学习机器人中，高性能控制器可以是 VLA、扩散策略或 MPC，基线控制器可以是减速保持、撤退到已知自由空间或受控放置。

切换条件应包含前瞻裕度，而非等到约束已经违反。例如以制动距离 $d_b(v)$ 和感知/计算延迟 $\tau$ 估计
\begin{equation}
d_{\mathrm{required}}=d_b(v)+v\tau+d_{\mathrm{margin}}.
\end{equation}
当障碍净距小于 $d_{\mathrm{required}}$ 时即抢占。这里的 $\tau$ 必须包括相机曝光、网络推理、队列、控制下发和执行器响应，而不只是模型推理时间。

\section{异常分类决定恢复，而不是统一“再试一次”}

恢复前先分类：可瞬态重试（短暂遮挡）、需主动感知（对象身份不确定）、需重规划（通道被占）、需补偿（抓取后目标变更）、需人工接管（风险/权限不明）、需锁存停机（驱动或安全链故障）。盲目重试可能把轻微滑移变成掉落，把碰撞接触变成持续施力。

\begin{lstlisting}[caption={运行时监督器教学伪代码：抢占、降级与恢复}]
def supervise(command, state, health):
    event = validate(command, state, health)
    if event.is_hard_fault:
        hardware_stop()
        log_latched(event)
        return STOPPED
    if event.is_stale_state or event.identity_uncertain:
        preempt_current_skill()
        return request_active_perception(event)

    safe_command, feasible = safety_filter(command, state)
    if not feasible:
        switch_to_verified_baseline("controlled_stop")
        return RECOVERING

    result = execute_with_deadline(safe_command)
    log_event(command, safe_command, state.version, result)
    if result.success and postcondition_observed(result.skill_id):
        return SUCCEEDED
    if result.code in RETRYABLE and retry_budget(result.skill_id) > 0:
        return retry_with_updated_state(result)
    if result.code in REPLANNABLE:
        return request_replan(result.evidence)
    switch_to_verified_baseline("safe_retreat")
    return request_human(result)
\end{lstlisting}

事件日志至少记录：原始候选命令、过滤后命令、状态包版本、约束裕度、抢占者、失败码、传感证据和恢复结果。仅记录“任务失败”无法支持第 26 章的数据回采，也无法区分模型错误、接口错误和硬件错误。

\section{系统案例：抓取策略只是运行时的一环}

假设 VLA 输出未来 16 步末端动作。技能层先检查目标对象 ID 和工作区；轨迹层把动作块转换为时间参数化参考；安全层按最新障碍与力限值投影；控制器执行第一小段；视觉和触觉确认接近、接触与滑移；若对象身份变化则抢占，若轻微滑移则降低速度并调整夹持，若力传感器饱和则受控撤退。模型的贡献是提供候选行为，系统成功来自每一层都能拒绝、解释和恢复。

\section{知识小结与综述判断}

具身运行时把任务结构变成受约束的物理执行。技能契约明确边界，多时间尺度调度分配权限，安全过滤修改或拒绝建议动作，Simplex 式基线控制提供降级路径，异常分类和事件日志支撑恢复与回采。综述判断是：任何“端到端”模型只要进入真实机器人，就必须回答谁能抢占、状态过期怎么办、约束无解怎么办、失败后如何确认环境；答不出这些问题，就还不是完整系统。



\part{VLA 与机器人基础模型综述}
\chapter[视觉—语言—动作模型的问题设定与分类]{视觉—语言—动作模型\\的问题设定与分类}
\label{chap:vla-taxonomy}

\section{统一问题设定：先写清五个空间}

VLA 不是单一网络结构，而是一类条件策略。比较前应同时声明观测窗口 $\mathcal O$、任务条件 $\mathcal G$、本体状态 $\mathcal S$、动作空间 $\mathcal A$ 与部署闭环。统一写成
\begin{equation}
\boldsymbol a_{t:t+H_p-1}\sim
\pi_\theta(\cdot\mid \boldsymbol o_{t-H_o+1:t},
\boldsymbol s_{t-H_o+1:t},g),
\label{eq:vla-unified-setting}
\end{equation}
其中 $H_o$ 是观测历史，$H_p$ 是预测视窗。还必须给出实际执行步数 $H_e\leq H_p$；系统每执行 $H_e$ 步便用新观测重规划。只说“模型输出动作块”而不报告 $(H_o,H_p,H_e)$，无法判断延迟、平滑与纠错能力。

动作向量本身还需声明坐标系和语义。例如七维输出可表示末端相对位姿增量加夹爪，也可表示六关节位置加夹爪；两者维数相同却不可互换。RT-1/RT-2 的动作 token 路线\cite{brohan2022rt1,brohan2023rt2}、OpenVLA 的均匀分箱\cite{kim2024openvla}、Octo 的连续动作块\cite{octo2024}和 Diffusion Policy 的去噪轨迹\cite{chi2023diffusion}必须放在同一控制接口下比较。

\begin{figure}[H]
\centering
\setlength{\tabcolsep}{3pt}
\begin{tabular}{c@{\ $\rightarrow$\ }c@{\ $\rightarrow$\ }c@{\ $\rightarrow$\ }c}
\fbox{\begin{minipage}{0.16\textwidth}\centering 图像、语言\\本体历史\end{minipage}} &
\fbox{\begin{minipage}{0.15\textwidth}\centering 共享主干\\条件表征\end{minipage}} &
\fbox{\begin{minipage}{0.17\textwidth}\centering 回归/token\\扩散/流式头\end{minipage}} &
\fbox{\begin{minipage}{0.19\textwidth}\centering 反归一化、执行\\$H_e$ 步并重规划\end{minipage}}
\end{tabular}
\caption{VLA 的统一输入—动作—执行坐标。动作头之后仍有反归一化、坐标变换与低层控制。}
\label{fig:vla-unified-pipeline}
\end{figure}

\section{动作表示的六条坐标轴}

第 23 章详细解释四类代表路线，本节给出比较字段。动作表示至少沿六轴描述：连续/离散，单步/动作块，确定/多峰，逐步生成/并行生成，关节/末端/技能空间，以及开环/滚动闭环。

\begin{table}[H]
\centering
\caption{动作头不能脱离执行接口比较}
\label{tab:vla-action-heads}
\begin{tabularx}{\textwidth}{p{0.15\textwidth}YYYY}
\toprule
动作头 & 训练目标 & 推理方式 & 优点 & 主要代价 \\
\midrule
连续回归 & MSE/似然 & 一次前向 & 快、接口直接 & 多峰行为被平均 \\
离散 token & 交叉熵 & 自回归或并行 & 复用语言解码器 & 量化误差、词表约定 \\
扩散动作块 & 去噪匹配 & 多步迭代 & 表达多峰、轨迹一致 & 采样时延、调度敏感 \\
流匹配 & 速度场匹配 & ODE/少步积分 & 连续路径、可并行块 & 求解步数与稳定性 \\
技能 token & 分类/层级目标 & 技能选择 & 长时程抽象 & 依赖技能库与终止器 \\
\bottomrule
\end{tabularx}
\end{table}

离散动作常先按训练集分位数或固定范围归一化，再量化为 token。分箱宽度为 $\Delta_j$ 时，第 $j$ 维反量化误差在均匀箱内满足 $|e_j|\leq\Delta_j/2$；但控制误差还取决于坐标雅可比和低层跟踪。FAST 用频域压缩改善动作 token 效率\cite{pertsch2025fast}，仍不能消除机器人间动作语义差异。

扩散与流式策略直接建模条件动作分布。前者学习噪声/得分并迭代去噪，后者学习从简单分布到数据分布的速度场；$\pi_0$ 是 VLA 中采用 flow matching 的代表\cite{black2024pi0}。二者共同优势是动作块的联合分布，差异在训练参数化和数值求解，不应统称为“连续回归”。

\section{训练目标与闭环位置}

行为克隆的一般负对数似然为
\begin{equation}
\mathcal L_{\mathrm{act}}=-\mathbb E_{(o,s,g,a)\sim D}
\log \pi_\theta(a\mid o,s,g).
\end{equation}
语言辅助、视觉语言预训练或价值目标可以共享主干，但动作损失对应的数据分布仍决定策略学到什么。训练目标必须连同动作有效掩码、机器人 ID、归一化统计和任务采样权重解释。

闭环也有三种层次：论文中的逐步环境交互；动作块执行后的滚动重规划；以及第 21 章安全控制器在更高频率上的独立闭环。VLA 输出频率不是电机控制频率。若模型 10 Hz 输出目标位姿，1 kHz 阻抗控制器仍承担绝大多数瞬时稳定性。

\begin{lstlisting}[caption={动作块滚动执行的最小协议}]
while not task_done():
    state = get_versioned_state(required_modalities)
    if state.stale or not state.calibration_valid:
        return request_recovery("invalid_state")
    normalized = adapter.normalize(state, robot_id)
    chunk = policy.sample(normalized, instruction)
    action = adapter.denormalize_and_transform(chunk[:H_EXEC])
    for command in action:
        safe_command = runtime_supervisor.filter(command, state)
        execute(safe_command)
        if runtime_supervisor.preempted():
            break
\end{lstlisting}

\section{评测坐标：开环误差不等于闭环能力}

开环 MSE 或 token 准确率能诊断拟合，却不能评价多峰动作是否可行；仿真闭环能批量比较，但受动力学和视觉域差异影响；真机成功率最接近部署，却高度依赖初始条件、失败判据和试验次数。第 24 章因此把证据按任务、对象、场景、本体和恢复能力拆开。

\section{知识小结与综述判断}

VLA 的统一比较单位不是模型名，而是“输入窗口—任务条件—动作语义—动作分布—执行视窗—低层闭环”。本章建立坐标，第 23 章在坐标上比较代表模型。综述判断是：凡是缺少坐标系、归一化、$H_p/H_e$、控制频率或失败处理的模型介绍，都不足以支持部署层结论。

\chapter{代表性 VLA 的设计比较}
\label{chap:vla-design-comparison}

\section{问题设定：比较的是闭环策略，不是模型名气}

视觉--语言--动作模型（Vision--Language--Action, VLA）把观测 $o_t$、语言条件 $l$、本体状态 $s_t$ 映射为动作或动作块。统一写成
\begin{equation}
a_{t:t+H-1}\sim\pi_\theta(\,\cdot\mid o_{t-k:t},s_{t-k:t},l\,),
\label{eq:vla-policy}
\end{equation}
其中 $H$ 是预测动作块长度，$k$ 是历史窗口。式\eqref{eq:vla-policy}隐藏了会直接改变系统行为的六个选择：动作坐标系、动作分布族、时间离散、解码次数、执行多少步后重规划，以及低层控制器如何跟踪动作。只比较参数量或论文成功率会把这些差异抹平。

RT-1 证明了大规模多任务真机数据与高容量序列模型可以形成可扩展的机器人策略\cite{brohan2022rt1}；RT-2 进一步把动作表示并入 VLM 的 token 输出空间\cite{brohan2023rt2}。但“把动作写成 token”只是一个设计分支。Diffusion Policy 和 Octo 直接建模连续动作分布\cite{chi2023diffusion,octo2024}，$\pi_0$ 使用 flow matching\cite{black2024pi0}，GR00T N1 则把视觉语言推理与 diffusion transformer 动作模块组织成双系统\cite{nvidia2025groot}。

\section{动作表示的四种基本路线}

\begingroup
\hbadness=10000
\begin{figure}[H]
\centering
\includegraphics[width=\textwidth]{v0.3-23-VLA动作表示坐标.png}
\caption{动作表示、动作块与低层闭环之间的关系。作者依据 RT-2、Diffusion Policy、FAST 与 $\pi_0$ 的公开方法改绘\cite{brohan2023rt2,chi2023diffusion,pertsch2025fast,black2024pi0}。}
\label{fig:vla-action-design-space}
\end{figure}

\subsection{逐维回归：快，但容易把多峰行为平均掉}

若策略输出高斯均值，常用目标为
\begin{equation}
\mathcal L_{\mathrm{MSE}}=\mathbb E\|a_t-f_\theta(o_t,s_t,l)\|_2^2.
\end{equation}
它只需一次前向传播，适合低延迟控制；但当绕障可以向左也可以向右、抓取存在多个等价姿态时，条件均值可能落在两种可行动作之间。Mixture Density、离散 latent 或生成式动作头的目的，不是“让模型更大”，而是恢复条件动作分布的多峰结构。

这个结论可以直接从损失函数得到。对固定条件 $c=(o_t,s_t,l)$，把标量输出记为 $y$，则
\begin{equation}
\frac{\partial}{\partial y}\,\mathbb E[(a-y)^2\mid c]
=2\left(y-\mathbb E[a\mid c]\right),
\end{equation}
所以平方误差的最优点估计必为条件均值 $y^*=\mathbb E[a\mid c]$。若演示者以相同概率选择 $a=-1$ 和 $a=+1$，两种动作都可行，回归器却会输出 $0$。在绕障任务中，$0$ 可能恰好指向障碍物；这不是训练不足，而是“单峰点估计”与“多峰决策”之间的结构冲突。

\begin{figure}[H]
\centering
\includegraphics[width=0.92\textwidth]{v0.3-23-回归多峰平均.pdf}
\caption{平方误差把双峰动作分布压成条件均值的教学示意。曲线为作者生成的合成分布，不是论文实验结果。}
\label{fig:regression-mode-averaging}
\end{figure}

\subsection{离散动作 token：复用语言模型目标，但量化不是免费的}

将归一化动作 $a^{(j)}\in[-1,1]$ 均匀量化为 $B$ 个区间，
\begin{equation}
z^{(j)}=\operatorname{clip}\!\left(\left\lfloor\frac{a^{(j)}+1}{2}B\right\rfloor,0,B-1\right),
\qquad
\hat a^{(j)}=-1+\frac{2z^{(j)}+1}{B}.
\label{eq:uniform-action-token}
\end{equation}
于是训练可复用自回归交叉熵。RT-1 采用离散动作 token；RT-2 和 OpenVLA 让动作 token 与语言 token 共处输出词表\cite{brohan2022rt1,brohan2023rt2,kim2024openvla}。优势是训练基础设施成熟、输出概率可解释；代价是每个维度/时刻都占 token，量化误差与序列长度随控制维度和频率上升。

令分箱宽度 $\Delta=2/B$。若用箱中心解码且动作未被裁剪，单维最大量化误差为 $\Delta/2=1/B$；$B=256$ 时归一化误差上界约为 $0.0039$。这个数字必须乘回各维物理量程：同一归一化误差，对 $\pm1$ rad 的关节角与 $\pm0.1$ m 的位移不是同一控制误差。动作块含 $H$ 个时刻、每时刻有 $D$ 个独立离散维度时，朴素自回归序列长为 $DH$；例如 $D=7,H=50$ 就是 350 个动作 token。量化精度、解码延迟和上下文占用因此构成三方权衡。

\begin{figure}[H]
\centering
\includegraphics[width=0.94\textwidth]{v0.3-23-动作token量化.pdf}
\caption{均匀动作分箱的阶梯解码与误差。作者依据 OpenVLA 的动作 tokenizer 逻辑绘制\cite{kim2024openvla,openvla2024code}；图中使用 $B=8$ 仅为放大展示量化效应。}
\label{fig:action-token-quantization}
\end{figure}

FAST 指出逐维逐时刻分箱在高频灵巧动作上表现不佳，因而先对动作块做离散余弦变换，再对频域系数编码；FAST+ 在多机器人动作轨迹上学习通用 tokenizer\cite{pertsch2025fast}。这揭示了一个更一般的原则：动作 tokenizer 应利用时间相关性，而不是把每个采样点当作相互独立的“词”。

\subsection{扩散动作：表达多峰分布，以迭代换取计算}

令 $A^0=a_{t:t+H-1}$ 为干净动作块，前向过程加噪得到
\begin{equation}
A^\tau=\sqrt{\bar\alpha_\tau}A^0+\sqrt{1-\bar\alpha_\tau}\epsilon,
\qquad \epsilon\sim\mathcal N(0,I),
\end{equation}
噪声预测目标为
\begin{equation}
\mathcal L_{\mathrm{diff}}=
\mathbb E_{A^0,\tau,\epsilon}
\left\|\epsilon-\epsilon_\theta(A^\tau,\tau,o,s,l)\right\|_2^2.
\label{eq:diffusion-loss}
\end{equation}
推理从高斯噪声开始迭代去噪。Diffusion Policy 把这一过程与视觉条件、时间序列动作块和滚动时域执行结合，在其 12 项任务/4 组 benchmark 的实验集合中报告相对既有方法平均提升 46.9\%\cite{chi2023diffusion}。该数值不能直接与 OpenVLA 或 Gemini Robotics 的不同任务成功率相减。

训练和推理方向不能混为一谈。训练时先从数据集取干净动作 $A^0$，随机采样噪声等级 $\tau$ 和高斯噪声 $\epsilon$，只需构造一个带噪样本并回归噪声；推理时没有 $A^0$，必须从随机噪声出发，按噪声日程多次调用网络。不同随机初值可以收敛到不同动作模式，因此它能保留图\ref{fig:regression-mode-averaging}中被均值回归抹掉的选择。Diffusion Policy 还采用滚动时域：预测一段、执行前若干步、再用新观测重算；原论文 Fig.~3 把这件事与视觉条件去噪网络并列呈现\cite{chi2023diffusion}。

\begin{figure}[H]
\centering
\includegraphics[width=\textwidth]{v0.3-23-扩散动作生成.pdf}
\caption{扩散动作的训练加噪与推理去噪方向。作者根据 Diffusion Policy 的方法与原论文 Fig.~1、Fig.~3 重新组织绘制\cite{chi2023diffusion}，不复制原图像素；轨迹为教学示意。}
\label{fig:diffusion-action-generation}
\end{figure}

Octo 采用 Transformer 主干和 diffusion readout，重点不是单一动作空间上的最高分，而是让不同相机、本体状态、任务条件和动作头以模块化 token group 接入，并在目标域小数据上微调\cite{octo2024}。扩散的工程代价是一次动作块需要多次网络评估；若动作块执行期间不重规划，快速生成并不等于快速闭环。

\subsection{流匹配：学习速度场，仍需数值积分}

取噪声 $A^0\sim\mathcal N(0,I)$、数据动作 $A^1\sim p_{data}$，线性概率路径
\begin{equation}
A^\tau=(1-\tau)A^0+\tau A^1,\qquad
u_\tau=A^1-A^0.
\end{equation}
流匹配训练速度场 $v_\theta$：
\begin{equation}
\mathcal L_{\mathrm{FM}}=
\mathbb E\left\|v_\theta(A^\tau,\tau,o,s,l)-u_\tau\right\|_2^2.
\label{eq:ch23-flow-matching}
\end{equation}
推理时积分 $dA/d\tau=v_\theta(A,\tau,\cdot)$。$\pi_0$ 在预训练 VLM 上接入 flow-matching 动作专家，以连续动作块覆盖单臂、双臂和移动操作平台\cite{black2024pi0}。与扩散相比，流匹配可使用较直接的常微分方程路径，但延迟仍取决于积分步数、主干缓存和硬件，不能由“flow”一词推断为单步控制。

具体地，以 $N$ 步显式 Euler 积分为例，$\Delta\tau=1/N$，
\begin{equation}
A_{i+1}=A_i+\Delta\tau\,v_\theta(A_i,\tau_i,o,s,l),
\qquad i=0,\ldots,N-1.
\label{eq:flow-euler}
\end{equation}
每一步都要在当前动作状态重新估计速度。$\pi_0$ 的公开论文实现以 PaliGemma 视觉语言主干配合约 300M 参数的动作专家，预测 $H=50$ 的动作块，并在所述实验中以 10 个 Euler 步完成积分\cite{black2024pi0}。其块状因果注意力允许视觉语言前缀被缓存，动作专家则在不同积分步更新；因此“10 步”不等于把整个大主干完整运行 10 次，但也绝不是一次前向传播。

\begin{figure}[H]
\centering
\includegraphics[width=\textwidth]{v0.3-23-流匹配动作生成.pdf}
\caption{流匹配的训练速度场与 Euler 推理。作者依据 $\pi_0$ 的 flow-matching 动作专家改绘\cite{black2024pi0}；路径为二维教学投影。}
\label{fig:flow-matching-action-generation}
\end{figure}

\section{动作块与闭环：$H$ 大不等于长时程智能}

预测 $H$ 步动作块能减少视觉主干调用并让动作在时间上连贯。部署时通常只执行前 $h\le H$ 步，再用新观测重规划。$h=H$ 接近开环动作块，延迟低但对扰动迟钝；$h=1$ 闭环最紧，但计算和通信压力最大。合理选择依赖三个时间尺度：相机/状态更新频率、VLA 推理时间和低层控制周期。

例如 50 Hz 动作数据、$H=50$ 表示 1 秒动作块。若模型用 200 ms 生成一次并执行前 10 步，则重规划周期为 200 ms；若执行全部 50 步，模型虽能提前 800 ms 空闲，却可能在物体被移动后继续执行旧动作。论文若只写“输出 50-step chunk”而不报告 $h$、推理延迟、低层跟踪和观测新鲜度，读者无法判断闭环能力。

\section{主干与动作头如何分工}

PaLM-E 把图像、连续状态估计和文本编码交错输入语言模型，证明具身传感可以进入大型多模态主干，但其代表性机器人任务更多落在具身推理和规划，并不等价于一个高频低层控制策略\cite{driess2023palme}。RT-2 则明确把动作写入 VLM 输出空间，使网页语义预训练与机器人控制共同训练\cite{brohan2023rt2}。这一区分阻止了常见误写：不能因为模型“看见机器人并输出计划”就把它归为直接动作 VLA。

RT-1 的输入端以 EfficientNet 提取图像特征，并用 FiLM 接入语言条件；TokenLearner 压缩视觉 token 后，Transformer 按时间建模并输出离散动作。RT-2 改变的关键并非“再加一层 Transformer”，而是把机器人动作写进 VLM 的文本化输出空间，以网页视觉语言任务与机器人轨迹共同训练。PaLM-E 则把连续传感编码成与文本 token 同维的嵌入插入语言模型，更适合作为“感知信息如何进入具身推理主干”的代表。三者都使用 Transformer，但输入组织、监督目标和控制责任不同。

OpenVLA 组合 DINOv2 与 SigLIP 视觉特征、Llama 2 语言主干和离散动作输出，在 970k 真实机器人示范上训练；论文报告其 7B 模型在指定 29 项任务设置中相对 RT-2-X 的绝对成功率优势，以及相对从头训练 Diffusion Policy 的目标域微调优势\cite{kim2024openvla}。这些结果支持其设计在该评测协议下有效，不支持“7B 普遍优于 55B”这种脱离数据与动作接口的结论。

OpenVLA 中，DINOv2 偏向局部视觉表征，SigLIP 提供视觉--语言对齐特征，两路特征拼接后经投影器送入 Llama 2；动作再由语言模型词表末端的一组 token 自回归产生。于是视觉编码器回答“看到了什么”，投影器解决维度/分布接口，语言主干融合指令与图像上下文，ActionTokenizer 才负责连续控制量的编解码。任何一层冻结或微调，都会改变数据需求与部署代价，不能用“Llama 2 输出动作”一句话概括。

图\ref{fig:openvla-original-architecture}给出原论文的完整接口关系。读图时应沿编号 $1\rightarrow2\rightarrow3$ 追踪：DINOv2 与 SigLIP 分别编码同一图像，特征经 MLP projector 投到语言嵌入空间，再与指令 token 一起进入 Llama 2；输出端的 action de-tokenizer 才把词表 token 还原成七维控制量。原图把“视觉融合”和“动作反量化”画在主干两端，说明 OpenVLA 不是把连续向量直接塞给语言模型，而是依赖两个机器人专用接口夹住通用 VLM 主干。

\begin{figure}[H]
\centering
\includegraphics[width=0.96\textwidth]{v0.3-23-OpenVLA-Fig2原图.pdf}
\caption{OpenVLA 模型架构原图。引自 OpenVLA 原论文 Fig.~2（arXiv:2406.09246v3）\cite{kim2024openvla}，CC BY 4.0；从论文第 4 页裁切，仅缩放，未翻译或改绘。}
\label{fig:openvla-original-architecture}
\end{figure}

\subsection{Octo：模块化不是一句“可扩展”}

Octo 的主干输入由若干 \emph{token group} 组成：任务侧可以是语言或目标图像，观测侧可以包含多相机图像和本体状态。各自 tokenizer 先把异构输入变成统一宽度的 token，再由块状注意力掩码规定哪些组能相互读取。主干不直接把最后一个普通 token 当动作，而是加入专用 readout token；readout 汇聚上下文后连接 diffusion action head，产生连续动作块\cite{octo2024}。

这种设计把“共享知识”和“本体接口”分开。迁移到新机器人时，可以增加新的观测 tokenizer 以接收力/力矩或不同相机，也可以替换动作 readout 以适配新的关节维度；共享 Transformer 仍提供从 25 个数据集、约 80 万段轨迹中学到的跨任务表征。它并不意味着零成本跨本体：新增模块仍需目标域数据训练，动作归一化、时间对齐和相机标定也必须重做。

原论文 Fig.~2 还表达了改绘图容易弱化的一点：预训练序列可以反复交错 observation 与 readout，而微调阶段以虚线模块加入新观测和新动作空间。图\ref{fig:octo-original-architecture}保留这一原始视觉编码；图\ref{fig:octo-modular-architecture}则将其重排成中文模块图，便于逐步解释 tokenizer、主干、readout 和动作头的职责。两图承担不同任务，不作为重复插图计数。

\begin{figure}[H]
\centering
\includegraphics[width=\textwidth]{v0.3-23-Octo-Fig2原图.pdf}
\caption{Octo 模型架构原图。引自 Octo 原论文 Fig.~2（arXiv:2405.12213v2）\cite{octo2024}，CC BY 4.0；从论文第 3 页裁切，仅缩放，未翻译或改绘。}
\label{fig:octo-original-architecture}
\end{figure}
\endgroup

\begin{figure}[H]
\centering
\includegraphics[width=\textwidth]{v0.3-23-Octo模块化架构.pdf}
\caption{Octo 的模块化输入、Transformer 主干、readout 与动作头，以及目标域新增接口。作者依据 Octo 原论文 Fig.~2 改绘\cite{octo2024}，不复制原图像素。}
\label{fig:octo-modular-architecture}
\end{figure}

GR00T N1 把 VLM System 2 与 diffusion transformer System 1 紧耦合并端到端训练，数据混合真实轨迹、人类视频和合成样本\cite{nvidia2025groot}。Gemini Robotics 报告直接控制 VLA；同一技术报告另设 Gemini Robotics-ER，负责空间/时间理解、指向、轨迹与抓取预测等具身推理能力\cite{gemini2025robotics}。二者都提供重要系统证据，但完整训练代码、数据和权重并未公开到足以复现论文主张的程度，开源程度必须与论文新颖性分开评价。

\begingroup
\hbadness=10000
\begin{table}[H]
\centering
\caption{代表性路线的统一设计比较；结果数值因任务与协议不同，不作横向排行榜}
\label{tab:vla-design}
\scriptsize
\begin{tabularx}{\textwidth}{p{0.105\textwidth}p{0.14\textwidth}p{0.13\textwidth}p{0.13\textwidth}p{0.14\textwidth}Y}
\toprule
路线 & 主干/融合 & 动作表示 & 时序与闭环 & 开放证据 & 主要适用条件与限制\\
\midrule
RT-1 & EfficientNet/\allowbreak FiLM、TokenLearner、Transformer & 逐维离散 token & 短历史、实时逐步策略 & 论文与部分官方代码 & 单平台大规模多任务；语义主干不来自通用 VLM\\
PaLM-E & 传感编码与文本 token 交错进 LLM & 以语言/规划输出为主 & 高层推理接口 & 论文，模型闭源 & 不能直接当作低层 VLA 基线\\
RT-2 & VLM 与机器人数据共同训练 & 动作作为文本 token & VLA 直接控制 & 论文，主模型闭源 & 语义迁移强；量化、解码与闭环细节受实现约束\\
Diffusion Policy & 视觉编码器+条件去噪网络 & 连续 diffusion chunk & receding horizon & 论文和官方代码 & 多峰连续动作；多步去噪增加延迟\\
Octo & 模块化 Transformer & diffusion readout & chunk+目标域微调 & 论文、权重、代码 & 新观测/动作空间适配；规模小于大 VLM 路线\\
OpenVLA & DINOv2+\allowbreak SigLIP+\allowbreak Llama 2 & 256-bin 离散 token & 自回归动作输出 & 论文、权重、MIT 代码 & 易复用语言模型栈；高频动作 token 长度和量化受限\\
$\pi_0$ & 预训练 VLM+动作专家 & flow-matching chunk & 数值积分后滚动执行 & 论文、部分权重与 Apache-2.0 代码 & 灵巧连续动作；推理步数和跨本体适配仍需工程验证\\
GR00T N1 & VLM System 2 + DiT System 1 & diffusion action & 双系统端到端 & 报告与开放模型范围 & 人形双臂；开放资产与论文全训练栈并非完全等价\\
Gemini Robotics & Gemini 2.0 基础与机器人专用后训练 & 报告为直接连续控制 & 强调反应式闭环 & 技术报告，核心闭源 & 能力覆盖广；第三方复现与接口细节有限\\
\bottomrule
\end{tabularx}
\end{table}
\endgroup

\section{真实源码导读：OpenVLA 的均匀动作分箱}

OpenVLA 的 \texttt{ActionTokenizer} 清楚展示了式\eqref{eq:uniform-action-token}如何进入代码：先裁剪到动作范围，再逐维分箱，并把分箱编号映射到词表尾部的低频 token\cite{openvla2024code}。

\begin{lstlisting}[language=Python,caption={OpenVLA ActionTokenizer 的关键路径（MIT；为节省篇幅删去类型与批处理分支）}]
self.bins = np.linspace(min_action, max_action, self.n_bins)
self.bin_centers = (self.bins[:-1] + self.bins[1:]) / 2.0

action = np.clip(action, a_min=float(self.min_action),
                 a_max=float(self.max_action))
discretized_action = np.digitize(action, self.bins)
token_ids = self.tokenizer.vocab_size - discretized_action

discretized = self.tokenizer.vocab_size - action_token_ids
discretized = np.clip(discretized - 1, 0,
                      self.bin_centers.shape[0] - 1)
actions = self.bin_centers[discretized]
\end{lstlisting}
\sourceanchor{仓库 \path{openvla/openvla}；许可证 MIT；commit \texttt{c8f03f48\allowbreak af692657\allowbreak d3060c19\allowbreak 588038c7\allowbreak 220e9af9}；文件 \path{prismatic/vla/action_tokenizer.py}；符号 \texttt{ActionTokenizer}。代码逻辑已按 2026-07-13 的 \texttt{main} 页面核验；2026-07-14 通过只读远端查询固定当日 \texttt{main} HEAD，并使用 immutable permalink。}

代码暴露了三个模型图中不易看见的假设。第一，所有动作维统一裁剪到 $[-1,1]$，因此数据统计和反归一化是控制接口的一部分。第二，分箱边界有 256 个而中心只有 255 个，解码需要边界裁剪；“256 个 action token”不等于 256 个等宽可还原中心。第三，映射占用基础 tokenizer 词表末端 token，动作与语言共享词表是工程约定，不是物理上自然的表示。

\section{如何读实验表：先校准任务，再看数字}

跨论文比较至少检查五项：训练数据是否包含测试机器人；任务是 zero-shot、目标域微调还是从头训练；成功判据和试验次数；动作空间与控制频率；是否允许人工复位、重试或高层规划器。Octo 的官方评测覆盖 9 个真机 setup，并报告新力/力矩观测和新关节动作空间的微调\cite{octo2024}；OpenVLA 的 29 项任务、RT-1 的 700+ 数据采集任务与 Gemini Robotics 的能力演示不是同一统计总体。

因此本章不制作“综合成功率排名”。可以比较的是同一论文内部的受控消融，例如 OpenVLA 的视觉特征组合、Octo 的数据/架构选择、Diffusion Policy 的滚动执行与时间序列设计；跨论文只能比较证据边界和接口属性。对闭源模型，作者视频可以证明某次演示发生过，却不能估计总体成功率或失败分布。

\section{知识小结与综述判断}

\textbf{知识小结。}离散 token 把连续动作转成语言模型可处理的分类目标，但引入量化与序列长度；扩散和流匹配能表达多峰连续动作，却需要迭代去噪或数值积分；action chunk 在时间一致性、推理摊销和闭环反应之间权衡。主干是否具备视觉语言知识、动作头是否承担低层控制、执行后何时重规划，是比模型参数量更有解释力的坐标。

\authorjudgement{VLA 研究的核心竞争并不是“谁把最大的 VLM 接到机器人上”，而是谁能在语义泛化、动作分布表达、闭环频率、跨本体归一化和可复现评测之间形成可部署的组合。当前证据尚不足以宣告离散 token、diffusion 或 flow matching 中任何一种取得普遍胜利。对接触密集任务，VLA 最合理的输出通常是受约束动作块或技能参数，仍由本书动力学与接触章定义的实时控制、安全限幅和恢复状态机执行。}

\chapter[VLA 数据、训练目标、评测与泛化]{VLA 数据、训练目标、\\评测与泛化}
\label{chap:vla-data-evaluation}

\section{多机器人数据的难点是语义对齐，不是文件拼接}

Open X-Embodiment 把多机构机器人数据组织为标准格式，并用 RT-X 探索跨机器人正迁移\cite{openx2024}。标准 schema 解决“字段放在哪里”，却不自动解决动作语义：末端增量的参考坐标、夹爪正负方向、控制周期、相机命名和成功标签仍可能不同。每个 episode 至少需要机器人 ID、任务语言、观测/动作 schema 版本、坐标约定、标定、时间戳、终止原因和来源许可证。

设数据集 $D_i$ 有 $N_i$ 个 episode，训练采样权重为 $w_i$，则模型实际优化的是
\begin{equation}
\mathcal L=\sum_i w_i\,\mathbb E_{\tau\sim D_i}
\left[\sum_t m_{i,t}\ell_\theta(\tau_t)\right],\qquad \sum_iw_i=1,
\label{eq:vla-mixture-loss}
\end{equation}
其中 $m_{i,t}$ 是有效动作维/时间掩码。若令 $w_i=N_i/\sum_jN_j$，最大数据源会主导；若数据集等权，小型噪声集会被过采样。Octo 的多数据集训练\cite{octo2024}和 OpenVLA 的大规模真实示范训练\cite{kim2024openvla}都必须连同混合权重与规范化解释。

\begin{table}[H]
\centering
\caption{跨机器人数据混合的对齐检查}
\label{tab:vla-data-alignment}
\begin{tabularx}{\textwidth}{p{0.16\textwidth}YYYY}
\toprule
字段 & 必须记录 & 可统一内容 & 不应强行统一 & 错配后果 \\
\midrule
动作 & 坐标、单位、频率、维数 & 通用目标 schema & 关节拓扑与执行器限制 & 方向反转、尺度失控 \\
图像 & 相机 ID、时间、内外参 & 张量布局、颜色空间 & 视角与遮挡分布 & 错把相机差异当任务差异 \\
语言 & 原文、重写来源、任务 ID & tokenizer 接口 & 粒度与歧义 & 重复措辞支配任务采样 \\
终止 & 成功/失败/截断原因 & 枚举字段 & 场景特定判据 & 截断被误作成功或负例 \\
\bottomrule
\end{tabularx}
\end{table}

\section{归一化是模型接口的一部分}

对动作维 $j$，常用分位数归一化
\begin{equation}
\tilde a_j=2\frac{\operatorname{clip}(a_j,q_j^- ,q_j^+)-q_j^-}
{q_j^+-q_j^-}-1.
\end{equation}
$q_j^\pm$ 必须与机器人、动作 schema 和训练版本绑定。若部署时使用另一数据集统计量，张量形状仍正确，却会系统改变速度或位移。离散 token 又在归一化之后分箱，因此错误会被合法 token 掩盖。第 23 章的 OpenVLA 源码导读展示了这一边界。

\section{目标函数：共享主干不代表监督等价}

动作行为克隆可与语言生成、图文对比、价值或成功预测联合：
\begin{equation}
\mathcal L_{\mathrm{total}}=\lambda_a\mathcal L_{\mathrm{act}}
+\lambda_l\mathcal L_{\mathrm{lang}}+\lambda_v\mathcal L_{\mathrm{vision}}
+\lambda_s\mathcal L_{\mathrm{success}}.
\end{equation}
辅助目标能保留语义或改善表征，但也可能占用容量并与精细控制冲突。RT-2 通过视觉语言与机器人数据共训练研究知识迁移\cite{brohan2023rt2}；这类结论应由动作任务消融支持，而不能从 VLM 指标直接外推。

后训练还包括目标域微调、低秩适配、失败数据回放和偏好/价值校正。比较时必须标明：基础模型是否相同、目标域示范数、哪些层更新、是否使用评测环境数据、是否按最佳 checkpoint 选择。否则“少样本适配”可能隐含大量预训练覆盖或测试集调参。

\section{把“泛化”拆成四个正交问题}

任务泛化是新指令/技能组合；对象泛化是新实例或类别；场景泛化是背景、布局、光照和干扰变化；本体泛化是新机器人几何、传感器与动作接口。四者可同时变化，但证据强度不同。若新对象出现在相同桌面、相同相机且动作相同，只能支持对象轴结论；不能写成“开放世界泛化”。

\begin{table}[H]
\centering
\caption{VLA 评测证据的可比性边界}
\label{tab:vla-evidence-comparability}
\begin{tabularx}{\textwidth}{p{0.17\textwidth}YYYY}
\toprule
证据 & 可回答问题 & 必报条件 & 不能单独回答 & 常见误用 \\
\midrule
离线动作指标 & 数据拟合、量化/时延诊断 & 划分、掩码、归一化 & 闭环成功、安全 & 用 MSE 排名多峰策略 \\
仿真 & 大样本闭环比较 & 版本、种子、物理/视觉设置 & 真机可靠性 & 忽略 simulator gap \\
真机成功率 & 指定协议下任务结果 & 次数、初态、失败判据 & 广泛泛化 & 小样本百分比当稳定结论 \\
恢复/安全 & 扰动后的系统韧性 & 扰动、介入、风险定义 & 单模型能力 & 把安全层功劳归给 VLA \\
\bottomrule
\end{tabularx}
\end{table}

\section{成功率必须带不确定性}

若 $n$ 次试验成功 $k$ 次，点估计 $\hat p=k/n$。Wilson 区间为
\begin{equation}
\frac{\hat p+z^2/(2n)\pm z\sqrt{\hat p(1-\hat p)/n+z^2/(4n^2)}}
{1+z^2/n}.
\label{eq:wilson-success}
\end{equation}
例如 10 次成功 8 次并不等价于大样本下稳定的 80\%。论文若只给平均百分比、不报告试验次数、跨种子/任务方差和人工介入，就不能支持微小模型差异。

\begin{lstlisting}[caption={可复核真机评测协议的最小记录循环}]
for trial in preregistered_trials:
    reset_to_recorded_initial_state(trial.reset_id)
    start = capture_state_and_versions()
    result = run_with_fixed_timeout(policy, trial.instruction)
    outcome = independent_success_check(result, trial.success_rule)
    write_record(trial_id=trial.id, outcome=outcome,
                 intervention=result.human_interventions,
                 recovery=result.recovery_events,
                 safety=result.constraint_events,
                 model_hash=MODEL_HASH, start_state=start)
report_counts_and_confidence_intervals(group_by=["task", "generalization_axis"])
\end{lstlisting}

\section{怎样核查 SOTA 声明}

逐项检查任务是否相同、初态难度是否相同、是否允许重试/人工介入、控制频率和硬件是否相同、基线是否使用相同目标域数据、模型选择是否接触测试集。只要任一项不同，应写“在各自协议下报告”，而不是排序。第 23 章中的模型数字因此只支持论文指定设置的设计判断。

\section{知识小结与综述判断}

VLA 证据链从 schema 和动作语义开始，经采样权重、归一化、训练目标与后训练，最后落到分轴泛化和带置信区间的闭环评测。综述判断是：模型规模或平均成功率都不是独立证据；只有数据暴露、适配预算、试验协议与失败处理同时可比时，性能排序才有意义。

\chapter[开源复现、代码路径与部署接口]{开源复现、代码路径\\与部署接口}
\label{chap:vla-reproduction}

\section{复现对象：结果、训练还是部署}

“跑通仓库”至少有三种含义：用官方权重完成离线推理；在指定数据上重现微调趋势；接入真实机器人闭环。三者成本和证据强度递增。本章选择 OpenVLA 与 Octo 两条 MIT 开源路线，固定到 OpenVLA commit \texttt{c8f03f48...} 和 Octo release \texttt{v1.5}；不把网页当前 \texttt{main} 当可复现版本\cite{openvla2024code,octo15code}。

\begin{table}[H]
\centering
\caption{两条开源 VLA 路线的可复核执行链}
\label{tab:open-vla-repro-paths}
\begin{tabularx}{\textwidth}{p{0.14\textwidth}YYYY}
\toprule
路线 & 固定版本 & 模型/训练入口 & 动作接口 & 复现重点 \\
\midrule
OpenVLA & commit \texttt{c8f03f48...} & \path{vla-scripts/finetune.py}；Prismatic VLA & 离散 token 反量化为连续动作 & tokenizer、数据统计、VLM 显存 \\
Octo & release \texttt{v1.5} & \texttt{OctoModel}、模块化 transformer/readout & diffusion 动作采样 & task、mask、历史窗、动作统计 \\
\bottomrule
\end{tabularx}
\end{table}

第 23 章已逐行解释 OpenVLA \texttt{ActionTokenizer} 的均匀分箱，并记录完整 commit、文件与符号。这里不重复源码，而把它放回执行链：数据适配器计算动作统计，训练目标生成 action token，推理端反量化，机器人适配器再恢复真实单位与坐标。任何一步版本错配都会输出“合法但危险”的动作。

\section{环境与权重必须成为 manifest}

复现 manifest 至少包含代码 commit/tag、Python/CUDA/JAX/PyTorch 版本、锁文件哈希、权重 URL 与校验和、数据版本、GPU 型号、精度、随机种子和命令行。权重能加载不代表兼容：tokenizer 词表、视觉预处理或配置字段变化都可能不报错地改变结果。

\begin{lstlisting}[caption={建议的复现 manifest 骨架，不以浮动 latest 作为版本}]
experiment: octo_offline_smoke_test
repository: https://github.com/octo-models/octo
revision: v1.5
license: MIT
python: "3.10.x"
framework: "jax/flax versions from locked environment"
checkpoint: "octo-base-1.5"
checkpoint_sha256: "<record after verified download>"
dataset_schema: "oxe/rlds + local adapter v3"
robot_adapter: "widowx_delta_ee_v2"
action_statistics_hash: "<sha256>"
\end{lstlisting}

尖括号字段表示构建者必须在实际下载后填写，不能由报告编造。初稿的价值在于定义放行条件；最终冻结前应在目标机器生成真实 manifest 和日志。

\section{Octo：task、观测掩码与动作采样}

Octo 的官方入口把模型加载、任务创建和动作采样分开\cite{octo15code}。下列关键调用依据 \texttt{v1.5} README/API 压缩，保留决定接口的语句：

\begin{lstlisting}[language=Python,caption={Octo v1.5 推理关键路径（MIT；按官方 API 压缩）}]
model = OctoModel.load_pretrained("hf://rail-berkeley/octo-base-1.5")
task = model.create_tasks(texts=["pick up the red cup"])

observation = {
    "image_primary": image[None, None],
    "timestep_pad_mask": np.array([[True]]),
}
actions = model.sample_actions(
    observation, task,
    unnormalization_statistics=dataset_statistics["action"],
    rng=jax.random.PRNGKey(seed),
)
\end{lstlisting}
\sourceanchor{仓库 \path{octo-models/octo}；许可证 MIT；release \texttt{v1.5}；主要入口 \path{octo/model/octo_model.py} 的 \texttt{OctoModel.load\_pretrained}、\texttt{create\_tasks} 与 \texttt{sample\_actions}。本段依据官方 release README/API 压缩；未声明逐字复制某一连续行段。}

\texttt{timestep\_pad\_mask} 不是装饰字段：历史窗口开头缺帧时，它阻止填充值进入注意力。\texttt{unnormalization\_statistics} 决定动作恢复尺度。模型预训练动作块长度为 4 时，可以执行全部四步，也可以只执行首步滚动重规划；两种部署协议的时延和扰动恢复不同。

\section{OpenVLA：数据微调与动作 token 接口}

OpenVLA 把图像和指令送入 VLM 主干，再输出动作 token\cite{kim2024openvla}。复现路径应按顺序验证：官方权重单样本前向；tokenizer 编解码往返；目标数据 RLDS/schema 适配；小批过拟合；固定任务微调；离线回放；仿真或影子模式；低速真机。直接从“权重加载成功”跳到真机，会同时暴露图像通道、动作统计、坐标变换和控制周期错误。

一个重要单元测试是
\begin{equation}
\epsilon_{\mathrm{roundtrip}}=
\left\|\boldsymbol a-operatorname{decode}(\operatorname{encode}(\boldsymbol a))\right\|_\infty,
\end{equation}
并分别在分箱中心、边界、裁剪外值测试。该误差只能验证量化实现，不能证明机器人动作语义正确；还需用已知末端运动做坐标与单位测试。

\section{机器人适配层：把模型输出变成可拒绝的候选}

\begin{lstlisting}[caption={部署适配器的最小失败链}]
def deploy_vla(model, raw_observation, robot_contract):
    obs = robot_contract.preprocess_and_timestamp(raw_observation)
    assert obs.schema_hash == model.expected_schema_hash
    normalized_action = model.predict_action(obs)
    action = robot_contract.denormalize(normalized_action)
    command = robot_contract.to_control_frame(action)
    if not robot_contract.dimension_unit_and_rate_valid(command):
        return reject("adapter_contract_violation")
    if not shadow_mode_agrees_with_limits(command):
        return reject("shadow_validation_failed")
    return runtime_supervisor.filter(command)
\end{lstlisting}

适配层必须知道模型动作是绝对量还是增量、关节还是末端、世界还是工具坐标、位置还是速度，以及夹爪编码。第 21 章的监督器仍对结果做安全投影；适配器不能把模型声明为“已验证”而绕过运行时。

\section{复现放行表}

\begin{table}[H]
\centering
\caption{从离线推理到真机的阶段性退出条件}
\label{tab:vla-repro-gates}
\begin{tabularx}{\textwidth}{p{0.16\textwidth}YYY}
\toprule
阶段 & 必须留存 & 退出条件 & 不通过时禁止 \\
\midrule
环境/权重 & lock、hash、设备日志 & 同一输入可重复加载与前向 & 宣称复现结果 \\
数据适配 & schema、统计、样本可视化 & 时间/相机/动作语义人工抽检通过 & 大规模训练 \\
训练验证 & 配置、曲线、checkpoint、种子 & 小批过拟合与固定验证集通过 & 用失败归咎模型 \\
影子/仿真 & 时延、动作分布、约束事件 & 无单位/坐标/限幅异常 & 真机上电 \\
低速真机 & 风险评估、急停、监督日志 & 预注册任务与恢复测试通过 & 扩大速度和工作区 \\
\bottomrule
\end{tabularx}
\end{table}

\section{知识小结与综述判断}

开源复现不是复制安装命令，而是固定源、权重、数据、统计和硬件接口，并让每一阶段都有可审计退出条件。OpenVLA 暴露动作 token 与 VLM 微调链，Octo 暴露模块化 task/observation/action 接口；两者都不能省略机器人适配和独立安全监督。综述判断是：没有 immutable 版本、权重哈希、动作统计、schema 与部署日志的“成功运行”，只能称演示，不能称可复现证据。


\part{数据、仿真、评测、部署与安全}
\chapter{机器人数据工程与持续训练}
\label{chap:robot-data-engineering}

\section{数据单位是可追溯 episode，而不是视频文件}

机器人学习数据把图像、深度、关节、本体状态、动作、力觉、语言和任务结果绑在同一时间轴上。一个 episode 应被视为带版本的事件记录：它既保存连续 step，也保存机器人、场景、任务、标定、控制模式、操作者、终止原因和来源许可。RLDS 将 episode/step 语义显式化，支持顺序决策数据的记录、回放与转换\cite{ramos2021rlds}；Open X-Embodiment 进一步组织多机构机器人数据\cite{openx2024}。格式标准化解决互操作，不保证物理语义已经一致。

\begin{table}[H]
\centering
\caption{机器人 episode 的最小可追溯字段}
\label{tab:episode-schema}
\begin{tabularx}{\textwidth}{p{0.16\textwidth}YYYY}
\toprule
层级 & 必备字段 & 版本锚点 & 质量检查 & 典型静默错误 \\
\midrule
episode & ID、任务、场景、来源、终止 & schema、任务定义、许可证 & 唯一性、完整终止 & 截断被写成失败 \\
step & 测量时刻、发布时刻、动作、状态 & 时钟域、控制模式 & 单调、丢帧、延迟 & 最新值拼成伪同步 \\
传感器 & 原始值、单位、坐标、有效性 & 内外参、固件、驱动 & 范围、饱和、时间 & 旧标定配新图像 \\
动作 & 命令、实际执行、限幅后值 & 控制器、频率、归一化 & 单位、维数、跟踪误差 & 只保存策略建议 \\
结果 & 成功、失败类、介入、恢复 & 判据与标注指南 & 证据帧、复核者 & “未完成”等同失败 \\
\bottomrule
\end{tabularx}
\end{table}

DROID 在多地点、真实场景中采集大规模操作轨迹，展示了分布广度与采集组织的重要性\cite{khazatsky2024droid}。但“更多场景”同时带来操作者风格、相机放置、网络、照明和任务措辞的混杂。数据说明应记录采集目的、组成、处理、用途和风险；Datasheets for Datasets 提供了系统记录这些信息的框架\cite{gebru2021datasheets}。

\section{五类来源不能按同一置信度使用}

真机自主数据最接近部署分布，却受当前策略偏差限制；遥操作能高效提供目标行为，但包含操作者延迟和控制设备映射；穿戴/人类视频覆盖广，却缺少机器人动作与力；仿真提供精确标签和可控扰动，却存在物理与视觉差距；互联网数据提供语义知识，通常不能直接监督控制。进入训练混合前，应给每条来源独立的用途和禁止用途。

例如，人类视频可训练对象/动作语义，但若没有可靠 retargeting，不能把人体手腕轨迹直接当机器人末端命令。仿真接触标签可训练事件检测，却不能证明真实传感器标定正确。失败数据也不能统一标为负例：碰撞、识别错、操作者主动终止、网络超时和任务本身不可行需要不同语义。

\section{同步与切片：先保存原始时间，再做对齐}

设传感器 $i$ 的硬件时间为 $t_i$，映射到主时钟可写为
\begin{equation}
t=\alpha_i t_i+\beta_i+\epsilon_i(t),
\label{eq:clock-mapping}
\end{equation}
其中 $\alpha_i$ 是频率漂移，$\beta_i$ 是偏移，$\epsilon_i$ 是抖动。只在启动时减去一个 offset 隐含 $\alpha_i=1$，长 episode 中可能逐渐失配。切片时既要保存原始消息时间，也要保存对齐算法、插值方法和最大允许延迟。

动作监督应区分三条流：策略提出的 $a_t^{\mathrm{policy}}$、安全/限幅后的 $a_t^{\mathrm{safe}}$、执行器实际响应 $a_t^{\mathrm{exec}}$。训练策略通常以安全后命令为目标，系统辨识需要执行响应，事故分析又必须保留原始建议。覆盖其中任何一条都会破坏因果追踪。

\begin{figure}[H]
\centering
\setlength{\tabcolsep}{3pt}
\begin{tabular}{c@{\ $\rightarrow$\ }c@{\ $\rightarrow$\ }c@{\ $\rightarrow$\ }c}
\fbox{\begin{minipage}{0.16\textwidth}\centering 原始消息与\\不可变对象\end{minipage}} &
\fbox{\begin{minipage}{0.16\textwidth}\centering 同步、标定\\与 episode 切片\end{minipage}} &
\fbox{\begin{minipage}{0.17\textwidth}\centering 标注、质检\\与版本 manifest\end{minipage}} &
\fbox{\begin{minipage}{0.18\textwidth}\centering 训练集、回归集\\失败回流\end{minipage}}
\end{tabular}
\caption{从采集到持续训练的数据血缘。派生集可以重建，原始对象和转换 manifest 不得覆盖。}
\label{fig:data-lineage}
\end{figure}

\section{质量门：自动检查与人工抽检各管一半}

自动检查适合时间单调、字段缺失、范围、饱和、哈希、重复 episode 和动作跳变；人工抽检负责任务语义、对象身份、成功证据、隐私与异常操作。模型分数不应成为唯一过滤器，否则会删除模型不熟悉但真实有效的长尾样本。

\begin{lstlisting}[caption={episode 入库质量门的教学伪代码}]
def validate_episode(raw, manifest):
    require_hash_match(raw, manifest.raw_sha256)
    require_schema_version(manifest.schema_version)
    require_monotonic_clocks(raw.messages)
    aligned = synchronize(raw, manifest.clock_models)
    require_max_skew(aligned, limit_ms=manifest.max_skew_ms)
    require_calibration_versions(aligned, manifest.calibration_ids)
    require_action_contract(aligned.actions, manifest.robot_contract)
    flags = detect_dropouts_saturation_and_jumps(aligned)
    episode = slice_without_overwriting_raw(aligned, manifest.task_rules)
    episode.quality_flags = flags
    episode.human_review_required = flags.high_risk or sample_for_audit()
    return episode
\end{lstlisting}

质量分不应只输出一个标量。至少分为同步、视觉、状态、动作、任务标注和合规六类，使训练采样器能按用途选择：轻微图像过曝可能不影响本体控制训练，却不适合视觉策略。

\section{血缘、划分与持续训练}

每个派生数据集以内容哈希和转换 DAG 标识：原始对象哈希、输入 episode 列表、过滤规则、标注版本、归一化统计、代码 commit 和输出哈希。训练/验证/回归划分应按场景、对象、操作者或时间段分组，而不是随机拆 step；否则同一 episode 的相邻帧会泄漏。

在线失败回流要经过隔离队列。先由第 21 章事件日志把故障归为感知、规划、动作、控制、硬件或任务不可行，再决定进入监督微调、价值/偏好数据、回归测试或仅事故档案。新模型必须在固定回归集和新增故障集上同时通过，才能发布。

\section{案例：一个滑移 episode 的全链路}

机器人抓取杯子时触觉检测滑移，安全层降低速度并将杯子放回。采集端保存原始图像、触觉、策略动作、安全后动作和执行反馈；切片器用事件前后各两秒形成 episode；标注者将主失败标为“抓取接触/滑移”，将恢复标为“安全放置成功”，而非笼统“任务失败”。该 episode 进入三处：策略训练的困难正例、滑移检测回归集、恢复状态机测试。若后续策略成功但安全层介入更多，不应只看最终成功率宣布改进。

\section{隐私、版权与现场合规}

家庭、医疗和工厂数据可能包含人脸、屏幕、声音、工艺和客户资产。采集前应定义目的、告知/授权、最小化字段、保留期、访问角色和删除流程。公开可见的互联网视频不自动具有可训练或可再分发许可；数据 manifest 应分别记录“可内部训练”“可发布派生统计”“可公开样本”等权利。

\section{知识小结与综述判断}

机器人数据工程的核心不是存储容量，而是保持时间、坐标、动作、任务和权利语义的血缘。episode schema、不可变原始对象、版本化转换、分层质量门和故障分类共同支撑持续训练。综述判断是：无法从训练样本回溯到原始消息、标定、控制器和终止证据的数据规模，不应被当作可靠的模型能力基础。


\chapter{仿真、数字孪生与场景生成}
\label{chap:simulation-digital-twin}

\section{仿真的交付物是可执行实验，不是漂亮画面}

机器人仿真至少承担四种任务：生成训练经验、执行软件在环测试、复现现场故障、比较算法。四者对物理、视觉、速度和接口的要求不同。数字孪生也不等于高精度三维模型；它必须把特定设备、标定、控制器、场景参数和实测误差绑定到版本，能够回答“这个仿真对应哪台机器人、哪次维护之后的状态”。

仿真状态转移可写为
\begin{equation}
\boldsymbol x_{t+1}=F_{\phi}(\boldsymbol x_t,\boldsymbol u_t,\boldsymbol\xi_t),
\qquad \boldsymbol y_t=G_{\psi}(\boldsymbol x_t,\boldsymbol\eta_t),
\label{eq:simulator-model}
\end{equation}
$\phi$ 包含质量、惯量、摩擦、接触与执行器参数，$\psi$ 包含相机、材质、光照和传感器模型，$\boldsymbol\xi,\boldsymbol\eta$ 表示过程与观测扰动。平台名称不能替代这些参数和求解设置。

\section{平台比较应从用途出发}

\begingroup
\hbadness=10000
MuJoCo 面向多关节动力学与接触控制，采用广义坐标和专门接触求解方法\cite{todorov2012mujoco}；Gazebo 从开放的三维多机器人仿真与传感器/中间件集成出发\cite{koenig2004gazebo}；Habitat 2.0 把可交互家居资产、并行仿真和长时程家务 benchmark 连接起来\cite{szot2021habitat2}；Isaac Sim 强调 Omniverse/USD 资产、PhysX、渲染、合成数据与 ROS/软件在环链路。当前官方文档还显示物理传感器 API 在 6.0 发生迁移，说明“平台能力”必须附版本\cite{nvidia2026isaacsim}。

\begin{table}[H]
\centering
\caption{仿真平台按实验职责比较，而非制作单一排行榜}
\label{tab:simulator-comparison}
\begin{tabularx}{\textwidth}{p{0.13\textwidth}YYYY}
\toprule
平台 & 主要长项 & 资产/传感 & 并行与接口 & 选型时优先验证 \\
\midrule
MuJoCo & 控制、动力学、接触研究 & MJCF、相机/基本传感 & 轻量 API、批量依赖上层框架 & 接触参数、执行器、模型转换 \\
Gazebo & ROS 系统集成、多机器人 & SDF/URDF、插件化传感 & 中间件与仿真时间集成 & 插件版本、实时因子、时钟 \\
Habitat 2.0 & 室内导航与移动操作 & ReplicaCAD、家居对象 & 高吞吐任务环境 & 任务/对象状态与真实布局差异 \\
Isaac Sim & USD 场景、渲染、合成数据 & RTX/物理传感、Replicator & GPU、ROS 2、SIL/HIL & GPU/版本、PhysX、传感 API 迁移 \\
\bottomrule
\end{tabularx}
\end{table}
\endgroup

表中“长项”不是精度保证。即使同一引擎，步长、求解迭代、碰撞网格、接触刚度和 GPU/CPU 后端不同也会改变结果。跨论文速度只在相同场景、硬件、传感器和同步策略下可比。

\section{视觉资产、碰撞资产和动力学资产必须分层}

渲染网格追求外观细节，碰撞网格追求稳定且计算可控的接触几何，动力学资产还需质量、惯量、关节、限制、摩擦、柔顺和执行器。直接把高面数视觉网格当碰撞体，会产生尖锐接触、穿透和性能问题；用过度简化碰撞体又可能虚增抓取裕度。

\begin{figure}[H]
\centering
\setlength{\tabcolsep}{3pt}
\begin{tabular}{c@{\ $\rightarrow$\ }c@{\ $\rightarrow$\ }c@{\ $\rightarrow$\ }c}
\fbox{\begin{minipage}{0.16\textwidth}\centering CAD/扫描\\视觉资产\end{minipage}} &
\fbox{\begin{minipage}{0.16\textwidth}\centering 碰撞、关节\\质量与材料\end{minipage}} &
\fbox{\begin{minipage}{0.17\textwidth}\centering 任务脚本、随机化\\传感器与故障\end{minipage}} &
\fbox{\begin{minipage}{0.18\textwidth}\centering 训练、SIL\\回归与现场复现\end{minipage}}
\end{tabular}
\caption{仿真资产进入工程闭环的分层链路。每层都应有版本、校准证据和适用边界。}
\label{fig:simulation-asset-pipeline}
\end{figure}

数字孪生参数可由实测轨迹辨识：
\begin{equation}
\phi^*=\arg\min_{\phi}\sum_{k}
\left\|\boldsymbol y^{\mathrm{real}}_k-
\boldsymbol y^{\mathrm{sim}}_k(\phi)\right\|_{\boldsymbol W}^2
+\lambda R(\phi).
\label{eq:twin-calibration}
\end{equation}
必须报告激励轨迹、观测变量和参数可辨识性。多个摩擦/阻尼组合可能产生相近轨迹；得到一个低残差解不代表参数具有真实物理意义。

\section{场景生成要保持任务因果结构}

参数化场景不只是随机摆物体。生成器需要满足可达性、稳定放置、无初始穿透、任务谓词、视野与安全约束，并记录随机种子。生成失败的场景不能静默重采样，否则训练分布会偏向“容易生成”的情况。

\begin{lstlisting}[caption={带约束和 provenance 的场景生成教学伪代码}]
def generate_scene(template, seed, distributions):
    rng = Random(seed)
    scene = instantiate_versioned_assets(template.asset_manifest)
    params = sample_mass_friction_lighting_latency(rng, distributions)
    apply_parameters(scene, params)
    place_objects_with_task_constraints(scene, template.task_graph, rng)
    if has_penetration(scene) or not task_reachable(scene):
        return RejectedScene(seed, params, reason="invalid_geometry_or_task")
    settle(scene, duration=template.settle_time)
    if not stable(scene):
        return RejectedScene(seed, params, reason="unstable_after_settle")
    return Scenario(scene, seed, params, template.version,
                    engine_version=ENGINE_VERSION)
\end{lstlisting}

合成数据还需防止标签“过于完美”：真实分割边界、深度孔洞、曝光、运动模糊和时间延迟都不等于引擎 ground truth。当前 Isaac Sim 文档明确区分物理 ground truth 与后处理噪声\cite{nvidia2026isaacsim}；是否加入噪声、噪声是否经过实测校准必须记录。

\section{接触、摩擦与传感器是主要失真源}

\begingroup
\hbadness=10000
刚体接触包含离散事件、摩擦锥近似、求解容差和时间步依赖。柔性物体、线缆、织物、吸盘和颗粒更难。若策略利用仿真器特有的穿透或高频抖动，训练回报会提高而现实失败。域随机化通过变化视觉与动力学参数扩大训练分布\cite{tobin2017domainrandomization,peng2018dynamicsrandomization}，但随机范围不应替代校准：范围过窄不覆盖现实，过宽则让策略保守或学习错误不变量。
\endgroup

建议把差距拆为参数差距、结构差距和观测差距。参数差距可校准或随机化；结构差距如刚体模型无法表达软接触，不能靠扩大摩擦范围修复；观测差距需传感器噪声、延迟和缺失模型。

\section{仿真如何进入测试与故障复现}

训练使用随机场景，发布回归应使用冻结场景集；现场故障则从第 26 章 episode 提取机器人状态、对象位姿、控制器版本和事件时间线，重建最小可复现场景。若精确复现不可能，应明确保持了哪些因果条件，例如“延迟突增后抓取目标过期”，而非声称像素级数字孪生。

仿真测试结果按三类记录：仿真内通过；需要真机确认；已由配对真机试验验证。只有第三类能支持特定现实结论，且有效期受硬件与软件版本限制。

\section{知识小结与综述判断}

仿真的价值来自可控、可重复和可观测，而不是默认“逼真”。平台选型、资产分层、参数校准、场景约束、传感/接触边界和版本化回归共同决定证据强度。综述判断是：如果一个仿真资产没有 engine、参数、种子、任务脚本和与现实的误差记录，它最多是演示场景，不能称数字孪生或回归基线。

\chapter[Benchmark、真机评测与复现纪律]{Benchmark、真机评测\\与复现纪律}
\label{chap:benchmark-reproducibility}

\section{先冻结任务，再运行模型}

评测协议必须先于待测模型确定。一个任务实例至少由初始状态分布 $p(x_0)$、目标 $g$、允许动作与资源、超时 $T$、成功判据 $S$、失败判据 $F$、允许的恢复/人工介入和随机种子组成。若看到模型表现后再修改物体位置、超时或成功判据，数字不再是独立证据。

\begin{equation}
\mathcal E=(p(x_0),g,\mathcal A,T,S,F,
\mathcal I,\mathcal R,v_{\mathrm{env}},v_{\mathrm{policy}}),
\label{eq:evaluation-contract}
\end{equation}
其中 $\mathcal I$ 是人工介入规则，$\mathcal R$ 是复位/恢复规则，两个版本字段分别锚定环境与策略。CALVIN 通过连续控制和长时程语言任务形成可重复仿真 benchmark\cite{mees2022calvin}；Habitat 2.0 把场景、任务和模拟器共同发布\cite{szot2021habitat2}。它们提供受控问题，不代表真实家庭全部分布。

\section{三类 benchmark 的证据等级}

\begin{table}[H]
\centering
\caption{数据集、仿真与真机评测的证据职责}
\label{tab:benchmark-evidence-levels}
\begin{tabularx}{\textwidth}{p{0.17\textwidth}YYYY}
\toprule
类型 & 最适合回答 & 可控性/规模 & 主要外部效度风险 & 必须补充 \\
\midrule
离线数据集 & 表征、拟合、数据效率 & 高、可重复 & 行为分布固定、无闭环纠错 & 数据血缘、泄漏检查 \\
仿真闭环 & 策略交互、扰动、消融 & 高、可大规模 & 物理、视觉与任务差距 & 版本、参数、真机配对 \\
实验室真机 & 指定硬件上的系统能力 & 中、成本高 & 场景窄、操作者和复位偏差 & 预注册协议、完整计数 \\
现场部署 & 真实工作流与运维边界 & 低、风险高 & 选择性上线、暴露量变化 & 暴露时长、介入、事故率 \\
竞赛 & 规则下的系统集成与对抗 & 中、公开规则 & 赛制策略与开放部署不同 & 规则版本、赛程与复现说明 \\
\bottomrule
\end{tabularx}
\end{table}

证据不是简单由下到上替代。仿真适合百万次扰动，真机适合校验关键差距，现场适合评估可靠性和人机流程。完整结论通常是三角验证，而不是只选最有利的一层。

\section{指标要对应失效机制}

任务成功率为 $\hat p=k/n$，应连同试验数和区间报告。第 24 章给出 Wilson 区间；本章进一步区分首次成功率、允许恢复后的最终成功率与无人工介入成功率：
\begin{align}
p_{\mathrm{first}} &= N_{\mathrm{first\ success}}/N,\\
p_{\mathrm{recover}} &= N_{\mathrm{success\ after\ recovery}}/N,\\
p_{\mathrm{autonomous}} &= N_{\mathrm{success,\ no\ human}}/N.
\end{align}
只报告最终成功会掩盖重试成本，只报告首次成功又忽略恢复系统价值。

长时程任务还应报告完成子任务数和失败位置；安全应报告约束接近、过滤、急停、碰撞、超力和人工接管；时延应报告感知、推理、排队、网络、控制下发的 $p50/p95/p99$，而不是平均推理时间。能耗、吞吐和人员分钟数在产业案例中同样是结果变量。

\begin{figure}[H]
\centering
\setlength{\tabcolsep}{3pt}
\begin{tabular}{c@{\ $\rightarrow$\ }c@{\ $\rightarrow$\ }c@{\ $\rightarrow$\ }c}
\fbox{\begin{minipage}{0.16\textwidth}\centering 预注册任务\\与冻结版本\end{minipage}} &
\fbox{\begin{minipage}{0.16\textwidth}\centering 盲化初态\\独立成功判定\end{minipage}} &
\fbox{\begin{minipage}{0.17\textwidth}\centering 完整事件日志\\失败与介入分类\end{minipage}} &
\fbox{\begin{minipage}{0.18\textwidth}\centering 计数、区间\\分层与复现实验\end{minipage}}
\end{tabular}
\caption{可复核评测的执行链。失败 trial 必须与成功 trial 一样保留。}
\label{fig:evaluation-protocol}
\end{figure}

\section{统计纪律：种子、层级与多重比较}

机器人试验往往按任务、对象、场景、日期和机器人形成层级数据。把所有 trial 混成一个百分比会让高频简单任务主导。应预先定义宏平均（任务等权）或微平均（trial 等权），并同时给出分层结果。学习算法对随机种子、超参数和实现细节敏感；《Deep Reinforcement Learning That Matters》系统讨论了这些可重复性问题\cite{henderson2018deeprl}。

模型选择使用验证集，最终测试只运行预注册版本。若比较许多 checkpoint、提示词和阈值后只报告最好结果，实际发生多重比较。至少公开候选数、选择规则和全部主实验。负结果、失败视频和无效 trial 的排除理由必须可审计。

\begin{lstlisting}[caption={预注册真机评测与完整计数的教学伪代码}]
protocol = load_signed_protocol(PROTOCOL_HASH)
policy = load_model(POLICY_HASH)
for trial in protocol.randomized_trials:
    assert environment_version() == protocol.environment_version
    reset = independent_operator_apply_reset(trial.reset_id)
    start = capture_initial_state(reset)
    result = run_until_success_failure_or_timeout(policy, trial)
    verdict = independent_judge(result.evidence, trial.rules)
    append_only_record(trial.id, start, result.events, verdict,
                       interventions=result.interventions,
                       exclusions=result.exclusion_reasons)
assert recorded_trials() == protocol.scheduled_trials
report_counts_intervals_latency_quantiles_and_failure_strata()
\end{lstlisting}

“传感器日志损坏”等无效试验可以排除，但必须在解盲前定义规则，并同时报告计划、完成、排除和分析的数量。否则失败 trial 容易被事后解释为“无效”。

\section{复现不是得到相同小数点}

计算复现要求同代码、数据和配置得到相容结果；方法复现允许独立实现；结论复现要求在新团队、硬件或场景中仍支持同一方向的判断。机器人领域硬件差异使逐点相同不现实，因此应定义容忍区间和不变量，例如“在三种摩擦设置下均优于基线”，而非必须重现 73.2\%。

复现包至少包括环境/任务版本、容器或锁文件、模型和数据哈希、运行命令、随机种子、原始 trial 表、分析脚本、失败媒体和硬件配置。只提供均值表无法重新计算置信区间或审查排除项。

\section{RoboCup：规则化系统试验，而非兴趣专题}

RoboCup 各联赛通过公开规则、裁判、对手和现场条件测试感知、协作、运动、恢复与运维；2026 赛事资源仍按联赛分别维护规则\cite{robocup2026rules}。它的学术价值在于把组件放入受约束系统竞争，并暴露实验室难遇到的通信、复位和对抗问题。

但比赛名次混合硬件、软件、战术、维护、赛程和对手强度，不能直接作为某个 VLA 或控制算法的因果证据。用于综述时应引用具体联赛、年份、规则版本、任务和技术报告；若要证明算法贡献，仍需同队内部消融或可复现实验。

\section{知识小结与综述判断}

评测从冻结任务契约开始，以完整 trial 记录、独立判定、分层指标、置信区间和复现包结束。离线、仿真、真机、现场和竞赛分别回答不同问题。综述判断是：没有分母、初始条件、失败判据、人工介入和版本信息的成功率，不足以成为横向模型排名或部署结论的证据。


\chapter{硬件、ROS 2、实时系统与云边端部署}
\label{chap:hardware-ros2-deployment}

\section{部署不是把模型复制到机器人上}

机器人部署的输入不是一个模型文件，而是一组同时受物理和时间约束的工作负载：伺服控制要在硬截止期内读状态、计算并写命令；状态估计要保持时间一致；视觉和 VLA 推理可以较慢，却必须处理过期观测；日志、训练和舰队分析可以放到云端，但网络断开时机器人仍要进入确定的安全状态。因而选型应从任务频率、最坏时延、数据率、功耗和失效后果反推本体、执行器、传感器、计算与通信。

设第 $i$ 条链路包含采样、排队、传输、计算和执行延迟，则
\begin{equation}
L_i^{\mathrm{e2e}}=L_i^{\mathrm{sense}}+L_i^{\mathrm{queue}}+
L_i^{\mathrm{net}}+L_i^{\mathrm{compute}}+L_i^{\mathrm{act}},
\qquad
M_i=T_i-L_{i,p99}^{\mathrm{e2e}},
\label{eq:deployment-latency-budget}
\end{equation}
其中 $T_i$ 是周期或截止期，$M_i$ 是用 $p99$ 端到端时延估算的裕量。平均推理时间不能替代 $M_i$：一次排队尖峰就可能让高速关节多运动数厘米。硬实时回路还要以最坏执行时间和调度分析验证，$p99$ 只适合工程监测，不构成形式保证。

\begin{table}[H]
\begingroup
\hbadness=10000
\centering
\caption{按时间尺度拆分部署职责，而不是为整机指定一个“算力”数字}
\label{tab:deployment-timescales}
\begin{tabularx}{\textwidth}{p{0.16\textwidth}p{0.15\textwidth}YYY}
\toprule
职责 & 典型尺度 & 首要约束 & 合适位置 & 必留证据 \\
\midrule
驱动/伺服 & 亚毫秒至数毫秒 & 抖动、急停、总线 & MCU/实时 CPU & 周期直方图、超期计数 \\
状态估计 & 数毫秒至数十毫秒 & 时间戳、同步、丢帧 & 板载 CPU/GPU & 延迟分解、创新量 \\
视觉与动作策略 & 数十至数百毫秒 & 显存、吞吐、观测新鲜度 & 板载加速器 & p50/p95/p99、功耗 \\
任务规划/记忆 & 百毫秒至秒 & 语义正确性、可取消 & 板载或边缘 & 超时、回退、工具日志 \\
训练与舰队分析 & 分钟至天 & 数据治理、成本、版本 & 数据中心/云 & 数据和模型 manifest \\
\bottomrule
\end{tabularx}
\endgroup
\end{table}

\section{硬件链：力、带宽、热与能量要同时闭合}

本体负载、减速器峰值扭矩、连续扭矩、关节速度和制动方式决定动作包络；编码器、相机、深度、IMU 和力矩传感器决定可观测性；板载计算又受到电池、散热和机械空间约束。峰值 TOPS 不能说明模型能否持续运行，因为热降频、内存带宽、预处理、数据搬运和并发控制都会改变有效吞吐。

一个可审计的功耗预算应满足
\begin{equation}
P_{\mathrm{avg}}=P_{\mathrm{act}}+P_{\mathrm{compute}}+P_{\mathrm{sense}}+
P_{\mathrm{comm}}+P_{\mathrm{aux}},\qquad
t_{\mathrm{mission}}\leq \frac{\eta E_{\mathrm{battery}}}{P_{\mathrm{avg}}},
\label{eq:power-mission-budget}
\end{equation}
其中 $\eta$ 汇总电池可用深度、变换损耗和安全余量。算例：可用能量为 $1.2\,\mathrm{kWh}$、平均功耗 $400\,\mathrm{W}$、$\eta=0.8$ 时，上界仅约 $2.4\,\mathrm{h}$；若只报告待机功耗，就会系统性高估班次覆盖。峰值制动和再生能量还要单独校核，不能由平均式推出。

\section{ROS 2 把数据面、控制面和硬件接口显式化}

ROS 2 的 topic、service 和 action 分别适合流式数据、短事务与可反馈/可取消的长操作。QoS 不是“网络调优开关”，而是通信契约：可靠性、队列深度、截止期、寿命和活性策略要按消息语义设定，发布者与订阅者策略不兼容时甚至无法通信\cite{ros2jazzyqos}。相机预览可以允许 best effort，安全状态心跳则需要明确 deadline 与 liveliness 监测；二者套用同一 profile 通常都不理想。

\texttt{ros2\_control} 用 Controller Manager 连接控制器与硬件抽象，Resource Manager 管理 System、Sensor、Actuator 及其状态/命令接口；主循环按 read--update--write 执行\cite{ros2control2026}。这建立了清晰责任边界，但不会自动提供确定性：驱动阻塞、动态分配、日志 I/O、锁竞争和非实时回调都可能破坏周期。

\begin{figure}[H]
\centering
\setlength{\tabcolsep}{3pt}
\begin{tabular}{c@{\ $\leftrightarrow$\ }c@{\ $\leftrightarrow$\ }c@{\ $\leftrightarrow$\ }c}
\fbox{\begin{minipage}{0.17\textwidth}\centering 云端训练\\版本仓库\end{minipage}} &
\fbox{\begin{minipage}{0.17\textwidth}\centering 边缘规划\\VLA/视觉\end{minipage}} &
\fbox{\begin{minipage}{0.18\textwidth}\centering 实时控制\\状态估计\end{minipage}} &
\fbox{\begin{minipage}{0.17\textwidth}\centering 驱动、传感\\安全 I/O\end{minipage}}
\end{tabular}
\vspace{0.7em}

\begin{tabular}{c@{\quad}c@{\quad}c}
\fbox{签名模型/回滚} & \fbox{超时则取消/降级} & \fbox{越界则安全停机}
\end{tabular}
\caption{云边端不是按“智能程度”分层，而是按截止期、失效后果和断网自治分配职责。}
\label{fig:cloud-edge-device-profile}
\end{figure}

managed node 通过 configure、activate、deactivate、cleanup 和 shutdown 把资源准备与业务运行分开，适合相机、雷达和电机驱动的有序启停\cite{ros2jazzylifecycle}。生命周期状态机仍需系统编排器：某个节点进入 active 不代表标定有效、依赖已就绪或安全互锁闭合。

\begin{lstlisting}[caption={依据 ros2\_control 官方字段写成的教学配置；频率必须由实测和安全分析确定}]
controller_manager:
  ros__parameters:
    update_rate: 1000
    joint_state_broadcaster:
      type: joint_state_broadcaster/JointStateBroadcaster
    arm_controller:
      type: joint_trajectory_controller/JointTrajectoryController

arm_controller:
  ros__parameters:
    joints: [j1, j2, j3, j4, j5, j6]
    command_interfaces: [position]
    state_interfaces: [position, velocity]
    state_publish_rate: 100.0
    action_monitor_rate: 20.0
\end{lstlisting}

配置中的 $1000\,\mathrm{Hz}$ 只是请求周期。验收必须记录真实 read/update/write 分解、超期次数、CPU 隔离、调度策略和总线状态；若控制器在桌面负载下达标、录包或推理并发时超期，部署仍不合格。

\section{时间同步与数据新鲜度是跨层接口}

不同传感器硬件时钟可写成 $t=\alpha_i t_i+\beta_i+\epsilon_i$。第 26 章讨论数据集对齐；在线系统还要传播采样时刻、接收时刻和估计不确定度。只在消息到达时盖主机时间，会把网络抖动误当作物体运动。控制器应拒绝超过新鲜度阈值的状态，融合器应监测时钟漂移，日志应保存同步源和校准版本。

实时线程应避免无界阻塞、缺页、动态内存和不可控锁；非实时 ROS 回调通过有界缓冲区向实时循环交换最新值。实时不是把所有模块塞入高优先级线程，而是把少量必须守期的路径隔离，并让其失败模式可测。

\section{云边端更新必须带回退条件}

云端适合训练、离线评测和舰队聚合；板载端必须拥有断网后的感知、制动、状态监督和最低任务回退。模型发布包至少包含权重哈希、代码/算子版本、输入输出 schema、标定兼容范围、硬件目标、已通过测试集和回滚版本。灰度上线按机器人、场景和时间窗限制暴露，不能只做随机用户 A/B 测试，因为物理风险与场景分布不均匀。

更新流程应为：签名验证 $\rightarrow$ 离线兼容检查 $\rightarrow$ 影子运行 $\rightarrow$ 有界任务灰度 $\rightarrow$ 指标与安全事件门控 $\rightarrow$ 扩大部署；任何阶段出现 schema 不兼容、时延超限、安全过滤激增或未知失败，回滚到最后已知良好版本。SROS2 可用 DDS-Security 的证书、keystore 和 enclave 提供身份、权限与加密机制\cite{ros2security2026}，但密钥轮换、最小权限、供应链签名和失陷恢复仍需产品流程承担。

\section{知识小结与综述判断}

部署剖面把每条任务链的截止期、时延分解、功耗、硬件接口、通信 QoS、生命周期和回退版本写成可测试契约。ROS 2 与 \texttt{ros2\_control} 提供接口和运行时骨架，却不会自动产生实时性或安全性。综述判断是：模型大小、TOPS 和平均 FPS 都是局部指标；只有端到端 $p99$/最坏时延、热稳态、超期计数、断网行为和版本回滚共同闭合，才说明系统具备部署条件。

\chapter{安全、验证、标准与治理}
\label{chap:safety-verification-governance}

\section{安全对象是完整应用，不是单个模型}

具身系统的危险来自能量、运动、工具、环境、通信、数据和策略共同作用。模型输出正确并不保证制动距离足够；机械本体合规也不保证远程更新不会关闭互锁。风险分析因此要从预期用途和可合理预见误用出发，沿任务、人员、环境、硬件、软件和运维生命周期识别危险。ISO 12100 给出机械风险评估与风险降低的通用方法\cite{iso12100_2010}，但具体产品仍需适用的 C 类标准和合格人员完成工程判断。

风险不宜被压成一个脱离语境的分数。可将危险场景 $h$ 记录为
\begin{equation}
R_h=(S_h,E_h,P_h,A_h),
\label{eq:risk-tuple}
\end{equation}
其中 $S$ 为伤害严重度，$E$ 为暴露，$P$ 为危险事件概率，$A$ 为避免或限制伤害的可能性。不同标准的分级方法并不完全相同；四元组的用途是迫使团队写清假设、受影响人群和证据，而不是自创认证等级。

\begin{table}[H]
\begingroup
\hbadness=10000
\centering
\caption{四类风险域需要不同控制与证据}
\label{tab:safety-domains}
\begin{tabularx}{\textwidth}{p{0.16\textwidth}YYY}
\toprule
风险域 & 典型危险 & 首选控制 & 验证证据 \\
\midrule
物理/机械 & 夹挤、碰撞、跌落、工具伤害 & 本质安全、限能、隔离、防护 & 力/速/距离、制动、结构试验 \\
功能安全 & 传感失效、控制卡死、急停失效 & 独立安全功能、诊断、冗余 & PL/SIL 分析、故障注入、证明测试 \\
网络与数据 & 未授权命令、更新篡改、隐私泄露 & 身份、最小权限、签名、分区 & 威胁模型、渗透测试、审计日志 \\
模型行为 & 错误理解、分布外动作、绕过约束 & 能力包络、运行时过滤、人工接管 & 对抗任务、未知类、过滤/接管率 \\
\bottomrule
\end{tabularx}
\endgroup
\end{table}

控制优先级通常是消除危险、通过设计降低风险、增加工程防护，再辅以使用信息和培训。把警告弹窗或“请谨慎使用”放在最前面，是把设计责任转嫁给操作者。

\section{标准先看适用范围，再看编号}

\begin{table}[H]
\begingroup
\hbadness=10000
\centering
\caption{截至 2026-07-14 的标准与框架定位；状态变化需在冻结前复核}
\label{tab:safety-standards-scope}
\begin{tabularx}{\textwidth}{p{0.22\textwidth}YYp{0.20\textwidth}}
\toprule
文件 & 主对象 & 能回答什么 & 当前边界 \\
\midrule
ISO 12100:2010 & 通用机械 & 风险评估与降低方法 & 现行、待修订，不替代产品标准 \\
ISO 10218-1:2025 & 工业机器人本体 & 本体安全设计与信息 & 不覆盖完整应用/单元\cite{iso10218_1_2025} \\
ISO 13849-1:2023 & 安全相关控制部件 & SRP/CS 设计集成方法 & 不替应用指定所需 $PL_r$\cite{iso13849_1_2023} \\
ISO/TS 15066:2016 & 协作工业机器人系统 & 协作运行补充要求 & 现行但待修订，非通用服务机器人标准\cite{isots15066_2016} \\
ISO 13482:2014 & 个人护理机器人 & 人机接触相关危险与措施 & 现行；第二版仍为 FDIS\cite{iso13482_2014} \\
IEC 62443-4-1:2018 & IACS 产品开发 & 安全开发、缺陷、补丁和退役 & 网络安全生命周期，不替代机械安全\cite{iec62443_4_1_2018} \\
NIST AI RMF 1.0 & 跨行业 AI 风险 & 治理、映射、度量、管理 & 自愿框架且正在修订\cite{nist2023airmf} \\
\bottomrule
\end{tabularx}
\endgroup
\end{table}

同一机器人可能同时落入机械、工业机器人、移动平台、医疗器械、电气、无线电、数据保护和网络安全要求。能否宣称“符合”取决于产品用途、地区、集成方式和合格评定路径；引用一个标准编号不能代替适用性分析。尤其不能把“参考了某标准”“实验中通过阈值”和“经独立机构认证”写成同一等级。

\section{安全架构必须独立于学习策略}

学习策略可以提出动作，但不应拥有最终、不可撤销的能量权限。安全监督器读取经验证的状态和硬件安全 I/O，对速度、力、区域、姿态、通信新鲜度和设备状态执行确定性约束；遇到越界或监督器自身故障时进入预定义安全状态。

\begin{figure}[H]
\centering
\begin{tabular}{c@{\ $\rightarrow$\ }c@{\ $\rightarrow$\ }c}
\fbox{\begin{minipage}{0.20\textwidth}\centering 任务规划/VLA\\候选动作 $u_{\mathrm{nom}}$\end{minipage}} &
\fbox{\begin{minipage}{0.23\textwidth}\centering 安全过滤与仲裁\\约束、时效、模式\end{minipage}} &
\fbox{\begin{minipage}{0.20\textwidth}\centering 实时控制/驱动\\允许动作 $u^*$\end{minipage}}
\end{tabular}

\vspace{0.7em}
\begin{tabular}{c@{\quad}c@{\quad}c}
\fbox{独立急停/安全 PLC} & \fbox{安全状态估计} & \fbox{事件记录与人工接管}
\end{tabular}
\caption{学习策略位于受约束通道内；独立急停不能依赖同一模型、进程或普通网络。}
\label{fig:independent-safety-supervisor}
\end{figure}

对可微安全集合 $\mathcal C=\{x:h(x)\geq0\}$，控制屏障函数可把名义命令投影到满足约束的命令\cite{ames2017cbf}：
\begin{equation}
u^*=\arg\min_u \frac12\|u-u_{\mathrm{nom}}\|_W^2
\quad\text{s.t.}\quad
\nabla h(x)^\top(f(x)+g(x)u)\geq-\alpha(h(x)),\quad u\in\mathcal U.
\label{eq:safety-filter-qp}
\end{equation}
它要求状态、动力学、约束和求解时间足够可信；感知漏检、模型误差或 QP 超时会破坏保证。因此工程系统还需速度/力硬限制、watchdog、制动和安全 I/O。Simplex 架构用高性能控制器与经过更强验证的基线控制器切换，是另一种隔离复杂性的路线\cite{sha2001simplex}；切换条件和回切滞环本身也必须验证。

\section{验证从组件证据走向异常工况}

正常任务成功率只覆盖运行域的一小块。验证矩阵应交叉任务阶段、人员位置、负载、速度、网络、照明、摩擦、传感故障、执行器故障、模型版本和恢复动作。至少包括静态分析与单元测试、软件/硬件在环、仿真扰动、受控真机、故障注入、红队、现场灰度和回归复测。

安全事件率必须带暴露量：
\begin{equation}
\lambda_k=\frac{N_k}{H_{\mathrm{operation}}},
\qquad
\lambda_k^{\mathrm{task}}=\frac{N_k}{N_{\mathrm{tasks}}},
\label{eq:safety-exposure-rate}
\end{equation}
其中事件 $k$ 可以是保护停机、超力、碰撞、人工接管或网络拒绝。零事故而只有十分钟运行不能证明低风险；保护停机下降也可能是检测失效，必须同时审查近失事件和诊断覆盖率。

红队不只测试提示注入。还应尝试模糊指令、互相冲突的多用户命令、被遮挡的人、伪造二维码、重放传感器、过期地图、恶意工具输出、断网和降压。每个发现要转成威胁/危险场景、修复、回归用例和剩余风险，而不是只保留演示视频。

\begin{lstlisting}[caption={证据包驱动的安全放行教学伪代码}]
release = load_candidate(model_hash, software_bill, calibration_set)
assert applicable_standards_reviewed(release.intended_use)
hazards = load_signed_hazard_log(release.system_version)
for hazard in hazards:
    assert hazard.control_is_implemented
    assert verify(hazard.test_protocol, release).passes_acceptance
assert all_safety_functions_independent_and_tested()
assert fault_injection_and_red_team_regressions_pass()
assert rollback_image_is_signed_and_boot_tested()
deploy_to_bounded_exposure(release)
if safety_event_rate_exceeds_gate() or unknown_failure_occurs():
    rollback_and_quarantine(release)
\end{lstlisting}

\section{网络、数据与模型治理进入同一安全案例}

机器人节点能运动真实硬件，身份和权限必须收敛到最小能力：视觉节点不应写电机命令，日志节点不应加载控制器，远程运维凭据不应长期共享。软件物料清单、签名制品、漏洞响应、补丁窗口和退役密钥属于产品生命周期；IEC 62443-4-1 正是把网络安全活动延伸到缺陷管理、补丁和产品退役\cite{iec62443_4_1_2018}。

AI 治理补充模型特有的不确定性：训练数据许可与隐私、能力边界、分布漂移、人类监督、事件透明度和版本责任。NIST AI RMF 提供组织级 Govern、Map、Measure、Manage 框架\cite{nist2023airmf}，但不能替代具体机械安全功能，也不是认证证书。安全案例应把每个主张连接到需求、实现、试验、结果和剩余风险，并记录谁有权接受风险和撤回版本。

\section{知识小结与综述判断}

安全工程从预期用途和危险场景出发，通过本质安全设计、独立安全功能、网络/数据控制、模型运行时约束和分层验证形成证据包。标准必须按对象、版本和适用范围引用。综述判断是：论文中的安全过滤、厂商的“符合设计原则”和经适用合格评定的产品能力不是同一证据等级；只有危险日志、独立控制、异常工况测试、暴露量指标和可回滚放行共同闭合，才能把实验性防护推进为可运营的安全能力。


\part{应用与系统案例}
\chapter{制造与物流：结构化高价值任务}
\label{chap:manufacturing-logistics-cases}

\section{案例要落到任务单元，而不是行业名称}

制造和物流更容易形成机器人交付，不是因为任务天然简单，而是因为对象、工位、流量、人员区域和失败处置可以被工程化。一个任务单元至少包含来料与初态、感知与标定、动作与工具、质量判定、异常分流、安全边界和上下游节拍。只展示一次抓取或行走，无法说明系统是否能进入生产。

单位合格产出可写为
\begin{equation}
Q=\frac{N_{\mathrm{good}}}{H_{\mathrm{scheduled}}},\qquad
I_h=\frac{T_{\mathrm{human\ intervention}}}{N_{\mathrm{good}}},
\label{eq:case-throughput-intervention}
\end{equation}
其中 $Q$ 把失败、复位和停机计入排程时间，$I_h$ 把人工介入归一化到每件合格品。只报告机器人运动时间会遗漏上料、质检、异常清除和等待上游的成本。

\section{案例一：双臂椅装配揭示“结构化”仍需完整闭环}

Suárez-Ruiz 等使用商用双臂、平行夹爪、3D 相机和腕部力传感器完成 IKEA 椅装配\cite{suarezruiz2018ikea}。系统把零件定位、抓取、双臂协调、插入和力觉终止连接起来，说明结构已知、零件有限的任务也需要视觉与接触共同闭环。

\begin{figure}[H]
\centering
\begin{tabular}{c@{\ $\rightarrow$\ }c@{\ $\rightarrow$\ }c@{\ $\rightarrow$\ }c}
\fbox{零件检测/位姿} & \fbox{抓取与交接} & \fbox{力控插入/固定} & \fbox{质量判定/回退}
\end{tabular}
\caption{装配任务单元。任何一步失败都要有可观测判据和退出路径。}
\label{fig:assembly-workcell-case}
\end{figure}

该工作的重要边界同样清楚：装配顺序由工程师预先给定，零件和工位受控；论文证明研究原型在指定任务上的系统集成，不提供跨班次 OEE、换型时间、维护工时或认证证据。它适合作为“可变零件装配”的技术基线，不应写成通用装配已经解决。

从交付角度，泛化的价值不是让机器人理解任意家具说明书，而是减少新型号的夹具、视觉模板和轨迹重写。这个收益必须与新增验证成本比较：若模型能生成新动作，但每次换型仍需逐项安全确认和长时间调试，经济价值可能低于专用工装。

\section{案例二：仓内物流靠任务分解与规模运维}

Amazon 将仓内机器人分成库存搬运、单件抓取、包裹分拣和开放区域移动等任务。其 2023 年官方材料称 Sparrow 已在指定站点处理订单，Proteus 与 Cardinal 在出库码头组合测试\cite{amazon2023robots}；2025 年官方材料又称机器人部署量达到一百万，并称 DeepFleet 使机器人行程效率改善 10\%\cite{amazon2025millionrobots}。这些数字是运营方自报事实，不是独立审计结果。

\begin{figure}[H]
\centering
\small
\begin{tabular}{c@{\ $\rightarrow$\ }c@{\ $\rightarrow$\ }c@{\ $\rightarrow$\ }c}
\fbox{订单释放} & \fbox{搬运与排队} & \fbox{抓取与分拣} & \fbox{异常恢复}
\end{tabular}
\caption{仓内物流的系统闭环。机器人数量不是产能，流量和异常处理才决定吞吐。}
\label{fig:warehouse-flow-case}
\end{figure}

规模化的核心资产包括地图和交通规则、任务调度、充电、备件、软件发布、现场诊断和人员培训。单机策略只占其中一段。开放地面 AMR 要处理人与车混行，抓取系统要处理长尾包装、透明/反光物体、破损和拥挤；失败若统一送往人工异常站，系统仍可能有效，但 $I_h$ 必须进入成本模型。

\begin{table}[H]
\begingroup\hbadness=10000
\centering
\caption{两个案例的证据职责不可互换}
\label{tab:manufacturing-logistics-evidence}
\begin{tabularx}{\textwidth}{p{0.18\textwidth}YYYY}
\toprule
案例 & 已证明 & 未证明 & 人工角色 & 下一层证据 \\
\midrule
IKEA 椅装配 & 指定零件/序列下双臂视觉力控闭环 & 产线节拍、换型、长期可靠性 & 设计任务、准备工位 & 多型号、跨班次、完整故障统计 \\
Amazon 仓内机器人 & 企业公开声称大规模部署与指定流程使用 & 独立效率因果、介入和事故分母 & 上料、异常、维护、监督 & 站点分层数据、暴露量、审计 \\
\bottomrule
\end{tabularx}
\endgroup
\end{table}

\section{结构化场景为何更易交付}

高价值结构化任务可通过工装、照明、包装、缓冲区和异常站降低环境熵；过程质量有可测阈值；产量足以摊销集成与维护。基础模型最现实的贡献是覆盖传统规则的长尾、降低换型工程量、改善异常分类和人机接口，而不是删除全部结构。

验收应同时报告合格产出、首次成功、恢复后成功、人工分钟、停机、能耗、损伤、安全事件和版本。若新模型提高抓取率却增加推理时延，使上游阻塞或下游饥饿，局部指标改善并未转化为系统收益。

\section{知识小结与案例判断}

装配原型显示视觉、双臂协调、接触控制和任务序列必须闭合；仓内案例显示规模能力主要来自任务分解、流量控制和运维系统。案例判断是：结构化不是智能含量低，而是把不可控变化转化为可验证接口。模型泛化只有减少换型、异常和人工介入的总成本时，才形成交付价值。

\chapter{家庭、服务、医疗与康复}
\label{chap:service-healthcare-case}

\section{辅助进食是完整服务系统，不是抓取演示}

辅助进食同时涉及食物感知、叉取/舀取、无碰转移、入口定位、用户时机、卫生和社会互动。失败可能是没有获取食物、食物掉落、用户咳嗽、头部移动、餐盘移位或交互设备不可用。与仓内抓取不同，最终执行对象靠近人的面部，速度、力、工具形状和用户控制权都是核心需求。

一次入口成功是多个阶段串联：
\begin{equation}
P_{\mathrm{bite}}=P(A)P(T\mid A)P(D\mid A,T)P(U\mid A,T,D),
\label{eq:feeding-chain-success}
\end{equation}
其中 $A$ 是成功获取食物，$T$ 是安全转移，$D$ 是正确递送到口部，$U$ 是用户在合适时机接受。即使每阶段成功率均为 0.95，串联成功率也只有约 0.81；因此必须设计中途检测和恢复，而不能只优化叉取。

\section{从 PR2 实验到院外系统}

Park 等的主动辅助进食系统使用 PR2、视觉引导行为、用户界面和安全工具，并由 10 名无运动障碍参与者和 9 名运动障碍参与者评价\cite{park2020feeding}。该研究把“机器人把食物送到嘴边”展开为用户选择、视觉定位、主动靠近和可感知控制，是比单组件精度更强的端到端用户证据。

后续院外系统进一步改变硬件和自主组织。Nanavati 等与运动障碍社区研究者共同设计便携系统，在办公室、食堂等地点由 5 名用户和一名社区研究者完成用餐，并由一名社区研究者在家中连续 5 天使用 10 餐；论文报告总计超过 15 小时的真实环境使用\cite{nanavati2025feeding}。其关键经验不是“更强模型消除异常”，而是可定制界面、用户参与选择和可变自主帮助处理异常。

\begin{figure}[H]
\centering
\begin{tabular}{c@{\ $\rightarrow$\ }c@{\ $\rightarrow$\ }c@{\ $\rightarrow$\ }c}
\fbox{用户选食物/方式} & \fbox{获取与确认} & \fbox{转移/入口时机} & \fbox{接受、取消或恢复}
\end{tabular}
\caption{辅助进食状态链。用户始终保留选择、暂停和取消能力。}
\label{fig:feeding-state-chain}
\end{figure}

\begin{table}[H]
\begingroup\hbadness=10000
\centering
\caption{证据从组件和实验室走向院外，但尚未等于产品交付}
\label{tab:feeding-evidence-ladder}
\begin{tabularx}{\textwidth}{p{0.18\textwidth}YYYp{0.20\textwidth}}
\toprule
层级 & 主要问题 & 代表证据 & 新暴露的失效 & 仍缺证据 \\
\midrule
组件试验 & 食物检测/获取是否可行 & 单食品、重复 trial & 类别与形变长尾 & 整餐、用户和清洁 \\
实验室用户研究 & 人能否控制并接受系统 & PR2、两类参与者 & 头部、界面、舒适性 & 家庭空间与长期使用 \\
院外/家庭研究 & 异常和情境能否处理 & 多地点与 5 天部署 & 座姿、活动、社交、布置 & 大样本、月级可靠性 \\
可持续产品 & 能否安全持续服务 & 尚需认证与运营材料 & 维护、责任、支付、培训 & 事故率、成本与服务网络 \\
\bottomrule
\end{tabularx}
\endgroup
\end{table}

\section{可变自主不是“不够自动”的妥协}

完全自主要求系统识别所有食物、用户状态和社交时机，开放世界代价很高。共享自主把机器擅长的轨迹、避障和稳定递送与用户擅长的偏好、时机和异常解释组合。用户介入率
\begin{equation}
I_u=\frac{N_{\mathrm{user\ corrections}}}{N_{\mathrm{bite\ attempts}}}
\label{eq:user-intervention-rate}
\end{equation}
不应简单越低越好：主动选择和取消体现控制权；被迫频繁纠错才是负担。评测需把意愿性交互与故障补救分开，并报告任务负荷、独立性和尊严等用户结果。

系统应提供机械限力/可脱离工具、速度限制、口部附近专用轨迹、急停、用户暂停和护理者接管。视觉或模型置信度不能取代这些机制。清洁和食物安全、设备安装、坐姿差异、误吸风险和紧急处置需要专业流程；研究原型的伦理审批也不等同于医疗器械准入。

\section{隐私、责任与成本进入系统定义}

面部、语音、饮食和健康数据高度敏感。最小化采集、端侧处理、访问控制、保留期限和删除机制应在设计期确定。远程支持若能查看视频或控制机械臂，必须记录授权、身份、动作范围和审计日志。

责任按阶段分配：模型提出食物或动作候选，运行时保证约束，用户决定偏好和时机，护理者处理超出能力边界的情况，供应方负责维护和更新。成本不能只算机械臂，还包括安装、清洁耗材、培训、远程支持、故障上门和替代照护。若系统每天只节省少量护理时间，却需要高频技术支持，独立性价值可能存在，商业可持续性仍未证明。

\section{知识小结与案例判断}

辅助进食把感知、操作、用户界面、安全和异常恢复绑定在一条靠近人体的服务链。院外研究比实验室成功率更能暴露空间和社会情境，但样本与时长仍有限。案例判断是：服务机器人最重要的泛化机制往往不是完全自动，而是让用户以低负担控制系统、识别异常并安全恢复；产品结论还需要长期可靠性、准入、责任和单位服务成本证据。

\chapter{农业、建筑、巡检与危险环境}
\label{chap:field-dangerous-environment-case}

\section{现场任务由环境与通信共同定义}

农业、建筑、矿山、能源巡检和灾害响应的共同点不是“户外”，而是地形、光照、粉尘、水、温度、通信和维护条件难以固定。实验室成功的导航器若依赖连续网络、洁净镜头或人工复位，在地下现场可能立刻失效。本章以 DARPA Subterranean Challenge 的 CERBERUS 系统为完整案例，不再逐项罗列行业。

DARPA 将任务定义为在未知地下环境中快速建图、导航、搜索并定位目标\cite{darpa2021subt}。竞赛同时引入隧道、城市地下和洞穴环境，把通信退化、垂直结构、黑暗、粉尘和无 GPS 变成系统约束。它仍是规则化竞赛，不代表矿山或消防的全部法规、耐久和人员流程。

\section{CERBERUS：异构系统而非单一最强机器人}

CERBERUS 组合 ANYmal 腿式机器人、常规与耐碰撞飞行器、多模态视觉/LiDAR/惯性定位、探索规划、目标检测、集中地图优化和通信中继，在 2021 决赛取得胜利\cite{tranzatto2022cerberus}。Tunnel/Urban Circuits 的系统论文记录了现场部署与限制\cite{tranzatto2022tunnel}。

\begin{figure}[H]
\centering
\begin{tabular}{c@{\ $\leftrightarrow$\ }c@{\ $\leftrightarrow$\ }c}
\fbox{腿式/飞行探索机器人} & \fbox{中继、投放通信节点} & \fbox{基站、操作员与地图融合}
\end{tabular}

\vspace{0.6em}
\begin{tabular}{c@{\quad}c@{\quad}c}
\fbox{局部自治与避障} & \fbox{失联继续/返航} & \fbox{任务分配与人工监督}
\end{tabular}
\caption{地下异构系统的通信与控制结构。自主性由机器人和基站共同组织。}
\label{fig:cerberus-system-case}
\end{figure}

腿式平台越障能力强但续航有限；飞行器进入狭窄/高处但通信和能量更紧；地面中继与投放通信节点扩展网络；集中优化改善多机器人地图一致性。异构性的代价是备件、软件、标定、操作和故障模式增加，因此统一的任务/地图接口比统一本体更重要。

任务效用可抽象为
\begin{equation}
J=w_c C+w_d N_{\mathrm{detected}}-w_e E-w_h H-w_r R,
\label{eq:field-mission-utility}
\end{equation}
其中 $C$ 是覆盖或信息增益，$N_{\mathrm{detected}}$ 是正确定位目标数，$E$ 是能耗，$H$ 是人工负担，$R$ 是机器人丢失或安全风险。竞赛分数主要显式奖励目标和时间，真实运营还需提高风险、耐久和人员安全权重。

\section{完全自主、遥操作与共享自主}

\begin{table}[H]
\begingroup\hbadness=10000
\centering
\caption{现场自治组织方式的适用条件}
\label{tab:field-autonomy-modes}
\begin{tabularx}{\textwidth}{p{0.18\textwidth}YYYp{0.20\textwidth}}
\toprule
模式 & 优势 & 脆弱点 & 合适任务 & 必报指标 \\
\midrule
完全自主 & 可在失联区持续工作 & 未知场景判断与恢复困难 & 高重复、边界清楚、低后果 & 无人时长、失败自恢复 \\
连续遥操作 & 人直接理解复杂语境 & 带宽、时延、操作者负荷 & 短程高价值精细操作 & 链路可用率、人员分钟 \\
共享自主 & 机器执行局部闭环，人管目标/异常 & 接管界面与责任切换复杂 & 地下探索、巡检、应急 & 介入原因、接管延迟、恢复率 \\
多机器人监督 & 一人覆盖更大区域 & 告警洪泛与地图冲突 & 搜索、测绘、巡检 & 人机比、覆盖、冲突和丢机 \\
\bottomrule
\end{tabularx}
\endgroup
\end{table}

CERBERUS 的“自主”并非没有人：操作员在基站分配任务、监控地图和机器人健康，并在需要时干预；机器人在局部执行探索、规划和避障。现场可扩展性的关键指标是每名操作员能可靠监督多少机器人、每小时需要多少干预，而不是是否存在人工。

\section{现场约束如何反向塑造系统}

\begin{table}[H]
\begingroup\hbadness=10000
\centering
\caption{环境约束到技术机制的映射}
\label{tab:field-constraint-mechanism}
\begin{tabularx}{\textwidth}{p{0.20\textwidth}YYp{0.23\textwidth}}
\toprule
现场约束 & 设计机制 & 失败检测 & 运维证据 \\
\midrule
无 GPS、退化视觉 & LiDAR/视觉/IMU 融合与地图优化 & 协方差、回环冲突、漂移 & 轨迹真值片段、重定位率 \\
间歇通信 & 本地自治、中继、存储转发 & 心跳、队列、链路质量 & 覆盖图、失联时长、数据丢失 \\
崎岖/狭窄地形 & 腿式与飞行异构平台 & 足端滑移、碰撞、卡滞 & 通过率、能耗、损伤 \\
有限续航 & 信息增益规划、返航阈值 & 电池、温度、剩余路径 & 有效任务时长、换电工时 \\
高维护成本 & 模块化载荷、远程诊断、备件 & 健康监测与故障码 & MTBF、MTTR、现场工具 \\
\bottomrule
\end{tabularx}
\endgroup
\end{table}

农业还会增加季节与生物变化，建筑会增加动态工序与临时结构，工业巡检会增加防爆和停产窗口。它们共享现场方法，但不能从 SubT 成绩直接外推。每个新行业都要重新定义危险、通信、耐久、清洁、责任和单位任务收益。

\section{从竞赛证据到长期运营还差什么}

SubT 提供公开规则、未知现场、时间压力和跨团队比较，证据强于剪辑 demo；但赛事持续时间短，团队可集中专家，机器人损失和维护成本与商业现场不同。运营需要月级/年级暴露量、天气和季节覆盖、备件、远程支持、培训、事故与近失事件、数据回传和合规记录。

商业模式也受自治组织反向牵引。若通信和异常要求专家长期在线，系统更像远程服务而非硬件销售；若一名操作员能监督多机且维护标准化，RaaS 才可能摊薄成本；若任务低频但危险，价值可能来自减少人员暴露而非提高吞吐。

\section{知识小结与案例判断}

CERBERUS 的现场能力来自异构本体、多模态定位、探索规划、通信工程、地图融合和人工监督共同作用。竞赛证明了受规则约束的高强度系统集成，不证明长期商业运营。案例判断是：危险环境最现实的路线通常是有界局部自治加共享监督；评价必须同时报告覆盖、目标、失联、能耗、人员负担、丢机和维护，而不能只给“自主完成”标签。


\part{产业、企业、供应链与政策}
\chapter{企业能力与交付证据矩阵}
\label{chap:company-evidence-matrix}

\section{比较对象不是公司名，而是可交付系统}

企业比较若从融资额、视频热度或参数表开始，结论会被披露策略支配。本章把“能力”限定为：企业能否在给定任务、现场和责任边界内，持续交付合格产出。统一记录任务单元、模型与数据、本体、客户部署、制造、运维、商业安排、风险和证据日期；缺失字段保留为空，不用宣传语补齐。

部署证据分五级：E0 为剪辑或现场 demo；E1 为指定工位试验；E2 为客户确认、采购或付费安排；E3 为跨班次持续运行并给出暴露量；E4 为跨客户/站点重复交付且存在可审计的可靠性和经济性。等级衡量公开证据，不等于企业真实能力。

\begin{figure}[H]
\centering
\scriptsize
\begin{tabular}{c@{\ $\rightarrow$\ }c@{\ $\rightarrow$\ }c@{\ $\rightarrow$\ }c@{\ $\rightarrow$\ }c}
\fbox{E0 展示} & \fbox{E1 工位试验} & \fbox{E2 客户/合同} & \fbox{E3 持续运行} & \fbox{E4 重复/审计}
\end{tabular}
\caption{交付证据梯度。升级需要新增分母，而不是更换宣传措辞。}
\label{fig:delivery-evidence-ladder}
\end{figure}

\section{五个样本在同一字段下的公开证据}

Agility 与 GXO 在 2024 年公布多年期 RaaS 安排，Digit 在仓内把 AMR 上的周转箱转移到输送线；企业随后公开称累计搬运超过 10 万箱\cite{agility2024gxo,agility2026totes}。前者提供客户、任务和商业安排，后者提供累计暴露量，但仍未公开排程小时、介入分钟和独立审计。

Figure 对 BMW 工位披露 1,250 小时运行、9 万余零件装载和 3 万辆车相关产出，并说明前臂是主要硬件失效点\cite{figure2025bmw}；2026 年又披露 Figure 03 的产量、在制品检查点和首过良率\cite{figure2026production}。这组材料比 demo 多出运行和制造分母，但数值仍来自供应商。

Apptronik 与 Mercedes-Benz 的 2024 材料明确使用“pilot”，与 Jabil 的 2025 合作同时包含试用和代工扩产意图\cite{apptronik2024mercedes,apptronik2025jabil}。因此可确认客户合作和制造伙伴，不能据此声称已经形成连续生产。UBTECH 的港交所公告披露 Walker S1/Walker C 小批量采购合同和预付款，2025 年年报进一步报告产能与工业场景进展\cite{ubtech2025contract,ubtech2025annual}；合同事实强于展示，产能和订单量仍需与实际交付、验收和续购分开。

Unitree 官方材料称 G1 推出量产版、H1 已量产，并展示汽车工厂搬运等场景\cite{unitree2024wrc}。公开材料能支持“平台可购买/量产与场景展示”，却缺少同口径的客户运行分母，故不把平台扩散直接写成自主任务交付。

\begin{table}[H]
\begingroup\hbadness=10000
\centering\small
\caption{代表企业公开证据矩阵（截至 2026-07-14）}
\label{tab:company-evidence-matrix}
\begin{tabularx}{\textwidth}{p{0.12\textwidth}p{0.17\textwidth}p{0.19\textwidth}p{0.16\textwidth}Y}
\toprule
样本 & 任务/本体 & 最高公开证据 & 制造/运维线索 & 不可外推边界 \\
\midrule
Agility & Digit；仓内料箱转运 & E3：客户 RaaS、累计箱数 & Arc 平台、现场集成 & 无统一 uptime、介入和经济性 \\
Figure & F.02/F.03；汽车上料 & E3：运行小时、零件与产出自报 & BotQ、失效点和良率自报 & 单客户/任务，缺独立审计 \\
Apptronik & Apollo；制造试点 & E2：客户 pilot 与制造合作 & Jabil 生产合作 & 未披露连续运行分母 \\
UBTECH & Walker；汽车制造 & E2：采购、预付款与年报披露 & 产能与量产自报 & 订单、交付、验收不可互换 \\
Unitree & H1/G1；平台与搬运展示 & E1--E2：量产产品和场景材料 & 平台可得性 & 缺客户持续运行和运维数据 \\
\bottomrule
\end{tabularx}
\endgroup
\end{table}

\section{模型、本体、制造和运维不能互相替代}

模型视频回答“某版本在某初态能否产生动作”；本体规格回答负载、速度和接口上限；制造信息回答能否稳定重复制造；客户数据才回答合格产出与人工成本。四者是串联系统，任何一项都不能代替其他项。特别是“年产能”是设备与节拍能力，不是实际产量；“订单”不是收入确认；“进入工厂”可能只是演示、实训或隔离试验。

最低可比披露应包括软件/硬件版本、任务与对象分布、排程小时、完成与合格数、人工介入、停机原因、安全暴露量、备件与更新时间。没有这些分母时，作者只判断证据成熟度，不排名技术优劣。

\section{知识小结与综述判断}

企业证据应沿展示、试验、合同、持续运行和重复交付逐级积累。五个样本分别显示任务聚焦、模型--本体整合、制造合作、采购与平台可得性，但公开口径仍不完整。综述判断是：当前最稀缺的不是机器人视频，而是跨版本、跨班次、含人工和失败分母的客户级数据；在其出现前，企业矩阵应保持“可核事实 + 明确缺口”，不制造总排名。

\chapter{供应链、规模化与单位经济性}
\label{chap:supply-chain-unit-economics}

\section{BOM 只是总成本的起点}

机器人交付同时消耗本体、末端、工装、集成、验证、运维、人工恢复和停机机会成本。供应商未披露的 BOM 不应被猜测；更可复算的对象是特定任务在合同期内的总拥有成本（TCO）：
\begin{equation}
C_{\mathrm{TCO}}=C_{\mathrm{cap}}+C_{\mathrm{integration}}+C_{\mathrm{maintenance}}+C_{\mathrm{energy}}+C_{\mathrm{human}}+C_{\mathrm{downtime}}-V_{\mathrm{residual}}.
\label{eq:robot-tco}
\end{equation}
若合同期合格产出为 $N_{\mathrm{good}}$，任务单位成本为 $c_u=C_{\mathrm{TCO}}/N_{\mathrm{good}}$。其中资本项需按寿命和资金成本摊销；人工项必须包含监督、异常清除、远程接管和维护，而不只计算上下料员。

制造 KPI 应有统一定义。ISO 22400-1 给出制造运营 KPI 的通用框架\cite{iso22400_1_2014}。本书采用
\begin{equation}
N_{\mathrm{good}}=H_s\,A\,r\,Y,
\qquad A=\frac{H_s-H_{\mathrm{down}}}{H_s},
\label{eq:effective-output}
\end{equation}
其中 $H_s$ 为排程小时，$A$ 为可用率，$r$ 为运行时理论产率，$Y$ 为良率。提高峰值动作速度只改变 $r$；若跌倒、热保护、换电和等待恢复拉低 $A$，总产出仍可能下降。

\section{关键部件的耦合约束}

执行器把电机、传动、轴承、编码器、驱动器、制动和热路径耦合在关节内。峰值扭矩不能代表连续能力；减速器间隙、效率、冲击寿命与反驱性共同决定控制品质。灵巧手把执行器密度、线缆寿命、触觉、污染防护和维修性压缩到更小体积，通常也是长尾故障源。

感知与算力同样受功耗和散热约束。以 Jetson AGX Orin 为例，官方规格给出 15--60 W 可配置功耗区间\cite{nvidia2026orin}；这只是计算模块而非整机功耗。提高视觉分辨率、模型频率或冗余感知会增加计算、内存带宽、热设计和电池负担，并可能减少有效工作时间。

\begin{table}[H]
\begingroup\hbadness=10000
\centering
\caption{供应链对象必须连到系统指标}
\label{tab:supply-chain-couplings}
\begin{tabularx}{\textwidth}{p{0.16\textwidth}YYY}
\toprule
对象 & 采购规格不够的原因 & 系统级验证 & 常见规模化失效 \\
\midrule
执行器/减速器 & 峰值不等于热稳态与寿命 & 工况谱、冲击、回差、热循环 & 批次离散、密封、润滑、线缆 \\
传感器 & 精度不等于污染/遮挡鲁棒性 & 标定漂移、同步、失效注入 & 供应替代导致重新标定 \\
算力/电池 & TOPS、kWh 不等于任务续航 & 端到端时延、峰值电流、热降频 & 模块 EOL、认证和热一致性 \\
灵巧手/末端 & 自由度不等于可维护操作能力 & 抓取谱、碰撞、清洁与换件 & 小件多、维修慢、寿命分散 \\
\bottomrule
\end{tabularx}
\endgroup
\end{table}

Figure 的 BotQ 材料说明其从 CNC 原型件转向压铸、注塑、冲压等模具工艺，并建立 MES、在制检测和供应商准入\cite{figure2025botq,figure2026production}。这说明“量产设计”改变零件数、连接方式、测试点和工艺，而不是把原型图纸简单复制；但产能、产量、首过良率和客户良率仍是四个不同指标。

\section{一个可复算的任务级算例}

设某三年任务单元的资本与集成摊销为 90 万元，运维和能源 30 万元，排程 12,000 小时，理论产率 50 件/小时，良率 $Y=0.98$。表\ref{tab:unit-economics-sensitivity}仅为教学假设，不代表任何企业报价。人工恢复按 80 元/小时、停机机会损失按 600 元/小时计算。

\begin{table}[H]
\centering
\caption{可用率对单位成本的敏感性（教学算例）}
\label{tab:unit-economics-sensitivity}
\begin{tabular}{lrrr}
\toprule
情景 & 可用率 $A$ & 合格产出/件 & 基础摊销/元每件 \\
\midrule
调试态 & 0.60 & 352,800 & 3.40 \\
稳定态 & 0.85 & 499,800 & 2.40 \\
高可用 & 0.95 & 558,600 & 2.15 \\
\bottomrule
\end{tabular}
\end{table}

基础摊销仅用 120 万元除以产出，尚未计人工和停机。若调试态比稳定态多停机 3,000 小时，停机损失增加 180 万元，远大于假设中的能源项。因此采购评审需要对 $A$、$Y$、人工介入率和故障修复时间做敏感性分析；用最低硬件售价做决定会把主要成本留在合同外。

\section{销售、集成、RaaS 与软件平台}

硬件销售把资产和利用率风险更多交给客户，适合任务稳定且客户有运维能力的场景；系统集成适合一次性工位改造，但收入与工程人力耦合；RaaS 把前期资本转换为服务费，供应商则承担 uptime、备件和残值风险。Agility--GXO 的多年期安排证明 RaaS 已进入人形机器人合同实践\cite{agility2024gxo}，却不自动证明其服务毛利为正。

软件/数据平台只有在多台、多站点复用时才可能脱离项目制。商业合同至少应定义任务和验收、可用率窗口、计划停机、人工责任、数据权利、软件更新、网络安全、事故责任、退出和残值；否则“按结果付费”无法计算结果。

\section{知识小结与综述判断}

供应链的技术问题是部件性能、工艺能力、批次一致性、可测试性和可维修性的联合优化；商业问题是把这些变量折算为合格产出成本。综述判断是：规模化的关键不在某个 BOM 数字，而在设计降复杂度、制造良率、现场可用率和低人工介入共同收敛。融资或规划产能只能说明资源，不能替代单位经济性。


\chapter{资本、政策、标准与区域格局}
\label{chap:policy-regional-landscape}

\section{先区分四类文本}

资本提供研发、库存和工厂爬坡的时间，却不证明技术成熟。政策目标表达资源方向；法律法规产生适用范围内的义务；标准提供设计、测试或管理基线，是否强制取决于引用和合同。以下比较截止 2026 年 7 月 14 日，后续适用日期不写成已经生效的义务。

\begin{figure}[H]
\centering
\small
\begin{tabular}{c@{\ $\rightarrow$\ }c@{\ $\rightarrow$\ }c@{\ $\rightarrow$\ }c}
\fbox{政策/资金方向} & \fbox{研发与场景计划} & \fbox{产品验证/标准} & \fbox{法律义务与市场监督}
\end{tabular}
\caption{政策到准入不是单向自动转化；每一步需要责任主体和证据。}
\label{fig:policy-to-market-access}
\end{figure}

\section{中国：供应链、公共平台与场景牵引}

工信部 2023 年《人形机器人创新发展指导意见》覆盖“大脑、小脑、肢体”、整机、关键部件、软件、示范场景、标准和治理，并提出 2025/2027 阶段目标\cite{miit2023humanoid}。它是产业指导文件，不是对单一产品的认证。

2026 年工信部与国务院国资委的专项行动把工业、服务和特种领域的真实任务单元列为实景实训载体，强调常态化部署和经济可行性\cite{miit2026fieldtraining}。这一变化把资源从“做出样机”进一步拉向场景、数据和迭代，但企业仍需分别满足机械、电气、网络、数据和行业监管要求。

\section{美国：创新基础设施、行业规则与自愿框架}

美国 2025 AI Action Plan 将政策分为加速创新、建设 AI 基础设施和国际输出三支柱\cite{whitehouse2025aiaction}，并非机器人产品的统一准入法。NIST AI RMF 1.0 是自愿、非行业特定的风险管理框架\cite{nist2023airmf}；其 Govern--Map--Measure--Manage 可用于组织证据，但不能替代机械安全认证、职业安全、产品责任或医疗器械等行业义务。

因此，美国路径更像“通用 AI/算力政策 + 行业监管 + 标准与责任诉讼”的组合。企业需要按使用场景确定主管机构和合同责任，不能把遵循自愿框架写成获得政府批准。

\section{欧盟：AI 与机械产品规则叠加}

欧盟 AI Act 已于 2024 年生效，但采用分阶段适用：部分条款已适用，主要规则计划自 2026 年 8 月 2 日起适用，部分嵌入受监管产品的义务更晚\cite{eu2024aiact}。截至本章日期，“即将适用”与“已经适用”必须区分。机器人是否属于高风险 AI 需根据用途、角色和附录范围判断，不能因使用神经网络自动归类。

Machinery Regulation (EU) 2023/1230 主要自 2027 年 1 月 20 日起适用\cite{eu2023machinery}。具身系统进入欧盟市场时，AI 风险、机械危害、网络安全、技术文档、合格评定和上市后监测可能叠加；模型更新若改变安全功能或预期用途，还会触发配置与责任问题。

\section{日本与韩国：实施环境、联盟和标准化}

日本 METI 2026 AI Robotics Strategy 把机器人友好环境、用户导入支持、研发/社会实施枢纽和分行业路线图放在一起\cite{meti2026robotics}。它反映日本从强机器人制造基础向 AI、数据和实际导入补链，评价成效仍需看部署、劳动力替代质量和维护网络。

韩国 MOTIE 2024 融合机器人国际标准战略计划推进国内标准和国际提案，覆盖康复与家庭服务等方向\cite{motie2024standards}；2025 K-Humanoid Alliance 则把整机、零部件、研究机构和基础模型开发组织起来\cite{motie2025humanoid}。联盟和谅解备忘录是协作基础，不等于采购、认证或量产。

\begin{table}[H]
\begingroup\hbadness=10000
\centering\small
\caption{区域路径的同口径比较（截至 2026-07-14）}
\label{tab:regional-policy-matrix}
\begin{tabularx}{\textwidth}{p{0.1\textwidth}YYYY}
\toprule
区域 & 主要抓手 & 对研发/交付的直接影响 & 准入结构 & 关键不确定性 \\
\midrule
中国 & 指导意见、公共平台、实景实训 & 场景、供应链、数据与中试资源 & 通用安全叠加行业监管 & 目标到可审计交付的转化 \\
美国 & 创新/算力政策、行业监管、自愿框架 & 资本与基础设施强，按行业合规 & 多机构、标准、合同与责任体系 & 机器人专门规则与州差异 \\
欧盟 & AI Act + Machinery Regulation & 文档、风险、更新与上市后责任前置 & 横向 AI 与产品法规叠加 & 分期适用和分类解释 \\
日本 & AI 机器人战略、友好环境、行业路线 & 强调用户导入与系统实施 & 既有产品/行业规则与标准 & 数据/软件补链和扩散速度 \\
韩国 & 产业联盟、基础模型与国际标准 & 联合研发、部件和标准提案 & 产品安全与标准体系 & 联盟成果到客户部署 \\
\bottomrule
\end{tabularx}
\endgroup
\end{table}

\section{资本与政策如何进入技术判断}

资本和政策只在改变约束时进入路线判断：是否提供长期数据采集、昂贵验证、模具与库存爬坡、公共测试场、采购需求或合规基础设施。融资额、估值、补贴目标和联盟成员数不进入模型能力评分。可验证的中间结果是测试设施建成、标准发布、采购完成、产品通过合格评定、部署给出暴露量和客户续约。

\section{知识小结与综述判断}

五个区域都在连接 AI 与机器人，但政策工具不同：中国偏供应链和场景牵引，美国偏创新基础设施与行业治理，欧盟以前置规则叠加产品法规，日本偏实施环境，韩国偏联盟和标准。综述判断是：区域竞争最终落在数据、部件、制造、验证、部署和责任体系的闭合速度；政策文本是约束和资源的证据，不是技术胜负结论。



\part{综合判断与开放问题}
\chapter{技术路线争议与开放研究问题}
\label{chap:technical-debates}

\section{争议必须改写成可比较实验}

“端到端还是模块化”“数据还是模型”若不固定任务、预算、失败代价和接口，就只是口号。一个有效争议至少给出竞争命题、共同评测集、控制变量、失败分层和否证阈值。本章不预测赢家，而是列出什么证据会迫使我们修改判断。

\section{端到端与模块化：学习边界放在哪里}

RT-2 与 OpenVLA 证明视觉语言表征可以直接进入动作预测\cite{brohan2023rt2,kim2024openvla}。SayCan 和行为树则把高层选择与可执行技能、显式控制流分开\cite{ahn2022saycan,colledanchise2018behaviortrees}。端到端路线减少手写接口，有机会利用联合数据；模块化路线便于局部验证、替换和故障隔离。二者并非二元：实际系统常把感知和策略学习化，同时保留低层控制、安全监控与恢复状态机。

现有论文很少在相同硬件、数据、算力和安全约束下比较两种边界。可证伪条件是：若端到端系统在未见对象/干扰上持续降低人工介入，同时保持最坏情况停止距离和故障定位能力，强模块分解的必要性下降；若小分布漂移反复产生不可诊断失效，且模块化替换能以较低数据成本恢复，则“统一模型自然更稳健”被否证。Simplex 与控制屏障函数说明学习策略外仍可放置可验证安全层\cite{sha2001simplex,ames2017cbf}。

\section{通用与垂类：复用收益能否覆盖验证成本}

Octo、OpenVLA 和 Open X-Embodiment 探索跨任务、跨数据源训练\cite{octo2024,kim2024openvla,openx2024}。ACT 与 Diffusion Policy 则在受控任务上展示动作建模\cite{zhao2023act,chi2023diffusionpolicy}。通用模型的价值不是参数更多，而是新任务所需数据、工程和停机验证更少；垂类系统的优势是分布、工具和验收边界清楚。

缺失实验应同时报告零样本成功、少样本适配曲线、旧任务回退、推理资源和安全再验证工时。若通用预训练在多个本体上把达到同一可靠性所需现场数据和工程时长稳定降低，且不增加关键失败，通用路线成立；若每个工位仍需大量重采、重新标定和安全审查，所谓通用性只是初始化。

\section{世界模型：显式预测何时值得}

PlaNet 与 DreamerV3 展示了在潜空间预测并学习价值的路线\cite{hafner2019planet,hafner2023dreamerv3}。显式世界模型适合反事实、长时回报和数据复用，但机器人接触、多主体和部分可观测环境会迅速放大模型偏差。反应式策略可避开长滚动误差，却难以解释长期资源和前置条件。

公平实验需固定真实交互预算，让世界模型与无模型/短视策略面对相同遮挡、接触和环境变化，并按预测时域报告校准误差。若更长想象时域没有改善闭环任务，或性能增益可由更大行为数据解释，则“预测是关键机制”证据不足；若在稀缺真机数据下显著减少危险试错并改善长时恢复，显式模型价值成立。

\section{动作表示：没有脱离控制频率的最佳 token}

动作块 ACT、连续扩散、流匹配与离散 token 分别改变多峰建模、平滑性、采样时延和错误传播\cite{zhao2023act,chi2023diffusionpolicy,lipman2023flowmatching,pertsch2025fast}。第 22--24 章已说明动作空间还受坐标系、频率、低层控制器和数据采样影响。

需要在相同观测、骨干、数据和低层接口下比较动作表示，报告端到端时延、接触峰值、重规划频率和恢复。若某表示只在离线误差上占优却增加接触冲击或闭环滞后，不能称为更优；跨任务稳定收益且对频率/本体扰动不敏感，才支持其通用优势。

\section{数据与模型规模：有效覆盖不是样本总数}

Open X-Embodiment、DROID、$\pi_0$ 和 GR00T 把跨机构数据、规模化采集与基础模型连接起来\cite{openx2024,khazatsky2024droid,black2024pi0,nvidia2025groot}。但 episode 数不能表示接触状态、失败、恢复、对象和本体覆盖；混合数据还会引入动作语义与标定不一致。

规模实验应绘制数据量、模型量、覆盖度和真实成功率的联合曲线，并保留来源/任务/本体外测试。若增加近重复数据或参数在分布外不再改善，瓶颈已转向覆盖、接口或控制；若控制覆盖后仍有稳定幂律收益，继续扩展规模才有依据。

\begin{table}[H]
\begingroup\hbadness=10000
\centering\small
\caption{五组争议的缺失实验与否证信号}
\label{tab:debates-falsification}
\begin{tabularx}{\textwidth}{p{0.15\textwidth}YYp{0.27\textwidth}}
\toprule
争议 & 共同实验 & 关键指标 & 否证信号 \\
\midrule
端到端/模块化 & 同硬件数据与安全约束 & 介入、诊断、停止边界 & 成功率高但失效不可隔离 \\
通用/垂类 & 多任务适配与回退 & 数据/工时、旧任务保持 & 每任务仍需完整重做 \\
世界模型 & 同真机交互预算 & 校准、长时回报、危险试错 & 长预测不改善闭环 \\
动作表示 & 同骨干与控制接口 & 时延、冲击、恢复 & 离线优而闭环差 \\
数据/模型规模 & 覆盖受控的缩放实验 & 分布外成功、边际收益 & 近重复扩展无收益 \\
\bottomrule
\end{tabularx}
\endgroup
\end{table}

\section{五个开放问题的可测门槛}

灵巧操作需要跨物体、材料和接触模式报告损伤与恢复；长时程自主需要在小时/天尺度报告无人工任务时长和故障链；跨本体迁移需要把适配数据、映射误差和低层重整定分开；端侧模型需要在功耗、热稳态和最坏时延下比较能力；持续学习需要证明新数据上线后旧能力、安全案例和可回滚性不退化。

\section{知识小结与综述判断}

路线优劣由任务分布、闭环频率、数据预算、失效代价和验证责任共同决定。综述判断是：短期最可信的系统仍会混合学习与显式工程边界；长期是否收敛到更统一模型，取决于其能否在真实暴露量下同时降低人工介入、适配成本和不可诊断风险。

\chapter{结论与可证伪研究议程}
\label{chap:conclusions-research-agenda}

\section{八条结论及其证据位置}

本章不重新复述模型和企业名录，只回答全书核心问题。表\ref{tab:conclusions-evidence-map}中的章节号是结论的主证据位置，不表示其他章节无关。

\begin{table}[H]
\begingroup\hbadness=10000
\centering\small
\caption{结论--证据章节映射}
\label{tab:conclusions-evidence-map}
\begin{tabularx}{\textwidth}{p{0.05\textwidth}Yp{0.2\textwidth}}
\toprule
序 & 可检验结论 & 主证据章节 \\
\midrule
1 & 具身智能是受物理、时序和安全约束的闭环系统，不是语言模型加机械外壳。 & 第 3、6、18--21 章 \\
2 & 接触、状态估计、规划和低层控制仍决定可执行边界；基础模型改变接口与泛化方式，不取消这些约束。 & 第 5--8、14--17 章 \\
3 & 机器人学习的主要瓶颈从单次成功转向数据覆盖、分布偏移、恢复和持续验证。 & 第 9--13、26--28 章 \\
4 & VLA 的差异必须沿动作表示、数据、目标函数、闭环频率和部署接口比较，模型名称本身没有解释力。 & 第 22--25 章 \\
5 & 可交付系统需要感知、状态、规划、技能、控制、安全和运维接口闭合；端到端策略也需承担明确责任边界。 & 第 18--21、29--30 章 \\
6 & 结构化高价值任务会先规模化，因为环境、质量和异常可工程化；开放环境的价值需与人工支持同时核算。 & 第 31--33 章 \\
7 & 企业成熟度由客户级持续运行、制造一致性和运维分母决定，融资、demo、订单和规划产能不可互换。 & 第 34--35 章 \\
8 & 政策和标准改变资源与准入约束，但技术胜负最终仍由可审计性能、单位经济性和责任闭环决定。 & 第 30、36--37 章 \\
\bottomrule
\end{tabularx}
\endgroup
\end{table}

跨本体数据和通用策略已展示迁移可能\cite{openx2024,octo2024}。真实部署材料也开始给出运行与产出分母\cite{figure2025bmw,agility2026totes}；但这些证据尚未共同证明任意环境、长期无人值守和正单位经济性。

\section{保守与激进情景只由触发条件定义}

\textbf{保守情景}被以下事实触发：跨站点成功率不随数据/模型扩大而提高；人工介入与硬件故障使 TCO 长期高于专机或人工；跨本体适配仍需逐任务重采和重验证；法规、保险和责任证据无法随软件更新闭合。此时市场主要停留在结构化制造、物流、危险作业和高价值辅助任务，通用本体更多作为研发平台。

\textbf{激进情景}要求同时出现：多客户系统以版本化口径公开跨季度运行、低介入和安全暴露量；通用预训练显著降低新任务数据与工程工时；量产良率、备件和远程运维使单位成本跨过客户阈值；模型更新可通过自动回归、风险案例和合格评定管理。只满足模型能力或硬件降价之一，不足以触发该情景。

\begin{figure}[H]
\centering
\small
\begin{tabular}{c@{\quad}c@{\quad}c@{\quad}c}
\fbox{跨站点低介入} & \fbox{适配成本下降} & \fbox{制造/运维收敛} & \fbox{更新可验证}
\end{tabular}
\caption{激进情景的联合触发器。四项是“与”关系，不是任选一项。}
\label{fig:aggressive-scenario-triggers}
\end{figure}

\section{可证伪研究议程}

下一阶段研究应优先建设八类公共证据：
\begin{enumerate}[leftmargin=2em]
  \item 以排程小时、动作次数和人机接触小时为分母的统一暴露量；
  \item 含失败、恢复、人工接管和损伤标签的数据集，而非只保留成功轨迹；
  \item 在接触、遮挡和部件退化下的状态估计与安全控制联评；
  \item 控制数据量、模型量和任务覆盖度的缩放实验；
  \item 报告适配数据、重整定和验证工时的跨本体迁移；
  \item 小时到周尺度、允许环境变化的长时程自主基准；
  \item 带功耗、热稳态、最坏时延和网络中断的端侧评测；
  \item 把合格产出、人工分钟、停机、能耗、损伤和 TCO 同时纳入技术评测。
\end{enumerate}

每项都可被否证。例如“持续学习降低维护成本”必须在上线新版本后证明旧任务、安全回归和人工工时不恶化；“跨本体基础模型降低适配成本”必须在多个本体上以同一成功阈值比较数据和工程量。

\section{最终判断}

具身智能已经形成可复现的学习方法、跨机构数据、开源模型和早期客户运行证据，但尚未形成覆盖开放环境、长期自主和统一经济性的通用答案。最有价值的研究不是继续扩大叙事，而是把物理机制、数据血缘、闭环评测、制造运维和责任证据接成同一条可审计链。任何未来判断，只要不能说明何种观测会使其失效，就不应进入本书结论。


